\documentclass[11pt]{ctexart}
\usepackage[a4paper,margin=1in]{geometry}
\usepackage{graphicx}
\usepackage{xcolor}
\usepackage{hyperref}
\definecolor{refgray}{gray}{0.45}
\title{\textbf{市场最新观点汇总}\\[0.4em]{\large 通胀结构分化、AI需求扩散与供给侧约束重塑资产定价逻辑}}
\author{KC桌面 自动生成}
\date{260612}
\begin{document}
\maketitle
\tableofcontents
\newpage
\section{一页摘要}
\begin{itemize}
  \item 5月中国通胀数据呈现PPI上行但核心CPI走弱的“结构性分裂”，上游成本传导与下游定价权断裂并存，市场关注点应从通胀读数转向利润分配格局。
  \item AI驱动的成本传导正在从半导体上游蔓延至消费电子终端，并可能比能源冲击更具持久性，成为重塑中国价格结构的新主线。
  \item 全球石油市场在霍尔木兹海峡中断背景下，需求破坏的规模远超预期，油价定价逻辑从供给瓶颈转向需求脆弱性。
  \item 中国“六网”投资框架加速落地，数据中心与城市更新成为财政扩张新方向，市场对政策执行节奏的定价存在错位。
  \item 全球电动车市场区域利润结构分化加剧，中国新能源渗透率突破60\%后，盈利驱动因素从国内规模转向海外溢价。
  \item AI模型商品化进程加速，竞争焦点从智能竞赛转向成本与信任，中国AI实验室凭借低成本有望占据全球市场35-40\%份额。
  \item 海外铝冶炼产能加速释放，印尼和中亚成为新供给变量，未来三年全球铝市场主要矛盾从中国供给约束转向海外产能扩张。
  \item 中国原油进口骤降并非需求崩塌，而是库存周期主动调节，新能源车替代正在结构性侵蚀汽油需求。
  \item AI驱动的半导体超级周期进入第二阶段，先进制程和HBM产能瓶颈成为未来2-3年更关键的定价因子。
  \item 市场资金流显示科技股拥挤度创历史新高，但通胀高于增长的环境下，资产配置的“看涨惯性”正在积累风险。
\end{itemize}
\section{宏观与利率：通胀结构分裂与政策信号再校准}
\textbf{中国通胀数据呈现PPI上行与核心CPI走弱的‘结构性分裂’，上游成本传导与下游定价权断裂并存，市场需从通胀读数转向利润分配格局。同时，人民币中间价模型显示贬值压力正在向升值压力切换，政策信号值得重新定价。}

\begin{itemize}
  \item 5月PPI同比3.9\%符合预期，但环比从4月的1.7\%骤降至0.5\%，驱动力从能源切换至AI需求与中下游成本传导，DB预计年底PPI升至5\%左右。
  \item 5月CPI同比1.2\%，环比-0.1\%，核心CPI走弱至1.1\%，食品通缩与汽车降价拖累明显，DB将全年CPI预测从1.6\%下调至1.5\%。
  \item BARC指出PPI上游原材料同比上涨9.2\%，采掘业上涨15.8\%，而下游消费品PPI仍在收缩（-0.8\%），定价权断裂导致制造业利润持续承压。
  \item MS认为5月数据确认中国处于‘再通胀退潮’早期，油价脉冲消退后核心CPI将进入温和下行通道，市场低估了这一结构性含义。
  \item Citi强调PPI再通胀正在从‘能源脉冲’切换为‘结构扩散’，剔除能源后PPI仍录得0.3\%环比正增长，投资企稳和AI超周期是新的支撑力量。
  \item NOM人民币中间价模型显示，无逆周期因子下预测值为6.7607，较前次大幅下调543个基点，模型误差收窄暗示贬值压力正在向升值压力切换。
  \item BofA全球资金流数据显示，美国通胀已超4\%，但科技基金单周流入123亿美元创历史新高，牛熊指标连续四周发出‘卖出信号’，市场处于‘看涨惯性’中。
  \item BofA指出1994年情景正在重演——通胀走高、就业紧张、政策滞后，资产配置需警惕‘通胀高于增长’环境下的重新定价风险。
\end{itemize}
\begin{figure}[htbp]
\centering
\includegraphics[width=0.88\linewidth]{figures/fig_001.jpg}
\caption{Figure 1 - Figure 1: PPI reflation was in line with our forecast}
\end{figure}
\begin{figure}[htbp]
\centering
\includegraphics[width=0.88\linewidth]{figures/fig_003.jpg}
\caption{Figure 3 - Figure 3: CPI was unchanged at 1.2\% YoY, while core CPI slowed to 1.1\% YoY}
\end{figure}
\begin{figure}[htbp]
\centering
\includegraphics[width=0.88\linewidth]{figures/fig_005.jpg}
\caption{Figure 5 - Figure 5: Core inflation declined on a MoM basis}
\end{figure}
\begin{figure}[htbp]
\centering
\includegraphics[width=0.88\linewidth]{figures/fig_007.jpg}
\caption{Figure 7 - Figure 7: We revise down the full-year CPI forecast to 1.5\% from 1.6\%}
\end{figure}
\begin{figure}[htbp]
\centering
\includegraphics[width=0.88\linewidth]{figures/fig_009.jpg}
\caption{Exhibit 1 - Exhibit 2: Sequentially oil push slowed, limited reflation elsewhere}
\end{figure}
\begin{figure}[htbp]
\centering
\includegraphics[width=0.88\linewidth]{figures/fig_190.jpg}
\caption{FIGURE 1 - Looking at the breakdown, PPI gains remain highly concentrated in upstream sectors. The producer goods PPI for mining (May: 15.8\% y/y, April: 10.6\%, March: 2.0\%) and raw materials (May: 9.2\%, April: 7.1\%, March: 1.1\%) accelerated, while manufacturing PPI picke}
\end{figure}
\begin{figure}[htbp]
\centering
\includegraphics[width=0.88\linewidth]{figures/fig_192.jpg}
\caption{FIGURE 3 - FIGURE 3. Energy remains a key driver of headline CPI... pp contribution to CPI inflation (\% y/y)}
\end{figure}
\begin{figure}[htbp]
\centering
\includegraphics[width=0.88\linewidth]{figures/fig_193.jpg}
\caption{FIGURE 4 - FIGURE 4. ...while core and services CPI moderated}
\end{figure}
\begin{figure}[htbp]
\centering
\includegraphics[width=0.88\linewidth]{figures/fig_196.jpg}
\caption{Exhibit 1 - Exhibit 1: We expect China's AFD to widen in H2 after narrowing in Q2}
\end{figure}
{\small\color{refgray}\textbf{References:} [R001] DB：市场正从“成本推动的通胀”转向“需求分化的定价游戏”; [R002] 摩根斯坦利：市场低估了“再通胀退潮”的结构性含义; [R003] Citi：PPI再通胀的真正含义不是能源，而是中国工业定价权的结构性拐点; [R026] BARC：通胀数据揭示的真正风险不是物价上涨，而是定价权断裂; [R027] NOM：人民币中间价模型正在揭示一个被低估的政策信号; [R028] NOM：人民币中间价模型正在释放一个被市场低估的政策信号; [R045] 美国银行：市场正在为一场“通胀高于增长”的资产重新定价做准备}

\section{AI/算力：从模型竞赛到成本信任，供应链紧俏重塑定价权}
\textbf{AI产业正经历从‘智能竞赛’到‘成本与信任竞赛’的范式转换，模型商品化加速；同时，AI成本传导从半导体上游蔓延至消费电子终端，供应链紧俏从产能层面升级为定价权转移。}

\begin{itemize}
  \item Bernstein提出AI模型商品化驱动力来自人类用户的‘够用’感知，而非技术收敛；一旦任务达到‘可靠可用’阈值，竞争焦点转向价格、可用性和开发者信任。
  \item Bernstein估算，即使考虑地缘政治限制，中国AI实验室凭借更低token成本，有望占据全球AI市场35-40\%份额。
  \item BofA上调2030年服务器CPU TAM至1700亿美元（此前1250亿），Agentic AI使CPU从‘配角’变为‘总指挥’，AMD、Intel和ARM均受益。
  \item JPM指出4月全球半导体营收同比增长106\%（1994年以来最高），AI驱动的超级周期进入第二阶段，先进制程和HBM产能瓶颈成为未来2-3年关键定价因子。
  \item Citi台湾科技月度追踪显示，台积电5月营收4170亿新台币（同比+30\%），先进制程与封测紧俏从产能层面蔓延至定价层面，全年营收指引有上调空间。
  \item MS台湾科技5月营收数据显示，IC代工集体miss预期（台积电miss 8\%），但数据中心硬件和手机供应链超预期，市场正从‘AI概念普涨’进入‘业绩兑现’阶段。
  \item NOM指出光器件核心原材料磷化铟（InP）衬底进入实质性紧缺，Lumentum一个月内从‘供应可控’转向‘寻找新供应商’，住友电工的‘材料+器件’双重优势被低估。
  \item JPM指出AI驱动的成本传导正在从数据中心建设端向消费电子终端扩散，手机和平板电脑价格环比上涨1.6\%和1.1\%，这一通胀脉冲可能比能源冲击更持久。
\end{itemize}
\begin{figure}[htbp]
\centering
\includegraphics[width=0.88\linewidth]{figures/fig_059.jpg}
\caption{EXHIBIT 2 - EXHIBIT 2: Our perception of AI model quality increasing centres around differentiable quality as perceived by humans, as opposed to underlying reasoning capabilities necessarily converging in a strict academic sense}
\end{figure}
\begin{figure}[htbp]
\centering
\includegraphics[width=0.88\linewidth]{figures/fig_060.jpg}
\caption{Exhibit 4 - Outside this perimeter however, for a growing envelope of lagging edge workloads, especially outside the US, we'd expect the Chinese AI labs' much lower token costs to prove attractive, and help them win share. Exhibit 4 shows directionally how we mentally bre}
\end{figure}
\begin{figure}[htbp]
\centering
\includegraphics[width=0.88\linewidth]{figures/fig_061.jpg}
\caption{EXHIBIT 4 - EXHIBIT 4: We expect the Chinese AI labs to win share on the basis of their much cheaper tokens, but the accessibility of the overseas market will vary by region EXHIBIT 5: We think at least 35-40\% of the long-term global TAM for A}
\end{figure}
\begin{figure}[htbp]
\centering
\includegraphics[width=0.88\linewidth]{figures/fig_161.jpg}
\caption{Exhibit 2 - Exhibit 2: Of the total \textbackslash{}\$2.1Tn CY30 DC Systems TAM, we see AI CPUs representing \textbackslash{}\textasciitilde{}\textbackslash{}\$140bn of the value (split roughly 50/50 across compute/head node and agentic AI node), non-AI CPUs \textbackslash{}\textasciitilde{}\textbackslash{}\$30bn Data Center Systems TAM – Breakout by:}
\end{figure}
\begin{figure}[htbp]
\centering
\includegraphics[width=0.88\linewidth]{figures/fig_162.jpg}
\caption{Exhibit 3 - Exhibit 3: We flag three key roles of CPUs in servers: 1. Traditional, 2. AI compute/head node, 3. AI agentic node Role of the CPU in different server architectures \#\# ROLE OF THE CPU IN DIFFERENT SERVER ARCHITECTURES ```mermaid}
\end{figure}
\begin{figure}[htbp]
\centering
\includegraphics[width=0.88\linewidth]{figures/fig_163.jpg}
\caption{Exhibit 5 - Exhibit 5: We expect INTC/AMD to each represent \textbackslash{}\textasciitilde{}25\% value share, ARM merchant \textbackslash{}\textasciitilde{}35\%, ARM custom \textbackslash{}\textasciitilde{}15\% by CY30 Server CPU TAM (\textbackslash{}\$mn) and Value Market Share Outlook}
\end{figure}
\begin{figure}[htbp]
\centering
\includegraphics[width=0.88\linewidth]{figures/fig_267.jpg}
\caption{Figure 1 - Figure 1: Semi ex-memory revenue YoY (May-24 onwards)}
\end{figure}
\begin{figure}[htbp]
\centering
\includegraphics[width=0.88\linewidth]{figures/fig_268.jpg}
\caption{Figure 3 - Figure 3: Semi revenue (monthly) and YoY}
\end{figure}
\begin{figure}[htbp]
\centering
\includegraphics[width=0.88\linewidth]{figures/fig_278.jpg}
\caption{Exhibit 1 - Exhibit 1: Taiwan Tech: Monthly revenue – May 2026}
\end{figure}
{\small\color{refgray}\textbf{References:} [R012] Bernstein：AI模型商品化进程比市场想象的更快，竞争格局正从智能竞赛转向成本与信任竞赛; [R020] 美国银行：市场低估的不是AI算力需求，而是CPU在Agentic时代的重新定价; [R022] Citi：AI供应链的紧俏不是短期噪音，而是新一轮定价权转移的信号; [R036] JPM：AI驱动的半导体超级周期才刚刚进入第二阶段; [R037] 摩根斯坦利：AI需求的结构性韧性正在重塑台湾科技股的定价逻辑; [R044] NOM：光器件供应链的瓶颈不在产能，而在核心衬底的锁定能力; [R010] JPM：中国通胀的真正信号不是PPI回升，而是AI定价权开始向消费端传导; [R011] JPM：市场真正低估的不是通胀读数，而是AI成本传导正在重塑中国的价格结构}

\section{能源与大宗：需求破坏对冲供给冲击，铝与原油的结构性变局}
\textbf{全球石油市场在霍尔木兹中断背景下，需求破坏规模远超预期，油价定价逻辑从供给瓶颈转向需求脆弱性；铝市场则面临海外产能加速释放，中国供给约束叙事被挑战。}

\begin{itemize}
  \item GS估算中东液体燃料生产二季度受冲击14-15百万桶/日，但全球供需赤字仅5-6百万桶/日，需求损失约5百万桶/日是最大缓冲变量。
  \item GS维持2026Q4布伦特90美元预测，但理由发生根本置换：更小的赤字（5-6百万桶/日）抵消了更长的中断（恢复推迟至8月底），油价地板由需求端重新锚定。
  \item GS将2027年布伦特均价预测下调5美元至80美元，原因包括阿联酋退出OPEC后增产、美洲供给增加以及中国电动车替代加速。
  \item MS指出中国原油进口从2月12.8百万桶/日骤降至5月7.8百万桶/日，但库存数据显示3-4月仍在建库，进口下降是前期过度进口后的主动调节，而非需求崩塌。
  \item MS数据显示中国4月表观石油需求同比下降8\%，LPG和燃料油贡献70\%降幅，新能源车渗透率提升正在结构性侵蚀汽油需求。
  \item GS专家电话会指出，印尼铝冶炼产能扩张条件显著改善，2026-2027年新增产能预期从年初142万吨大幅上调至215万吨和197.5万吨，两年增量超400万吨。
  \item GS指出中国1-4月原铝产量年化已达47.1百万吨，持续超过4500万吨产能红线，违规超产约1.0-1.2百万吨，但近期环保检查属于常规安排。
  \item MS看好铝、铜、锂、铀、金和玻璃纤维，对2026年铝价预测较市场共识高11\%，核心逻辑是供给约束从‘周期性扰动’演变为‘结构性天花板’。
\end{itemize}
\begin{figure}[htbp]
\centering
\includegraphics[width=0.88\linewidth]{figures/fig_119.jpg}
\caption{Exhibit 1 - Exhibit 1: Investor Positioning Moderation Has Been One Driver of the Decline in Oil Prices}
\end{figure}
\begin{figure}[htbp]
\centering
\includegraphics[width=0.88\linewidth]{figures/fig_120.jpg}
\caption{Exhibit 1 - Exhibit 2: We Maintain Our 2026Q4 Brent Oil Price Forecast at \textbackslash{}\$90 But Reduce Our 2027 Average Forecast by \textbackslash{}\$5 to \textbackslash{}\$80}
\end{figure}
\begin{figure}[htbp]
\centering
\includegraphics[width=0.88\linewidth]{figures/fig_123.jpg}
\caption{Exhibit 5 - Exhibit 5: We Estimate a 5-6mb/d Q2 Deficit, Which Is Smaller Than the 14-15mb/d Hit to Mideast Production Reconciling 14-15mb/d Mideast Supply Loss With 5-6mb/d 2026Q2 Global Deficit}
\end{figure}
\begin{figure}[htbp]
\centering
\includegraphics[width=0.88\linewidth]{figures/fig_128.jpg}
\caption{Exhibit 10 - Exhibit 10: We Have Revised Higher Our 2027 Global Oil Surplus to 3.5mb/d Given Our Expectation of Higher Production in the UAE and Americas GS 2027 Global Oil Surplus}
\end{figure}
\begin{figure}[htbp]
\centering
\includegraphics[width=0.88\linewidth]{figures/fig_065.jpg}
\caption{Exhibit 1 - Exhibit 1: Crude net imports fell 2.4 mb/d MoM in April to 9.3 mb/d. Net imports then declined to 7.8 mb/d in May, down 3.2 mb/d YoY.}
\end{figure}
\begin{figure}[htbp]
\centering
\includegraphics[width=0.88\linewidth]{figures/fig_066.jpg}
\caption{Exhibit 2 - Exhibit 2 : Observable Chinese crude inventories continued to build over March and April, based on Vortexa data. Average build rate was \textbackslash{}\textasciitilde{}1.1 mb/d. However, inventories switched to draws in May, with an average draw rate of 750 kb/d}
\end{figure}
\begin{figure}[htbp]
\centering
\includegraphics[width=0.88\linewidth]{figures/fig_067.jpg}
\caption{Exhibit 3 - Exhibit 3: China's apparent oil demand fell 1.4 mb/d MoM and was 1.2 mb/d lower YoY (-8\%) in April China Apparent Oil Demand (mb/d)}
\end{figure}
\begin{figure}[htbp]
\centering
\includegraphics[width=0.88\linewidth]{figures/fig_063.jpg}
\caption{Exhibit 1 - Exhibit 1: Monthly China aluminum ouptut run rate reached 47.1mnt in April 2026, 1.8 mnt higher than the average outout of 45.3mnt in 2025A China primary aluminum output (annualized, mmt)}
\end{figure}
\begin{figure}[htbp]
\centering
\includegraphics[width=0.88\linewidth]{figures/fig_064.jpg}
\caption{Exhibit 2 - Net addition of aluminum capacity/output in Indonesia}
\end{figure}
{\small\color{refgray}\textbf{References:} [R016] GS：市场低估的不是地缘冲突的烈度，而是需求破坏的持久性; [R014] 摩根斯坦利：中国原油进口的骤降并非需求崩塌，而是库存周期的主动调节; [R013] GS：铝市场真正超预期的不是中国产量，而是海外供给的加速释放; [R024] 摩根斯坦利：市场低估了供给侧的“结构性断裂”，而非仅仅需求回暖}

\section{地产与消费：分化中的结构性机会与成本端出清}
\textbf{房地产信用分化加速，市场从总量博弈转向个券alpha挖掘；消费行业正经历成本端的分化式出清，原奶价格企稳和IP零售‘高基数疲劳’是两大关键信号。}

\begin{itemize}
  \item BofA指出6月香港一手房成交仅约240宗（前5月月均超2000宗），但放缓主因是开发商主动控制推盘节奏，而非需求坍塌，离岸投资规则不确定性是核心变量。
  \item BofA认为中国高收益地产债收益率上行18bp至9.58\%，但信用主体分化加速，个别香港地产公司因债务重组空间和资产出售进展获得偏乐观看法。
  \item BofA指出中央财政安排970亿中央预算和1600亿超长期特别国债支持城市更新，但对比每年超3万亿的总投资，增量影响相对有限。
  \item MS大宗商品价格图鉴显示，原奶价格5月环比微涨0.2\%，供给削减后趋于稳定，预计2026年逐步回升，伊利和蒙牛有望凭借成本优势抢占中小品牌份额。
  \item MS指出原奶价格企稳意味着存货拨备压力大幅减少，净利率修复弹性可能超预期，成本上行期将加速中小乳企退出。
  \item GS指出鲜零食零售的兴起验证了零食消费韧性，但头部价值零售商（鸣鸣很忙、万辰）凭借网络密度和供应链优势，反而获得品类升级和店型升级的催化剂。
  \item GS指出IP零售核心玩家正经历‘高基数疲劳’，泡泡玛特国内线上增速从4月95\%骤降至5月29\%，名创优品从43\%降至21\%，增长引擎从流量红利切换至IP运营能力。
  \item GS数据显示泡泡玛特美国市场5月信用卡销售同比下降36\%，海外扩张故事中基数效应和消费者情绪波动成为不可忽视的变量。
\end{itemize}
\begin{figure}[htbp]
\centering
\includegraphics[width=0.88\linewidth]{figures/fig_148.jpg}
\caption{Exhibit 2 - Exhibit 2: China Property Developers' bond yields vs. share prices The average bond yield widened 18bp WoW to \$9.58\textbackslash{}\%\$ , and the share price \$-0.9\textbackslash{}\%\$ WoW.}
\end{figure}
\begin{figure}[htbp]
\centering
\includegraphics[width=0.88\linewidth]{figures/fig_149.jpg}
\caption{Exhibit 3 - Exhibit 3: Hong Kong property issuers' bond spreads vs. share prices The average bond spread (OAS) tightened 3bp WoW to 61bp and the share price -7.2\% WoW.}
\end{figure}
\begin{figure}[htbp]
\centering
\includegraphics[width=0.88\linewidth]{figures/fig_150.jpg}
\caption{Exhibit 4 - Exhibit 4: YTM (\%) of benchmark bonds for select Chinese HY property developers The average yield of benchmark bonds widened 10bp WoW on average basis.}
\end{figure}
\begin{figure}[htbp]
\centering
\includegraphics[width=0.88\linewidth]{figures/fig_310.jpg}
\caption{Exhibit 35 - Exhibit 35: Raw Milk Price (Monthly)}
\end{figure}
\begin{figure}[htbp]
\centering
\includegraphics[width=0.88\linewidth]{figures/fig_251.jpg}
\caption{Exhibit 1 - Exhibit 1: Pop Mart's online sales growth in May and 2Q26-to-date growth were both slower sequentially based on the data we track Pop Mart China online sales growth}
\end{figure}
\begin{figure}[htbp]
\centering
\includegraphics[width=0.88\linewidth]{figures/fig_257.jpg}
\caption{Exhibit 7 - Exhibit 7: Pop Mart US credit card sales decelerated sequentially in 2Q26-to-date Pop Mart US credit card sales (quarterly)}
\end{figure}
\begin{figure}[htbp]
\centering
\includegraphics[width=0.88\linewidth]{figures/fig_262.jpg}
\caption{Exhibit 15 - Exhibit 15: Miniso US credit card sales growth re-accelerated to \$+31\textbackslash{}\%\$ in May from \$+19\textbackslash{}\%\$ yoy in Apr}
\end{figure}
{\small\color{refgray}\textbf{References:} [R018] 美国银行：市场低估了房地产信用分化中的结构性机会，而非整体复苏; [R038] 摩根斯坦利：大宗商品周期的分化正在重塑中国消费股的盈利格局; [R034] GS：鲜零食零售的兴起不是威胁，而是价值零售龙头品类升级的催化剂; [R035] GS：IP零售的“高基数疲劳”正在重塑市场对增长的定价逻辑}

\section{区域市场：欧元弱势与人民币汇率信号切换}
\textbf{欧元面临技术面破位、基本面背离和历史模式重复三重压力，弱势远未结束；人民币中间价模型则释放贬值压力向升值压力切换的信号，政策框架正在系统性再校准。}

\begin{itemize}
  \item BofA指出欧元在6月5日非农数据后跌破上升趋势线支撑，周线图正在构建头肩顶形态，颈线在1.1411-1.1392，若跌破下行目标指向1.11甚至1.0980。
  \item BofA发现美元指数在特朗普第二任期走势与第一任期高度相似，若2018年剧本重演，美元指数可能向103-105区间上行。
  \item BofA认为美欧增长差被低估，美国就业市场韧性超预期，市场对美联储加息的定价不够充分，而欧洲加息因经济疲弱难以持续。
  \item NOM人民币中间价模型显示，无逆周期因子下预测值为6.7607，较前次大幅下调543个基点，模型误差从年初-1800点收窄至接近0，预测框架在适应新环境。
  \item NOM指出模型投影下降由多因子驱动，欧元和澳元贡献正向变动，韩元反向拖累，中间价正从‘盯住美元’转向‘盯住一篮子货币’的精细化操作。
  \item NOM列出关键时间节点：7月底政治局会议、10月黄金周、11月APEC、12月中央经济工作会议，这些事件可能影响汇率政策方向。
\end{itemize}
\begin{figure}[htbp]
\centering
\includegraphics[width=0.88\linewidth]{figures/fig_053.jpg}
\caption{Exhibit 1 - Exhibit 1: USD still trading within its previous 12m range, while labor surprise measures reach fresh multi-year highs DXY \& Bloomberg US Labor Surprise Index}
\end{figure}
\begin{figure}[htbp]
\centering
\includegraphics[width=0.88\linewidth]{figures/fig_054.jpg}
\caption{Exhibit 3 - Exhibit 3: Policy differentials now priced for modest but steady narrowing (towards EUR)... Fed \& ECB policy rate}
\end{figure}
\begin{figure}[htbp]
\centering
\includegraphics[width=0.88\linewidth]{figures/fig_056.jpg}
\caption{Exhibit 2 - Exhibit 2: Relative US vs. Euro area data trends have been stark, suggesting further space for the EUR to depreciate further Economic Change Indices: EA - USA}
\end{figure}
\begin{figure}[htbp]
\centering
\includegraphics[width=0.88\linewidth]{figures/fig_057.jpg}
\caption{Exhibit 4 - Exhibit 4: ... Yet post-“Liberation Day” structural break between FX and policy rates has continued ECB-Fed policy differential (and OIS Implied) vs. EUR}
\end{figure}
{\small\color{refgray}\textbf{References:} [R009] 美国银行：欧元弱势远未结束，市场低估了技术面与基本面的共振; [R027] NOM：人民币中间价模型正在揭示一个被低估的政策信号; [R028] NOM：人民币中间价模型正在释放一个被市场低估的政策信号}

\section{风险偏好：资金流拥挤与‘通胀高于增长’的再定价}
\textbf{全球资金流显示科技股拥挤度创历史新高，但通胀高于增长的环境下，资产配置的‘看涨惯性’正在积累风险，历史类比指向剧烈调整的可能性。}

\begin{itemize}
  \item BofA数据显示上周全球股票基金净流入315亿美元，美国股票连续11周净流入，科技板块独占123亿美元创历史纪录，但牛熊指标已连续四周发出‘卖出信号’。
  \item BofA指出美国通胀已超4\%，特朗普在通胀问题上支持率跌至27\%（低于拜登任内低点），政策可能从降息转向加息，市场对‘3P’（持仓、盈利预期、政策）风险定价不足。
  \item BofA指出1994年情景正在重演：通胀>4\%时标普500未来3个月平均跌4\%，6个月跌7\%，当时美联储被迫大幅加息导致债券收益率飙升和股市多月的交易区间。
  \item Citi全球主题排名体系显示，排名驱动力从‘价格动量’转向‘盈利可见性’与‘质量因子’，云计算、可穿戴技术和AI赋能者领跑，奢侈品和太阳能排名下滑。
  \item Citi指出市场正在重新定价‘增长的质量’而非增长本身，同一AI赛道内部公司命运分化才刚刚开始。
  \item BofA认为如果地缘冲突缓和，消费股、REITs、欧洲和新兴市场货币可能受益，当前被资金抛弃但基本面改善的角落存在逆向机会。
\end{itemize}
\begin{figure}[htbp]
\centering
\includegraphics[width=0.88\linewidth]{figures/fig_016.jpg}
\caption{Figure 1 - As part of our annual update cycle, we have recently completed a comprehensive thematic reassessment and stock remapping, reflected in our latest publication, Global Theme Machine: Reassessment of Themes for 2026. Leveraging tools refined over many years, we a}
\end{figure}
\begin{figure}[htbp]
\centering
\includegraphics[width=0.88\linewidth]{figures/fig_018.jpg}
\caption{Figure 4 - © 2026 Citi Inc. No redistribution without Citi's written permission. \#\# Model Performance Top quintile themes rose \$3.6\textbackslash{}\%\$ in the latest month, modestly ahead of the \$3.4\textbackslash{}\%\$ return from bottom quintile themes, although both lagged the MSCI World's \$4.6\textbackslash{}\%\$ gai}
\end{figure}
\begin{figure}[htbp]
\centering
\includegraphics[width=0.88\linewidth]{figures/fig_020.jpg}
\caption{Figure 8 - Figure 8. Bottom Performing Themes in Last Month (red denotes the bottom performing theme for the corresponding period) © 2026 Citi Inc. No redistribution without Citi's written permission. Figure 9. Top 10 \& Bottom 10 Themes by Last month's Returns}
\end{figure}
\begin{figure}[htbp]
\centering
\includegraphics[width=0.88\linewidth]{figures/fig_021.jpg}
\caption{Figure 11 - © 2026 Citi Inc. No redistribution without Citi's written permission. Figure 11. Top 10 \& Bottom 10 Themes by Last 12 months' Returns}
\end{figure}
{\small\color{refgray}\textbf{References:} [R045] 美国银行：市场正在为一场“通胀高于增长”的资产重新定价做准备; [R005] Citi：市场正在重新定价“增长的质量”，而非增长本身}

\section{电动车与汽车：区域利润结构重构与燃油车定价崩塌}
\textbf{全球电动车市场区域利润结构分化加剧，中国新能源渗透率突破60\%后盈利驱动因素转向海外溢价；燃油车库存定价崩塌正在形成自我强化的下行螺旋。}

\begin{itemize}
  \item JPM数据显示5月中国新能源乘用车零售84.9万辆，渗透率63\%创历史新高，燃油车零售57.9万辆同比暴跌38\%，消费者在新能源车更贵的情况下仍加速转向。
  \item JPM指出零里程二手燃油车数量正在上升，燃油车新车价格和销量之间形成自我强化的下行螺旋，保值率对新能源的优势正快速消失。
  \item JPM数据显示5月中国乘用车出口80.9万辆（同比+73\%）创历史新高，其中新能源出口43.5万辆（+110\%），巴西和澳大利亚成为新重点。
  \item JPM指出欧洲五国5月EV+PHEV渗透率同比跳升8个百分点至32\%，法国表现最猛（+13ppt），补贴政策和平价车型是主要拉动力量。
  \item JPM数据显示美国EV渗透率仅7\%同比下滑2个百分点，车企砍新车型和收缩营销导致消费者选择减少，形成负反馈循环。
  \item JPM指出储能系统成为电池产业链中定价权最集中的环节，全球4M26出货量同比翻倍以上，AIDC储能正在成为新变量。
  \item Citi数据显示5月新能源车CR5为57.0\%，环比微增但同比下降2.0个百分点，龙头集中度开始松动，比亚迪份额27.9\%环比+2.2ppt但新势力分化明显。
  \item Citi指出渗透率突破60\%后增长引擎从‘增量红利’切换为‘存量争夺’，纯电增速（+17\%）明显好于插混（+1\%），B级车市占率持续提升。
\end{itemize}
\begin{figure}[htbp]
\centering
\includegraphics[width=0.88\linewidth]{figures/fig_033.jpg}
\caption{Figure 3 - Figure 3: EU10 + US + China EV/PHEV sales and penetration trend}
\end{figure}
\begin{figure}[htbp]
\centering
\includegraphics[width=0.88\linewidth]{figures/fig_034.jpg}
\caption{Figure 5 - Figure 5: EU-5+US+China EV/PHEV sales mix trend}
\end{figure}
\begin{figure}[htbp]
\centering
\includegraphics[width=0.88\linewidth]{figures/fig_035.jpg}
\caption{Figure 7 - Figure 7: EU-5 EV/PHEV penetration trend}
\end{figure}
\begin{figure}[htbp]
\centering
\includegraphics[width=0.88\linewidth]{figures/fig_036.jpg}
\caption{Figure 4 - Figure 4: EU10 + US + China EV sales and penetration trend}
\end{figure}
\begin{figure}[htbp]
\centering
\includegraphics[width=0.88\linewidth]{figures/fig_037.jpg}
\caption{Figure 6 - Figure 6: US EV/PHEV penetration trend}
\end{figure}
\begin{figure}[htbp]
\centering
\includegraphics[width=0.88\linewidth]{figures/fig_038.jpg}
\caption{Figure 8 - Figure 8: China EV/PHEV penetration trend}
\end{figure}
\begin{figure}[htbp]
\centering
\includegraphics[width=0.88\linewidth]{figures/fig_180.jpg}
\caption{Figure 1 - Figure 1: China Passenger Car Monthly Wholesale Volume 1,000 units, \%}
\end{figure}
\begin{figure}[htbp]
\centering
\includegraphics[width=0.88\linewidth]{figures/fig_247.jpg}
\caption{Figure 1 - Figure 1. NEV Wholesale Volume (units) by Major OEMs © 2026 Citi Inc. No redistribution without Citi's written permission. Figure 2. NEV-PV Monthly Sales Volume and YoY Change}
\end{figure}
\begin{figure}[htbp]
\centering
\includegraphics[width=0.88\linewidth]{figures/fig_248.jpg}
\caption{Figure 3 - © 2026 Citi Inc. No redistribution without Citi's written permission. Figure 3. BEV vs PHEV growth rate}
\end{figure}
{\small\color{refgray}\textbf{References:} [R008] JPM：全球电动车市场的真正分歧不在渗透率，而在区域利润结构; [R023] JPM：中国汽车市场真正的拐点不是电动化渗透率，而是燃油车库存的定价崩塌; [R033] Citi：中国新能源车市真正的拐点不是渗透率破60\%，而是龙头集中度开始松动}

\section{中国政策与投资：‘六网’落地与资本管控信号再解读}
\textbf{中国‘六网’投资框架加速落地，数据中心与城市更新成为财政扩张新方向，市场对政策执行节奏的定价存在错位；资本管控新规核心目标是保护企业出海而非限制个人资金。}

\begin{itemize}
  \item Citi指出三峡新通道（772亿）和数据中心建设规划（2万亿）是‘六网’落地的标志性项目，国家发改委承诺今年‘六网’投资规模超7万亿元，市场定价不足。
  \item Citi测算2万亿数据中心投资若加上电网配套，总规模至少5万亿，融资路径包括超长期国债、政策性金融工具（今年5000亿）和政府引导基金。
  \item GS指出‘2万亿数据中心投资计划’并非新消息（3月‘十五五’规划已纳入），真正新信息是政策沟通频率和项目准备工作的明显提速。
  \item GS观察到5月底以来政策性银行新融资工具加速落地，二季度财政支出放缓后政策制定者正从‘等待数据’转向‘提前准备’。
  \item MS指出国务院新规发布会上只有司法部、发改委和商务部出席，无金融监管机构参与，政策目标群体是进行海外实业投资的企业而非个人。
  \item MS认为个人5万美元换汇额度未变，家庭金融资产增速>10\%但收入仅增4-5\%，数据不支持大规模资本外逃，政策本质是制度升级而非方向逆转。
  \item Citi指出四月固定资产投资增速转负是结构转换期典型特征，旧动能（房地产）收缩速度超过新动能（新基建、制造业升级）扩张速度。
  \item Citi认为五月建筑业PMI反弹和水泥开工率改善是复苏初现信号，但整体固投回升需要时间，大概率是K型复苏——新经济投资领跑。
\end{itemize}
\begin{figure}[htbp]
\centering
\includegraphics[width=0.88\linewidth]{figures/fig_015.jpg}
\caption{Figure 1 - Figure 1. State Council has set explicit targets for urban renewal in its latest guidance © 2026 Citi Inc. No redistribution without Citi's written permission. Figure 2. IT FAI has outpaced headline FAI and stayed buoyant in April FAI, Headline and IT}
\end{figure}
\begin{figure}[htbp]
\centering
\includegraphics[width=0.88\linewidth]{figures/fig_196.jpg}
\caption{Exhibit 1 - Exhibit 1: We expect China's AFD to widen in H2 after narrowing in Q2}
\end{figure}
\begin{figure}[htbp]
\centering
\includegraphics[width=0.88\linewidth]{figures/fig_197.jpg}
\caption{Exhibit 1 - Exhibit 1: Household financial asset growth remained above 10\%... China household financial assets, Rmb bn}
\end{figure}
\begin{figure}[htbp]
\centering
\includegraphics[width=0.88\linewidth]{figures/fig_198.jpg}
\caption{Exhibit 3 - Exhibit 3: The PBOC withdrew liquidity over the past three months Monthly PBoC Liquidity Injection/(withdrawal), Rmb mn}
\end{figure}
{\small\color{refgray}\textbf{References:} [R004] Citi：市场低估了“六网”投资对中国固定资产投资的托底力度; [R029] GS：市场真正低估的不是2万亿数据中心投资，而是“六张网”的宏观含义; [R030] 摩根斯坦利：市场误读了中国的资本管控信号，真正焦点在实体经济出海而非个人资金外逃}

\section{工业与材料：供给约束与新兴需求驱动的再定价}
\textbf{中国材料行业正经历由供给约束（产能天花板、ESG成本、地缘风险）与新兴需求（储能、AI）共同驱动的再定价，市场对供给弹性的评估不足。}

\begin{itemize}
  \item MS指出中国电解铝产能逼近4500万吨政策天花板，海外能源成本上升抑制复产意愿，对2026年铝价预测较市场共识高11\%，2027年高13\%。
  \item MS看好铜、锂、铀、黄金和玻璃纤维，中国铜消费指数显示电力（占比47\%）是绝对主力，预测2026年电力用铜增速8\%高于2025年的5\%。
  \item MS对2026年碳酸锂价格预测22840美元/吨，较市场共识高17\%，但认为长期价格会回落；黄金预测4916美元/盎司，铀价93.4美元/磅。
  \item MS指出印度钢铁强势周期尚未结束，国内热轧卷价格从12月低点反弹26\%后回落11\%（雨季因素），保障性关税延期和中国‘反内卷’是供给侧双重锁定。
  \item MS数据显示印度钢铁股年初至今上涨约12\%，而Sensex下跌13\%，上调多家钢铁公司盈利预测，看好JSW Steel、Jindal Steel和Tata Steel。
  \item MS指出中国聚乙烯净进口量才是美国合同价格真正的定价锚，4月中国净进口从110万吨骤降至18万吨，向全球释放约9\%额外供应，几乎对冲中东供应中断。
  \item MS推测中国在冲突前积累了约230-350万吨过剩聚乙烯库存（相当于30天国内需求），若6-7月消耗完恢复进口节奏，每月可能额外需要约25万吨美国出口。
  \item Citi指出光伏国内价格仍在低位但最剧烈下跌阶段已过去，出口端框架协议大规模落地成为新亮点，隆基、晶科等与中东、非洲买家签GW级协议。
\end{itemize}
\begin{figure}[htbp]
\centering
\includegraphics[width=0.88\linewidth]{figures/fig_063.jpg}
\caption{Exhibit 1 - Exhibit 1: Monthly China aluminum ouptut run rate reached 47.1mnt in April 2026, 1.8 mnt higher than the average outout of 45.3mnt in 2025A China primary aluminum output (annualized, mmt)}
\end{figure}
\begin{figure}[htbp]
\centering
\includegraphics[width=0.88\linewidth]{figures/fig_064.jpg}
\caption{Exhibit 2 - Net addition of aluminum capacity/output in Indonesia}
\end{figure}
\begin{figure}[htbp]
\centering
\includegraphics[width=0.88\linewidth]{figures/fig_181.jpg}
\caption{Exhibit 1 - Exhibit 1: China Net Import Reliance for PE China Net Import Reliance for PE}
\end{figure}
\begin{figure}[htbp]
\centering
\includegraphics[width=0.88\linewidth]{figures/fig_182.jpg}
\caption{Exhibit 2 - Exhibit 2: MSe 2026 US PE contract expectations, alongside (i) a scenario where incremental US exports drive higher prices and (ii) another where the reversion to a lower mean is accelerated}
\end{figure}
\begin{figure}[htbp]
\centering
\includegraphics[width=0.88\linewidth]{figures/fig_183.jpg}
\caption{Exhibit 3 - Exhibit 3: Sinopec and Petrochina Polyolefin Inventories (PE \& PP)}
\end{figure}
\begin{figure}[htbp]
\centering
\includegraphics[width=0.88\linewidth]{figures/fig_368.jpg}
\caption{Exhibit 2 - Exhibit 2: India: Expect domestic production growth to pick up as new capacities ramp up through the year; demand likely to remain good}
\end{figure}
\begin{figure}[htbp]
\centering
\includegraphics[width=0.88\linewidth]{figures/fig_369.jpg}
\caption{Exhibit 3 - Exhibit 3: Gross exports of finished steel picked up recently, but should normalize as demand picks up beyond monsoons (million tons)}
\end{figure}
\begin{figure}[htbp]
\centering
\includegraphics[width=0.88\linewidth]{figures/fig_371.jpg}
\caption{Exhibit 5 - Exhibit 5: Finished steel inventory (non-alloy) on an absolute basis has come off (mnt)...}
\end{figure}
\begin{figure}[htbp]
\centering
\includegraphics[width=0.88\linewidth]{figures/fig_175.jpg}
\caption{Figure 1 - Figure 1. PRC weekly polysilicon prices in the week ended Jun 10 © 2026 Citi Inc. No redistribution without Citi's written permission. Figure 2. PRC weekly polysilicon price}
\end{figure}
\begin{figure}[htbp]
\centering
\includegraphics[width=0.88\linewidth]{figures/fig_176.jpg}
\caption{Figure 3 - © 2026 Citi Inc. No redistribution without Citi's written permission. Figure 3. PRC weekly wafer price}
\end{figure}
{\small\color{refgray}\textbf{References:} [R024] 摩根斯坦利：市场低估了供给侧的“结构性断裂”，而非仅仅需求回暖; [R043] 摩根斯坦利：印度钢铁的强势周期远未结束，但真正值得关注的是供给侧的再定价; [R025] 摩根斯坦利：中国聚乙烯净进口量才是美国合同价格真正的定价锚; [R021] Citi：光伏市场的真正拐点不在国内价格战，而在海外框架协议的规模与结构}

\section{生物医药与健康：中国创新资产需求刚性，阿尔茨海默病成本曲线待重构}
\textbf{全球药企对中国创新资产的结构性需求远超美国政策噪音，地缘风险更多是情绪层面的；阿尔茨海默病正成为人口老龄化下最集中的成本放大器，‘脑穿梭’技术或重构治疗路径。}

\begin{itemize}
  \item JPM美国分析师团队明确区分‘影响投资者情绪’和‘改变全球药物开发轨迹’，认为美国国会信函和药明康德清单属于前者而非后者，是‘噪音’而非结构性障碍。
  \item JPM调查显示60\%以上投资者认为今年并购会比去年多，大药企弹药充足重点在肿瘤、免疫、心血管代谢和神经科学，但买方对估值和风险分担更严格。
  \item JPM指出PD-1/VEGF等IO资产临床价值太大，监管方若仅因数据来自中国就阻止美国患者使用将面临巨大压力，康方生物依沃西单抗是典型案例。
  \item MS估算当前美国阿尔茨海默病相关市场规模约700亿美元，基于约700万患者，疾病修饰疗法约500亿，症状控制约200亿。
  \item MS指出美国人口2050年见顶后缓慢下降，但老年人比例从18.8\%升至30.5\%，仅年龄结构变化就能将个人医疗支出占GDP比重从11.7\%推至13.1\%。
  \item MS认为口服司美格鲁肽在早期阿尔茨海默的3期试验失败，代谢通路和神经退行之间的鸿沟比想象更深，‘脑穿梭’技术平台能否重构成本曲线才是真正拐点。
  \item MS指出Lecanemab和Donanemab清除淀粉样蛋白效果已接近极限，临床获益约27\%延缓（18个月CDR-SB评分约0.45分），但不够。
\end{itemize}
\begin{figure}[htbp]
\centering
\includegraphics[width=0.88\linewidth]{figures/fig_384.jpg}
\caption{Exhibit 2 - Exhibit 2: Longevity of the Brain: Alzheimer's, the GLP-1 Boundary, and the Next Spend Curve ```mermaid graph LR}
\end{figure}
\begin{figure}[htbp]
\centering
\includegraphics[width=0.88\linewidth]{figures/fig_385.jpg}
\caption{Exhibit 6 - Exhibit 6: CBO US Population forecast Exhibit 7: CBO US Population forecast}
\end{figure}
\begin{figure}[htbp]
\centering
\includegraphics[width=0.88\linewidth]{figures/fig_386.jpg}
\caption{Exhibit 9 - Exhibit 8: Average PHC spending age segment}
\end{figure}
\begin{figure}[htbp]
\centering
\includegraphics[width=0.88\linewidth]{figures/fig_387.jpg}
\caption{Exhibit 9 - Exhibit 9: Seniors drive PHC spend}
\end{figure}
\begin{figure}[htbp]
\centering
\includegraphics[width=0.88\linewidth]{figures/fig_388.jpg}
\caption{Exhibit 12 - Exhibit 12: Care for Seniors with dementia is meaningfully higher than Seniors without dementia. Dementia is the deciding cost driver.}
\end{figure}
\begin{figure}[htbp]
\centering
\includegraphics[width=0.88\linewidth]{figures/fig_389.jpg}
\caption{Exhibit 13 - Exhibit 13: Healthspan (years without severe impairment) ends in the last decade of life Global Lifespan vs. Healthspan (High income countries)}
\end{figure}
\begin{figure}[htbp]
\centering
\includegraphics[width=0.88\linewidth]{figures/fig_390.jpg}
\caption{Exhibit 15 - Exhibit 14: AD prevalence continues to increase, and is estimated to surpass 10M US adults by 2035 AD \& MCI US Prevalence Forecast}
\end{figure}
\begin{figure}[htbp]
\centering
\includegraphics[width=0.88\linewidth]{figures/fig_391.jpg}
\caption{Exhibit 15 - Exhibit 15: AD prevalence, gross spend, and per capita spend (MSe vs Alzheimer's association forecast)}
\end{figure}
{\small\color{refgray}\textbf{References:} [R007] JPM：市场低估了全球药企对中国创新资产的刚性需求，而非地缘风险; [R046] 摩根斯坦利：阿尔茨海默病的真正拐点不在新药，而在“脑穿梭”能否重构成本曲线}

\section{新兴市场专题：印度电力时段性短缺与日本半导体设备份额反转}
\textbf{印度电力短缺的真正瓶颈从总量转向时段，非日照时段供给刚性未来三年几乎无法缓解；日本半导体设备在中国市场份额流失主要由汇率和订单节奏驱动，反转正在酝酿。}

\begin{itemize}
  \item Bernstein指出印度电力现货市场总量数据看似宽松（均价3.8卢比/千瓦时），但拆解时段后早晨与晚间价差已扩大至7卢比/千瓦时（历史接近0），晚间买卖比达2.2倍。
  \item Bernstein预计未来3年仅新增15GW火电，而傍晚需求将增长50GW，除非每年加15GW电池储能否则缺口无解，去年大量BESS招标可能因电池价格和汇率不利无法落地。
  \item Bernstein指出印度电力需求去年仅增1\%，今年预计6\%已算保守，降雨预测低于正常叠加数据中心和电气化结构性需求，下半年可能继续走强。
  \item Bernstein指出日本半导体设备在中国进口市场名义份额从26\%降至23\%，但经汇率调整后2024年实际份额为31\%，2025年为26\%，与2021-2023年均值28.1\%几乎持平。
  \item Bernstein认为日元贬值31\%让日本设备性价比突出，中国客户正重新增加对日订单，TEL已明确表示在中国有涨价空间，Screen预计CXMT订单2年后回归。
  \item Bernstein指出过去2-3个月半导体设备跑输存储/模拟等周期股，但未来几年资本支出上行周期中日本设备商有望从2026年下半年起重新收复失地。
\end{itemize}
\begin{figure}[htbp]
\centering
\includegraphics[width=0.88\linewidth]{figures/fig_202.jpg}
\caption{Exhibit 1 - Uma Menon Aman Jain We see the Indian power sector a structural and not a cyclical story, but the sector has traded more like air-conditioner stocks - buy before summer and sell before onset of monsoons. The sector peaked in mid-CY24, post which demand slowed }
\end{figure}
\begin{figure}[htbp]
\centering
\includegraphics[width=0.88\linewidth]{figures/fig_204.jpg}
\caption{EXHIBIT 4 - EXHIBIT 4: But on disaggregating by time slot, we continue to see the evening situation being tight in the power sector}
\end{figure}
\begin{figure}[htbp]
\centering
\includegraphics[width=0.88\linewidth]{figures/fig_205.jpg}
\caption{EXHIBIT 5 - EXHIBIT 5: Spot power prices: While overall spot prices have come down, evening continues to be high, and also the day-evening arbitrage is increasing}
\end{figure}
\begin{figure}[htbp]
\centering
\includegraphics[width=0.88\linewidth]{figures/fig_206.jpg}
\caption{EXHIBIT 6 - EXHIBIT 6: Power demand growth vs. GDP growth: FY26 saw another year of moderated demand growth, however FY27 demand so far looks strong. We retain our 1x real gdp long term view for power demand growth}
\end{figure}
\begin{figure}[htbp]
\centering
\includegraphics[width=0.88\linewidth]{figures/fig_207.jpg}
\caption{EXHIBIT 7 - EXHIBIT 7: Power demand: May and June 2026 saw over 10\% power demand YoY growth benefitting from a low base last year Power demand YoY growth by month (\%)}
\end{figure}
\begin{figure}[htbp]
\centering
\includegraphics[width=0.88\linewidth]{figures/fig_213.jpg}
\caption{Exhibit 14 - EXHIBIT 13: Limited thermal capacity additions next 3 yrs}
\end{figure}
\begin{figure}[htbp]
\centering
\includegraphics[width=0.88\linewidth]{figures/fig_214.jpg}
\caption{EXHIBIT 14 - EXHIBIT 14: Power demand-supply: We see power shortages continue next 3 years, and a need for 15 GW/yr of BESS to reduce the gap... India - Peak Demand-Supply (GW)}
\end{figure}
\begin{figure}[htbp]
\centering
\includegraphics[width=0.88\linewidth]{figures/fig_219.jpg}
\caption{Exhibit 1 - EXHIBIT 1: WFE imports to China has been on a constant rise... Total China WFE imports (USD bn)}
\end{figure}
\begin{figure}[htbp]
\centering
\includegraphics[width=0.88\linewidth]{figures/fig_220.jpg}
\caption{EXHIBIT 2 - EXHIBIT 2: ...but Japan has been constantly dropping market share since 2022.}
\end{figure}
\begin{figure}[htbp]
\centering
\includegraphics[width=0.88\linewidth]{figures/fig_221.jpg}
\caption{Exhibit 3 - EXHIBIT 3: JPY has depreciated 31.5\% over the last 5 years.}
\end{figure}
\begin{figure}[htbp]
\centering
\includegraphics[width=0.88\linewidth]{figures/fig_222.jpg}
\caption{EXHIBIT 4 - EXHIBIT 4: FX may have some impact on Japan's market share loss as they charge customers mostly in JPY.}
\end{figure}
{\small\color{refgray}\textbf{References:} [R031] Bernstein：市场低估的不是电力需求，而是非日照时段的供给刚性; [R032] Bernstein：日本半导体设备在中国市场的份额流失并非结构性，反转正在酝酿}

\section{固定收益与信用：CLO套利空间收窄与信用分化中的alpha机会}
\textbf{美元CLO套利空间出现结构性收窄，‘自有股权资本’崛起改变了投资者行为逻辑；中国房地产信用分化加速，市场从总量博弈转向个券alpha挖掘。}

\begin{itemize}
  \item GS构建的交易匹配套利指标显示，美元CLO套利空间过去两年收窄幅度远超传统指标，并非周期性现象而是结构性变化，根源在于‘自有股权资本’的增长。
  \item GS指出‘自有股权资本’追求连续部署和年份分散化，即使即期套利空间收窄也不会轻易退场，降低了股权需求对利差变化的敏感度。
  \item GS发现美元市场高评级经理反而能维持更宽套利空间（负债端优势），欧元市场则相反低评级经理套利更宽（资产端差异），区域基本面偏弱使欧元超额利差可能只是信用损失补偿。
  \item BofA指出中国高收益地产债收益率上行18bp至9.58\%，但年初至今仍收窄162bp，香港投资级地产债利差收窄3bp至61bp。
  \item BofA认为政策子弹边际效用递减，但信用主体分化加速，过去两年‘买整个板块’的交易策略正在失效，取而代之的是更依赖主体信用判断的配置逻辑。
  \item BofA对个别香港地产公司持偏乐观看法，理由包括债务重组条款仍有空间、利息覆盖约1倍、资产出售持续推进去杠杆、公司主动卖资产态度明确。
\end{itemize}
\begin{figure}[htbp]
\centering
\includegraphics[width=0.88\linewidth]{figures/fig_111.jpg}
\caption{Exhibit 1 - Exhibit 1: The deal-level arbitrage spreads generally correlate with the new issue calculations Deal-matched arbitrage spread is trailing 3-month deal-weighted figure; new Issue spread uses new issue loan spreads net monthly new is}
\end{figure}
\begin{figure}[htbp]
\centering
\includegraphics[width=0.88\linewidth]{figures/fig_112.jpg}
\caption{Exhibit 2 - Exhibit 2: The trend towards tighter arbitrage spreads coincides with more managers in a more mature market}
\end{figure}
\begin{figure}[htbp]
\centering
\includegraphics[width=0.88\linewidth]{figures/fig_113.jpg}
\caption{Exhibit 3 - Exhibit 3: USD Middle Market CLO arbitrage spreads have tightened towards 300bp Direct Lending spread calculated as difference between CDLI yield-to-maturity and 3-month LIBOR (SOFR after December 2020) USD Middle Market CLO arbi}
\end{figure}
\begin{figure}[htbp]
\centering
\includegraphics[width=0.88\linewidth]{figures/fig_114.jpg}
\caption{Exhibit 4 - Exhibit 4: EUR IG tranches offer a spread pickup relative to USD EUR/USD tranche spread ratios (adjusted with 5-year cross-currency basis); median since 2021 EUR/USD primary tranche spread ratio}
\end{figure}
\begin{figure}[htbp]
\centering
\includegraphics[width=0.88\linewidth]{figures/fig_115.jpg}
\caption{Exhibit 5 - Exhibit 5: USD deal economics have improved by less than 20bp on reset since 2023 Change in USD CLO arbitrage spread for resets by new issue vintage}
\end{figure}
\begin{figure}[htbp]
\centering
\includegraphics[width=0.88\linewidth]{figures/fig_148.jpg}
\caption{Exhibit 2 - Exhibit 2: China Property Developers' bond yields vs. share prices The average bond yield widened 18bp WoW to \$9.58\textbackslash{}\%\$ , and the share price \$-0.9\textbackslash{}\%\$ WoW.}
\end{figure}
\begin{figure}[htbp]
\centering
\includegraphics[width=0.88\linewidth]{figures/fig_149.jpg}
\caption{Exhibit 3 - Exhibit 3: Hong Kong property issuers' bond spreads vs. share prices The average bond spread (OAS) tightened 3bp WoW to 61bp and the share price -7.2\% WoW.}
\end{figure}
\begin{figure}[htbp]
\centering
\includegraphics[width=0.88\linewidth]{figures/fig_150.jpg}
\caption{Exhibit 4 - Exhibit 4: YTM (\%) of benchmark bonds for select Chinese HY property developers The average yield of benchmark bonds widened 10bp WoW on average basis.}
\end{figure}
{\small\color{refgray}\textbf{References:} [R015] GS：CLO套利空间的真正故事不在利差本身，而在资本结构如何重塑定价权; [R018] 美国银行：市场低估了房地产信用分化中的结构性机会，而非整体复苏}

\section{结语}
综合来看，当前市场主线正从总量叙事转向结构分化：通胀的‘分裂’、AI的‘扩散’、能源的‘需求破坏’、以及区域利润的‘重构’，共同指向一个需要更精细分析框架的定价环境。理解这些结构性变化，比追逐单一数据点或宏观叙事更为重要。
\newpage
\section{报告覆盖清单}
\begin{itemize}
  \item [R001] 宏观与利率：通胀结构分裂与政策信号再校准 - DB：市场正从“成本推动的通胀”转向“需求分化的定价游戏”
  \item [R002] 宏观与利率：通胀结构分裂与政策信号再校准 - 摩根斯坦利：市场低估了“再通胀退潮”的结构性含义
  \item [R003] 宏观与利率：通胀结构分裂与政策信号再校准 - Citi：PPI再通胀的真正含义不是能源，而是中国工业定价权的结构性拐点
  \item [R004] 中国政策与投资：‘六网’落地与资本管控信号再解读 - Citi：市场低估了“六网”投资对中国固定资产投资的托底力度
  \item [R005] 风险偏好：资金流拥挤与‘通胀高于增长’的再定价 - Citi：市场正在重新定价“增长的质量”，而非增长本身
  \item [R006] 未被主题摘要引用 - JPM：市场低估的不是原油价格，而是美国西海岸成品油的结构性短缺
  \item [R007] 生物医药与健康：中国创新资产需求刚性，阿尔茨海默病成本曲线待重构 - JPM：市场低估了全球药企对中国创新资产的刚性需求，而非地缘风险
  \item [R008] 电动车与汽车：区域利润结构重构与燃油车定价崩塌 - JPM：全球电动车市场的真正分歧不在渗透率，而在区域利润结构
  \item [R009] 区域市场：欧元弱势与人民币汇率信号切换 - 美国银行：欧元弱势远未结束，市场低估了技术面与基本面的共振
  \item [R010] AI/算力：从模型竞赛到成本信任，供应链紧俏重塑定价权 - JPM：中国通胀的真正信号不是PPI回升，而是AI定价权开始向消费端传导
  \item [R011] AI/算力：从模型竞赛到成本信任，供应链紧俏重塑定价权 - JPM：市场真正低估的不是通胀读数，而是AI成本传导正在重塑中国的价格结构
  \item [R012] AI/算力：从模型竞赛到成本信任，供应链紧俏重塑定价权 - Bernstein：AI模型商品化进程比市场想象的更快，竞争格局正从智能竞赛转向成本与信任竞赛
  \item [R013] 能源与大宗：需求破坏对冲供给冲击，铝与原油的结构性变局 - GS：铝市场真正超预期的不是中国产量，而是海外供给的加速释放
  \item [R014] 能源与大宗：需求破坏对冲供给冲击，铝与原油的结构性变局 - 摩根斯坦利：中国原油进口的骤降并非需求崩塌，而是库存周期的主动调节
  \item [R015] 固定收益与信用：CLO套利空间收窄与信用分化中的alpha机会 - GS：CLO套利空间的真正故事不在利差本身，而在资本结构如何重塑定价权
  \item [R016] 能源与大宗：需求破坏对冲供给冲击，铝与原油的结构性变局 - GS：市场低估的不是地缘冲突的烈度，而是需求破坏的持久性
  \item [R017] 未被主题摘要引用 - 摩根斯坦利：北美油气高管薪酬设计已传递出比产量更重要的信号
  \item [R018] 地产与消费：分化中的结构性机会与成本端出清 / 固定收益与信用：CLO套利空间收窄与信用分化中的alpha机会 - 美国银行：市场低估了房地产信用分化中的结构性机会，而非整体复苏
  \item [R019] 未被主题摘要引用 - BARC：校园网络市场正在回归正常增长，真正的竞争焦点是份额再分配
  \item [R020] AI/算力：从模型竞赛到成本信任，供应链紧俏重塑定价权 - 美国银行：市场低估的不是AI算力需求，而是CPU在Agentic时代的重新定价
  \item [R021] 工业与材料：供给约束与新兴需求驱动的再定价 - Citi：光伏市场的真正拐点不在国内价格战，而在海外框架协议的规模与结构
  \item [R022] AI/算力：从模型竞赛到成本信任，供应链紧俏重塑定价权 - Citi：AI供应链的紧俏不是短期噪音，而是新一轮定价权转移的信号
  \item [R023] 电动车与汽车：区域利润结构重构与燃油车定价崩塌 - JPM：中国汽车市场真正的拐点不是电动化渗透率，而是燃油车库存的定价崩塌
  \item [R024] 能源与大宗：需求破坏对冲供给冲击，铝与原油的结构性变局 / 工业与材料：供给约束与新兴需求驱动的再定价 - 摩根斯坦利：市场低估了供给侧的“结构性断裂”，而非仅仅需求回暖
  \item [R025] 工业与材料：供给约束与新兴需求驱动的再定价 - 摩根斯坦利：中国聚乙烯净进口量才是美国合同价格真正的定价锚
  \item [R026] 宏观与利率：通胀结构分裂与政策信号再校准 - BARC：通胀数据揭示的真正风险不是物价上涨，而是定价权断裂
  \item [R027] 宏观与利率：通胀结构分裂与政策信号再校准 / 区域市场：欧元弱势与人民币汇率信号切换 - NOM：人民币中间价模型正在揭示一个被低估的政策信号
  \item [R028] 宏观与利率：通胀结构分裂与政策信号再校准 / 区域市场：欧元弱势与人民币汇率信号切换 - NOM：人民币中间价模型正在释放一个被市场低估的政策信号
  \item [R029] 中国政策与投资：‘六网’落地与资本管控信号再解读 - GS：市场真正低估的不是2万亿数据中心投资，而是“六张网”的宏观含义
  \item [R030] 中国政策与投资：‘六网’落地与资本管控信号再解读 - 摩根斯坦利：市场误读了中国的资本管控信号，真正焦点在实体经济出海而非个人资金外逃
  \item [R031] 新兴市场专题：印度电力时段性短缺与日本半导体设备份额反转 - Bernstein：市场低估的不是电力需求，而是非日照时段的供给刚性
  \item [R032] 新兴市场专题：印度电力时段性短缺与日本半导体设备份额反转 - Bernstein：日本半导体设备在中国市场的份额流失并非结构性，反转正在酝酿
  \item [R033] 电动车与汽车：区域利润结构重构与燃油车定价崩塌 - Citi：中国新能源车市真正的拐点不是渗透率破60\%，而是龙头集中度开始松动
  \item [R034] 地产与消费：分化中的结构性机会与成本端出清 - GS：鲜零食零售的兴起不是威胁，而是价值零售龙头品类升级的催化剂
  \item [R035] 地产与消费：分化中的结构性机会与成本端出清 - GS：IP零售的“高基数疲劳”正在重塑市场对增长的定价逻辑
  \item [R036] AI/算力：从模型竞赛到成本信任，供应链紧俏重塑定价权 - JPM：AI驱动的半导体超级周期才刚刚进入第二阶段
  \item [R037] AI/算力：从模型竞赛到成本信任，供应链紧俏重塑定价权 - 摩根斯坦利：AI需求的结构性韧性正在重塑台湾科技股的定价逻辑
  \item [R038] 地产与消费：分化中的结构性机会与成本端出清 - 摩根斯坦利：大宗商品周期的分化正在重塑中国消费股的盈利格局
  \item [R039] 未被主题摘要引用 - Bernstein：人形机器人真正的竞争壁垒不是技术，而是分化能力
  \item [R040] 未被主题摘要引用 - JPM：AI服务器与PC的背离正在重塑台湾ODM的投资逻辑
  \item [R041] 未被主题摘要引用 - 摩根斯坦利：中国机场股的核心问题不是客流，而是盈利结构尚未完成重置
  \item [R042] 未被主题摘要引用 - 摩根斯坦利：市场低估了eBay和Etsy在智能体电商时代的独特优势
  \item [R043] 工业与材料：供给约束与新兴需求驱动的再定价 - 摩根斯坦利：印度钢铁的强势周期远未结束，但真正值得关注的是供给侧的再定价
  \item [R044] AI/算力：从模型竞赛到成本信任，供应链紧俏重塑定价权 - NOM：光器件供应链的瓶颈不在产能，而在核心衬底的锁定能力
  \item [R045] 宏观与利率：通胀结构分裂与政策信号再校准 / 风险偏好：资金流拥挤与‘通胀高于增长’的再定价 - 美国银行：市场正在为一场“通胀高于增长”的资产重新定价做准备
  \item [R046] 生物医药与健康：中国创新资产需求刚性，阿尔茨海默病成本曲线待重构 - 摩根斯坦利：阿尔茨海默病的真正拐点不在新药，而在“脑穿梭”能否重构成本曲线
\end{itemize}
\section{逐篇报告摘录}
\subsection{[R001] DB：市场正从“成本推动的通胀”转向“需求分化的定价游戏”}
{\small\color{refgray}归类：宏观与利率：通胀结构分裂与政策信号再校准}

这份来自DB的最新通胀监测报告，表面上是5月CPI与PPI数据的例行更新，但仔细拆解其内部结构变化后，一个更关键的判断浮现出来：中国经济正在经历一轮价格信号的“结构性分裂”。上游PPI的持续回暖并非简单的成本传导，其驱动力正在从能源价格切换至AI相关需求与中下游的议价能力重塑。与此同时，CPI的停滞与核心通胀的走软，则揭示出终端消费需求并未同步复苏。报告给出的核心洞察是，市场真正需要关注的，不是通胀数字本身的高低，而是这一轮价格分化如何重新定义不同产业的利润分配与竞争格局。

这份报告发布于2026年6月10日，覆盖了5月份的通胀数据。其价值不在于数据点的罗列，而在于指出了驱动力的切换。对于产业决策者和投资者而言，理解这种切换，比预测下个月的CPI数字重要得多。

1. PPI的再通胀已从“能源单引擎”切换至“AI与传导双引擎”

5月PPI同比上涨3.9\%，符合预期，环比仍为正增长0.5\%。但报告明确指出，4月之前驱动PPI上涨的核心变量——国际能源价格——在5月趋于平稳，石油相关行业价格甚至出现小幅环比下跌。这意味着，PPI的韧性并非来自能源的二次冲击，而是来自两个新的结构性力量。

第一个力量是前期能源价格上涨的“成本传导”效应正在向中下游扩散。化学制品、黑色金属、纺织等中下游行业在5月录得了更快的环比价格涨幅。这并非新一轮成本冲击，而是产业链消化前期成本后的重新定价。

第二个力量更具长期意义：AI相关需求持续推高锡、铜等有色金属价格，并带动电气机械、光纤电缆、计算机等终端产品价格上行。这不是短期的库存回补，而是由资本开支驱动的结构性需求。报告将其单独列出，暗示这一力量在未来几个季度可能成为PPI的新锚点。

这意味着，PPI的后续走势将不再主要取决于布伦特原油的波动，而是取决于两个变量：成本传导的广度与AI相关需求的强度。DB维持PPI年末升至5\%左右的预测，但驱动力的切换意味着，不同行业面临的成本环境和定价能力将出现显著分化。

2. CPI的停滞揭示的不是“通胀可控”，而是“需求收缩”的底色

与PPI形成鲜明对比的是，5月CPI同比维持在1.2\%，但环比转为负增长0.1\%。报告使用的措辞是“softening domestic demand”。服务、食品、家庭用品、汽车等类别均出现环比下跌，同比涨幅也在收窄。

这组数据指向一个不容忽视的现实：居民消费的复苏动能正在减弱。以旧换新补贴对商品消费的拉动效应正在边际递减。报告因此下调了全年CPI预测至1.5\%，年末目标从2.0\%下调至1.8\%。

这里的关键洞察是，CPI的低迷并非因为供给过剩，而是因为需求侧缺乏自发性扩张的动力。补贴驱动的消费脉冲过后，内生需求并未跟上。这与PPI的结构性回暖形成了“温差”。对于企业而言，这意味着上游成本压力能否向下游传导，将不再是产业链内部的谈判问题，而是取决于终端消费者是否愿意接受涨价。从目前的CPI数据看，答案并不乐观。

3. AI需求是唯一同时推高PPI与CPI的“异类”，但其规模仍待检验

报告中最值得关注的细节，是AI相关需求在两端价格指数中的体现。在PPI端，它推高了有色金属、电气机械和计算机设备的价格；在CPI端，它推高了手机和电脑的价格，5月环比分别上涨1.6\%和1.1\%。这在当前通缩压力弥漫的消费品市场中，几乎是唯一的价格亮点。

这引出一个更重要的判断：AI相关的资本开支和消费升级，正在成为一个独立于宏观经济周期的价格驱动力。它不受居民收入预期疲软的影响，也不受传统库存周期的制约。如果这一趋势持续，它将重塑相关产业链的定价权和利润分配。

但报告对此并未完全展开。AI需求对PPI和CPI的拉动幅度、可持续性、以及对传统制造业的挤出效应，都是尚未回答的问题。我们只知道它正在发生，但不知道它的量级

[... middle omitted ...]

力。这背后的逻辑是，在需求疲软的环境下，企业缺乏议价权。它们要么选择牺牲利润来维持市场份额，要么选择提价并承受销量下滑的风险。

报告下调CPI预测，本质上是对“传导失败”的一种确认。对于投资者而言，这意味着需要重新审视那些处于产业链中游、且下游需求缺乏弹性的行业。它们可能面临“增收不增利”甚至“量价齐跌”的困境。而那些能够将成本完全传导、甚至通过技术升级实现定价溢价的行业，将获得超额收益。AI相关产业链显然是后者的候选。

基于这份报告，一个更实用的分析框架是放弃对CPI和PPI具体数字的过度关注，转而绘制一张“利润分配地图”。这张地图的横轴是产业链位置（上游、中游、下游），纵轴是需求属性（AI驱动、补贴驱动、内生消费驱动）。

报告提供的数据已经勾勒出轮廓：上游受益于AI需求和供给约束，利润稳定；中游面临成本传导与需求疲软的双重挤压，利润承压；下游只有AI相关消费品（手机、电脑）和季节性商品（服装）表现较好，其余品类均面临需求收缩。补贴驱动的消费品类正在失去支撑。

\subsection{[R002] 摩根斯坦利：市场低估了“再通胀退潮”的结构性含义}
{\small\color{refgray}归类：宏观与利率：通胀结构分裂与政策信号再校准}

5月的通胀数据刚刚发布，PPI同比冲至3.9\%，创下近四年新高。如果只看这个数字，很容易得出“通胀压力正在积聚”的判断。但摩根斯坦利这份最新报告揭示了一个完全不同的图景：通胀的驱动力正在快速切换，而真正重要的不是PPI同比有多高，而是环比动能在哪里减速。

报告的核心主张可以概括为一句话：5月数据确认了中国正处于“再通胀退潮”的早期阶段，油价脉冲消退后，下游缺乏接力，核心CPI将进入温和下行通道。 这不是一个短期波动，而是一个结构性信号——意味着企业定价权的分布正在发生根本性变化。

为什么现在需要认真对待这份报告？因为市场对通胀的讨论仍然停留在“油价涨了多少”“PPI同比会不会破4\%”这些表层问题上。而摩根斯坦利把焦点拉到了更关键的层面：在油价推动消退之后，还有哪些力量能支撑价格？答案是，几乎没有。这种“单一引擎熄火”的格局，对资产定价、行业配置和宏观判断都有深远含义。

1. 5月PPI环比动能的骤降，暴露了“低基数幻觉”背后的真实疲弱

5月PPI同比从2.8\%跳升至3.9\%，表面上看非常强劲。但摩根斯坦利报告明确拆解了背后的结构：环比从4月的1.7\%骤降至0.5\%，降幅高达1.2个百分点。而环比才是衡量当前价格动能的真实指标。

这个环比减速几乎完全由油价驱动。4月油价对PPI环比的贡献高达1.7个百分点中的绝大部分，而5月全球油价趋于稳定后，石油相关的价格推升效应归零。换言之，4月的“通胀脉冲”本质上是油价的一次性冲击，而不是需求驱动的持续性上涨。

对于投资者来说，这意味着什么？如果你在4月PPI数据公布后认为“通胀上行趋势已经确立”，那么5月的数据就是一个重要的修正信号。PPI同比在未来两个月可能因为低基数继续走高至4\%以上，但环比动能才是判断趋势的关键——而环比已经开始走弱。

摩根斯坦利报告做了一个非常有价值的拆解：剔除石油和石化行业后，5月PPI环比仅微升至0.3\%，相比4月的0.2\%几乎没有加速。这意味着，油价熄火之后，其他行业没能形成有效的接力。

报告指出，少数几个行业确实有涨价迹象。煤炭价格环比上涨3.2\%，主要受夏季用电高峰预期的季节性拉动。有色金属冶炼环比上涨1.1\%，与铝相关的资本开支需求有关。但这些驱动因素要么是季节性的，要么是局部的，无法替代油价对整个工业品价格体系的广泛拉动。

更深层的问题是下游。尽管出口表现强劲，下游工厂出厂价格环比仅持平于0.2\%，没有出现任何加速迹象。摩根斯坦利将原因归结为“普遍的产能过剩”——这是报告中最值得关注的判断之一。出口的扩张没有转化为定价权，因为供给端的竞争仍然激烈，企业缺乏提价的能力。

这意味着，即使宏观需求有所改善，利润向中下游传导的路径仍然受阻。对于制造业投资者而言，这是一个需要持续跟踪的信号：产能出清的速度决定了下一轮定价权回归的时间窗口。

3. 核心CPI的“AI需求脉冲”正在消退，消费疲弱才是底色

如果说PPI反映的是工业端的价格压力，那么CPI则更直接地反映了终端消费的冷暖。5月CPI同比持平于1.2\%，看似稳定，但结构变化值得警惕。

报告特别指出，核心CPI剔除黄金后，经季节性调整的环比从4月的0.2\%降至0\%。这个“归零”背后有两个原因。第一，AI驱动的个人电子产品换机需求正在减弱——此前这一需求曾对核心CPI形成小幅提振，但脉冲正在消退。第二，也是更根本的，消费疲弱仍然在压制价格。就业市场疲软和房地产市场的持续调整，使得居民消费意愿和能力都受到抑制。

值得注意的是，报告提到“资本密集型出口的支撑有限”。这是一个容易被忽视的判断。出口强劲通常被认为会通过就业和收入渠道提振消费，但摩根斯坦利认为，资本密集型出口对普通劳动者收入的传导效应较弱，因此对消费的拉动作用有限。

对于宏观投资者而言，核心CPI的走势比P

[... middle omitted ...]

慢于预期。

但这个风险到底有多大？报告没有给出具体的概率评估，也没有讨论不同油价情景下对通胀的量化影响。这是一个关键的未解问题。如果油价再次大幅上涨，那么“再通胀退潮”的判断就需要修正。但如果油价维持在当前区间，那么核心CPI的温和下行将成为主导趋势。

对于读者来说，这意味着需要建立自己的情景分析框架：跟踪油价走势、地缘政治动态以及霍尔木兹海峡的通行数据，才能判断摩根斯坦利这份报告的核心判断是否仍然成立。

基于摩根斯坦利报告的分析框架，我们可以提炼出三个核心观察维度，用于跟踪后续数据的含义。

第一，关注PPI环比而非同比。同比数据受基数效应干扰太大，未来两个月PPI同比可能继续走高，但这不意味着通胀压力在上升。环比是否能够企稳甚至回升，才是判断工业价格动能的真实指标。

第二，关注“排除油价后的PPI环比”。这是报告最值得借鉴的分析方法。只有剔除油价的一次性冲击，才能看清工业价格的底层趋势。如果排除油价后的环比持续低于0.3\%，说明产能过剩的压力仍然主导市场。

\subsection{[R003] Citi：PPI再通胀的真正含义不是能源，而是中国工业定价权的结构性拐点}
{\small\color{refgray}归类：宏观与利率：通胀结构分裂与政策信号再校准}

Citi：PPI再通胀的真正含义不是能源，而是中国工业定价权的结构性拐点

五月通胀数据公布后，市场注意力几乎被CPI的小幅不及预期和能源价格回落所吸引。但这份Citi研报提出的核心判断值得认真对待：PPI再通胀正在从“能源脉冲”切换为“结构扩散”——这不是一次短暂的价格波动，而是中国工业部门定价逻辑正在发生改变的早期信号。

Citi维持CPI和PPI全年预测分别为1.0\%和2.8\%同比，这本身并不激进。真正有意思的是报告在数据细节中捕捉到的三个结构性变化：投资企稳的初步迹象、AI超周期开始显性化为价格力量、以及K型分化下下游通缩与上游通胀的持续拉锯。

这些事实合在一起，指向一个比市场共识更复杂的叙事：中国的价格信号正在从“总量问题”演变为“结构问题”。对于产业决策者和资产配置者来说，理解这种结构性分化，远比争论CPI是否见底更有价值。

五月PPI同比3.9\%，略高于Citi预期的3.8\%，但环比增速从四月的1.7\%大幅回落至0.5\%。表面上看，这是中东冲突溢出的价格冲击正在消退——石油和天然气开采PPI环比从18.5\%骤降至-0.1\%，燃料加工PPI环比归零。

但报告明确指出，真正值得关注的不是能源的退潮，而是再通胀的“广度”在扩大。Citi估算，能源和化工对五月PPI环比的合计贡献仅为0.2个百分点，远低于四月的1.6个百分点。这意味着，即使剔除能源冲击，PPI仍然录得了0.3个百分点的环比正增长。

这个“剔除能源后的PPI”才是报告的核心洞察。它说明工业品价格上涨的动力正在从单一的地缘政治冲击，切换为更广泛的内生性因素——包括投资需求的边际改善、上游原材料成本的传导、以及部分行业供给侧的收紧。

对于投资者而言，这意味着不能再用“能源价格见顶”来简单判断PPI拐点。五月的数据表明，即使能源贡献消退，其他行业的价格动能仍在累积。PPI的韧性可能比市场预期的更持久。

报告中最值得关注的信号来自建筑业相关价格。黑色金属矿采选PPI同比3.3\%、环比0.9\%；非金属矿采选PPI同比-3.4\%，但环比已经转正为0.4\%。Citi估算，建筑业相关项目对PPI环比的合计贡献升至0.1个百分点——虽然绝对值仍然很小，但这已经是过去一年中的第二高读数。

报告将此描述为“投资企稳的试探性信号”，并与建筑业PMI的改善相互印证。这个判断的谨慎之处在于“试探性”三个字：价格改善的幅度仍然有限，且主要集中在采矿环节，尚未向下游的冶炼和建材制造充分传导。

但从另一个角度看，这正是政策最有可能发力的领域。如果建筑业价格的改善能够持续，Citi认为“政策行动可能已经开始”。报告将劳动力市场状况视为进一步定向宽松的触发因素——这个判断隐含的逻辑是：投资企稳本身不足以解决需求问题，但如果投资改善能带动就业，政策的空间就会打开。

对于关注基建和地产周期的读者，这里有一个重要的观察框架：五月的数据可能只是起点，未来两到三个月的建筑业价格走势，将是判断“投资是否真正企稳”的关键验证窗口。

这是报告中一个容易被忽略但极具前瞻性的判断：AI超周期正在成为通胀的新驱动力，不仅影响科技PPI，也开始传导至电信CPI。

具体数据支持了这一论点。五月通信设备CPI环比上涨1.5\%，同比达到6.6\%，这是历史第二高读数。与此同时，计算机和其他电子设备PPI同比2.1\%、环比0.6\%，虽然涨幅不如上游能源，但趋势明确向上。报告还特别提到，光纤制造和电缆制造的价格动能也在增强。

与之形成鲜明对比的是汽车行业。汽车CPI环比下降0.4\%，已是连续第三个月下滑，同比跌幅维持在-2.0\%。Citi将这一分化归因于“AI超周期的展开”——需求结构正在从传统耐用消费品向科技硬件转移。

这个判断意味着什么？中国的通胀结构正在被技术周期重塑。过去市场习

[... middle omitted ...]

著回落——车辆租赁和机票价格分别环比下降6.8\%和6.3\%。

下游制造业的通缩更加顽固。报告整理的数据显示，在27个下游PPI行业中，有超过一半仍处于同比负增长状态。汽车制造业PPI已连续44个月处于通缩区间，非金属矿物制品业连续44个月，通用设备制造业连续40个月。这些数字本身就在讲述一个故事：中国制造业的产能过剩和需求不足，并非短期现象。

Citi将这种分化概括为“经济的K型分化”——上游受益于能源、资源和AI周期，下游则持续承受需求疲软的压力。这种分化在五月不仅没有收敛，反而有固化的趋势。上游价格向中下游的传导仍然受阻，核心CPI的走弱就是最直接的证据。

对于产业决策者而言，这意味着“成本传导”的假设需要被重新审视。不是所有行业都能将上游涨价转嫁给下游，那些缺乏定价权、产能过剩的行业，可能面临利润被持续挤压的风险。

Citi的报告在识别结构性变化方面表现出色，但它也留下了一个未解的核心问题：如果投资企稳和AI周期是供给侧的力量，那么需求侧的拐点何时到来？

\subsection{[R004] Citi：市场低估了“六网”投资对中国固定资产投资的托底力度}
{\small\color{refgray}归类：中国政策与投资：‘六网’落地与资本管控信号再解读}

Citi：市场低估了“六网”投资对中国固定资产投资的托底力度

中国四月固定资产投资增速转负，市场情绪一度转向悲观。但Citi在最新研报中提出了一个值得重新审视的判断：政策正在加速落地，以“六网”为核心的新一轮基建投资，其规模、节奏和融资路径，都比市场预期的更为清晰和坚决。这份报告的价值不在于罗列项目清单，而在于它揭示了投资结构正在发生的深层切换——从传统的房地产和旧基建，转向由AI算力、数据中心和城市更新驱动的新经济投资。这意味着，市场对投资增速的线性外推可能低估了结构性变量的影响。

四月固定资产投资累计同比增速转负，是这一轮经济复苏节奏中最刺眼的数据点。按照传统框架，这意味着内需疲弱，政策刺激效果递减。但Citi的分析提供了一个不同的视角：投资总量的短期失速，恰恰是结构转换期的典型特征。旧动能（房地产）的收缩速度超过了新动能（新基建、制造业升级）的扩张速度，造成了统计上的“真空期”。而近期一系列高规格项目的启动，正是政策层试图填补这一真空期的主动作为。

三峡新通道、两万亿级别的数据中心投资计划、以及城市更新的明确量化目标，这三件事在同一个月内密集发生，不是巧合。它们共同指向了一个顶层设计的落地：即“六网”建设。Citi的报告点出了一个关键事实——国家发改委年初承诺今年“六网”投资规模超过7万亿元，而市场并未给予足够定价。

1. 三峡新通道与数据中心：两个标志性项目打开了投资的想象空间

六月八日，三峡新通道项目正式开工，总投资约772亿元，建设周期约十年。这是“十五五”规划中首个启动的超级工程。它的意义不仅在于投资体量，更在于它释放了一个明确的信号：中央层面对大型基建项目的审批和推进节奏正在加快。

与此同时，彭博报道中国计划在未来五年投入约两万亿元建设数据中心。Citi在报告中进一步指出，如果将配套的电网整合纳入考虑，总投资规模可能至少达到五万亿元。这是一个值得反复咀嚼的数字。五万亿元意味着什么？它相当于2025年全年固定资产投资总额的约8\%左右，且集中在五年内释放。对于IT基础设施投资、电力设备、冷却系统、以及相关建筑服务行业，这构成了一个中期维度上极为确定的增长主线。

这两个项目，一个代表传统基建的延续与升级（水利、交通），一个代表新经济基础设施的扩张（算力、能源）。它们并行推进，说明“六网”不是对旧模式的简单重复，而是一次投资重心的系统性位移。

市场对基建投资的一个常见疑虑是资金从何而来，尤其是地方政府债务压力尚未缓解的背景下。Citi的报告对此给出了较为清晰的回答：融资渠道是多元且充足的。

报告指出，资金来源包括一般公共预算、特别国债、政策性金融工具和政府引导基金。具体来看，中央一般预算已为2026年城市更新安排970亿元；特别国债中明确为地下管网改造安排了1600亿元，数据中心建设也可能获得类似支持。此外，政策性金融工具在2025年和2026年分别提供了5000亿和8000亿元的初始额度。

这些数字加在一起，构成了一个规模可观的财政支持矩阵。更重要的是，这些资金是“定向”的，即专门用于“两重”项目（国家重大战略和重大工程）和“六网”建设。这降低了资金被挪用或低效使用的风险。对于产业投资者而言，这意味着“六网”相关领域的项目落地确定性，要高于普通的地方政府基建项目。

Citi的报告没有停留在描述政策本身，而是追问了一个“为什么是现在”的问题。答案有两个层面，一个是周期性的，一个是结构性的。

周期性层面，四月固定资产投资累计同比增速转负，是直接的政策触发信号。尽管社会消费品零售总额等数据表现尚可，但投资端的失速引起了决策层的高度关注。Citi引用去年底中央经济工作会议的表述——“要稳住投资增速”，说明政策目标是一贯的，而四月的数据让这一目标的紧迫性显著上升。

结构性层面，AI超级

[... middle omitted ...]

需求仍然偏弱的背景下。这里的关键判断是“K型复苏”——新经济部门（尤其是“六网”相关领域）将引领投资回升，而传统部门可能继续承压。这意味着，总量数据的改善将是渐进的、结构分化的，而非V型反转。

对于投资者而言，关注总量固定资产投资增速的拐点，不如关注IT固定资产投资、数据中心相关招标、以及电网投资等细分领域的趋势。Citi报告中的图表显示，IT固定资产投资增速在2025年和2026年前四个月持续保持在18\%以上的高位，远高于固定资产投资整体水平。这一分化本身就是K型复苏最直接的证据。

Citi的这份报告在宏观判断上具有前瞻性，但作为一份侧重于政策解读的研报，它对几个微观层面的问题并未充分展开，这恰恰是后续需要跟踪的重点。

第一个问题是，两万亿至五万亿的数据中心投资，其真实的回报率如何？AI算力需求是否足以支撑如此大规模的前期资本开支？如果需求增速低于预期，这部分投资可能面临产能过剩的风险。报告没有提供对AI应用端需求的具体测算，这是一个需要进一步验证的关键假设。

\subsection{[R005] Citi：市场正在重新定价“增长的质量”，而非增长本身}
{\small\color{refgray}归类：风险偏好：资金流拥挤与‘通胀高于增长’的再定价}

当绝大多数投资者还在争论AI是泡沫还是革命时，一份来自Citi的年度主题评估框架给出了一个更微妙的信号：市场已经不再为“故事”买单，而是在为“可持续性”和“质量溢价”重新排序。

Citi的Global Theme Machine覆盖了83个全球主题、近3600只股票，建立了超过2万个个股-主题关联。这份报告的价值不在于告诉你哪些主题涨了，而在于揭示了一个关键的结构性变化——主题排名的驱动力正在从单纯的“价格动量”转向“盈利可见性”与“质量因子”的组合。这意味着，即便是在同一个AI赛道内部，不同公司的命运分化才刚刚开始。

这不是一个关于“买什么”的报告，而是一个关于“用什么框架判断买什么”的报告。对于产业决策者和高净值投资者而言，理解这个框架的演化，可能比追逐任何一个具体主题都更重要。

1. 主题排名的驱动力正在发生根本性切换：从“动量”到“质量”

Citi的排名体系并非简单的涨跌幅排序。它整合了估值、增长、价格动量、质量、低风险、盈利动量六个维度，最终给出综合排名。最新数据显示，排名靠前的主题不再单纯依赖价格动量，而是呈现出“质量+盈利可见性”的双轮驱动特征。

以排名第二的“Risky Business”主题为例，其价格动量排名仅为53（共84个主题），但质量排名第2，低风险排名第1。同样，排名第三的“Pension Shortfalls”主题，估值排名第1，质量排名第7，低风险排名第2。这两个主题的共同特点是：它们不是市场上最“性感”的叙事，但它们在盈利稳定性、估值安全垫和风险控制上得分极高。

这意味着什么？Citi的框架在告诉我们：当前市场的定价逻辑正在从“谁的故事更宏大”转向“谁的盈利更可预测”。对于投资者而言，这要求重新审视组合中每个持仓的“质量溢价”——那些盈利波动大、但靠价格动量维持排名的主题，正在面临被降级的风险。

2. AI与工业主题的领先并非“全面开花”，而是“结构性分化”

从月度表现看，排名靠前的主题高度集中在数字化和工业转型领域：Greening the Home（+8.8\%）、Cloud Computing（+6.5\%）、Wearable Technology（+5.2\%）、AI Enablers（+4.6\%）、Manufacturing Onshoring（+3.7\%）。这些主题相对MSCI World的超额收益非常显著。

但Citi的框架揭示了一个更深层的分化：即便在AI主题内部，不同细分领域的排名驱动因子也完全不同。AI Enablers整体排名第6，其增长因子排名第4，但质量排名仅为25，低风险排名44。这说明AI Enablers仍然高度依赖增长预期和价格动量，而非基本面的稳健性。相比之下，Cloud Computing的排名从7提升至4，其改善主要来自质量因子从13提升至8，以及价格动量从44提升至35。

这里的insight是：市场正在对AI产业链进行“质量筛选”。那些能够将技术优势转化为稳定现金流、并具备定价权的公司（如云基础设施提供商），正在获得更高的排名溢价；而那些仍处于投入期、盈利路径不清晰的AI应用公司，其排名正面临来自质量因子的压制。

排名下滑最剧烈的主题包括：Luxury Spend（从36降至76）、Global Tourism（从35降至81）、Solar Energy（从62降至79）。这些主题的共同特征是：价格动量仍然存在（Luxury Spend的价格动量排名59，Solar Energy为11），但盈利动量已经显著恶化（Luxury Spend的盈利动量从17降至76，Global Tourism从33降至82）。

这是Citi框架中最危险的信号之一。当价格动量与盈利动量出现“剪刀差”时，意味着市场定价已经脱离了基本面支撑。

[... middle omitted ...]

维度：个股对多个主题的暴露度。报告列出了同时暴露于多个Top 15和Bottom 15主题的股票。例如，泰国的Gulf Development同时覆盖了10个Top 15主题（包括AI Enablers、Cloud Computing、Wearable Technology等）和5个Bottom 15主题（包括Solar Energy、Nuclear Energy、Biodiversity等）。

这种“多主题交叉暴露”是一个尚未被充分讨论的风险点。一方面，它意味着这些公司具备更强的主题多样性和业务韧性；另一方面，当市场风格切换时，它们可能同时承受来自多个方向的压力。Citi的框架目前没有给出“如何加权计算多主题暴露度”的方法论，这可能是报告最值得深挖的盲点。

对于投资者而言，这意味着：单纯看一个公司属于哪个“热门主题”是不够的，必须拆解其业务组合中哪些主题在驱动估值、哪些在拖累排名。一个同时暴露于AI和清洁能源的公司，其风险收益特征可能完全不同于一个纯粹AI公司。

\subsection{[R006] JPM：市场低估的不是原油价格，而是美国西海岸成品油的结构性短缺}
{\small\color{refgray}归类：未被主题摘要引用}

JPM：市场低估的不是原油价格，而是美国西海岸成品油的结构性短缺

这份JPM最新发布的石油市场周报，标题意味深长——“The first thousand miles”（最初的一千英里）。它指向一个反直觉的事实：自冲突爆发以来，历史上最大的供应中断，却只带来了一个“平淡无奇”的价格图表。布伦特原油均价仅101美元，几乎完全贴合该行早在3月初就划定的价格路径。

市场普遍关注的是，为什么原油价格没有暴涨？但这份报告真正值得读的判断是：市场已经通过需求损失和库存释放消化了供应冲击，但成品油市场的结构性紧张，尤其是美国西海岸，正在成为新的、更隐蔽的定价锚点。 对于关注通胀、供应链和资产定价的决策者而言，真正需要跟踪的，不是布伦特能否站上120美元，而是加州每加仑5.89美元的汽油和7.16美元的柴油，如何通过物流成本传导至全美。

1. 报告自己推翻了3月的核心假设：成品油短缺并未发生，但“向下游转移”的逻辑失效了

3月初，JPM提出了一个精妙的框架：如果原油价格无法充分反映供应中断，那么调整一定会“向下游转移”——即原油价格维持100美元附近，但短缺会体现在汽油、柴油和航空煤油的裂解价差上。这个框架的第一部分（原油价格路径）被完美验证，但第二部分（成品油短缺）被证伪。

从2月27日至6月8日的数据来看，汽油、柴油和航空煤油价格在经历最初飙升后，已显著回落，炼油利润率不但没有扩大，反而收窄。这意味着，市场吸收冲击的机制比预想的更复杂：需求损失（冲突区、亚洲、中国、非洲）和库存释放（政府和商业库存）共同作用，使得成品油市场并未出现极端短缺。 这个“证伪”本身就是一个重要洞察：它告诉我们，全球石油市场的缓冲能力，至少在成品油层面，比许多模型假设的要强。

但这份报告紧接着揭示了一个更深层的结构性矛盾——全国平均数据掩盖了区域性的极端分化。

2. 美国西海岸的成品油市场，正在经历“独立于全球”的紧缩周期

当全国平均汽油价格约为4.16美元/加仑时，加州的价格是5.89美元，柴油价格更是达到7.16美元。这不是简单的运输成本差异，而是多重结构性因素叠加的结果。

首先，加州拥有全美最严格的环保法规（CARB标准），这意味着它无法轻易从相邻州获取成品油供应，必须依赖亚洲炼厂进口。但自3月1日以来，霍尔木兹海峡的供应中断导致许多亚洲炼厂被迫减少成品油出口。加州的汽油进口中，近三分之二来自亚洲，占全州汽油供应总量的20\%至30\%。这条供应链的断裂，直接推高了区域价格。

其次，加州本土炼油产能正在萎缩。Valero的Benicia炼厂和Phillips 66的洛杉矶炼厂相继关闭，合计占加州炼油产能的17\%。在需求并未大幅下降的情况下，供给侧的收缩是刚性的。

*更深层的含义在于：美国西海岸不仅是加州的问题。** 全美约五分之一的进口货物通过洛杉矶港和长滩港进入。这些货物在运往全美各地仓库、配送中心和消费者之前，最初的数百甚至数千英里，都依赖西海岸的柴油和汽油。这意味着，美国供应链的“第一英里”成本，是由西海岸燃料价格决定的。加州7.16美元的柴油价格，会通过货运成本传导至全美商品价格。

3. 汽油和柴油面临截然不同的供需格局：一个靠进口缓解，一个因出口加剧

汽油方面，5月的价格上涨已触发部分供给响应：部分炼厂从航空煤油转产汽油，加拿Daiwa西北欧的进口开始抵达东海岸。但这种再平衡可能只是暂时的。欧洲汽油市场本身紧张，随着夏季需求高峰到来，欧洲炼厂保留本地桶油的意愿会增强。如果美国汽油进口放缓，国内库存将维持在较低的季节性水平。

柴油的处境则更为严峻。美国柴油出口在1-2月平均为120万桶/日，比去年同期高出15万桶/日。冲突爆发后，出口进一步升至150万桶/日。背后的驱动因素包括：巴西对俄罗斯柴油的替代需求、欧洲

[... middle omitted ...]

决的核心问题：当前的市场平衡是可持续的吗？还是说，它依赖于一次性的库存释放和尚未被充分量化的需求损失？

报告指出，政府库存释放和商业库存消耗是吸收冲击的关键。但库存是有限的。如果霍尔木兹海峡的局势持续，库存释放的节奏能否维持？需求损失（尤其是亚洲和非洲）是永久性的结构变化，还是暂时性的价格响应？如果是后者，当价格回落时，需求可能迅速反弹，重新暴露供给缺口。

此外，报告对“中国消费者比预期更快地替代石油”这一判断，并没有展开具体分析。这种替代是电动汽车渗透率提升的结果，还是短期行为？如果是前者，它对全球石油需求的影响将是长期的。

5. 对决策者的观察框架：从跟踪原油价格转向跟踪成品油裂解价差和区域价差

基于这份报告的逻辑，我们认为投资者和企业决策者应该调整自己的观察框架。

第一，放弃对布伦特原油绝对价格的过度关注。JPM的预测显示，2026年余下时间布伦特原油均价将在100美元附近波动。但真正影响通胀预期和企业成本的，是成品油价格，尤其是柴油和汽油的裂解价差。

\subsection{[R007] JPM：市场低估了全球药企对中国创新资产的刚性需求，而非地缘风险}
{\small\color{refgray}归类：生物医药与健康：中国创新资产需求刚性，阿尔茨海默病成本曲线待重构}

JPM：市场低估了全球药企对中国创新资产的刚性需求，而非地缘风险

投资决策者过去半年最纠结的问题，莫过于中国生物医药资产是否已被地缘政治风险“污染”。美国国会信函、FDA对中国临床数据的审查信号、药明康德被列入1260H清单——每一则新闻都足以让交易员本能地调低中国创新药的估值权重。

但JPM在6月11日发布的一份研报中，给出了一个与市场情绪相反的核心判断：全球大型药企对中国创新资产的结构性需求，其影响力远超美国政策噪音。这份报告并非来自JPM的中国团队，而是由其美国制药/生物技术分析师团队在ASCO和ADA会议后，面向全球投资者进行的专题电话会纪要。美国分析师们从全球创新流动、并购纪律、临床数据价值三个维度，系统性地论证了一个结论：真正决定中国生物医药资产价格的，不是华盛顿的政治信号，而是波士顿、新泽西和旧金山的研发管线缺口。

这不是一份“中国资产看多”的宣传稿。它的价值在于，给出了美国产业界如何看待中国创新的第一手视角——这些分析师每天与大型药企管理层对话，他们的判断直接反映了资本配置的真实逻辑。

JPM美国分析师团队对政策风险的定性，值得每一个关注中国医药资产的投资者仔细咀嚼。他们明确区分了“影响投资者情绪”和“改变全球药物开发轨迹”两个层次，并认为政策因素属于前者，而非后者。

具体而言，美国分析师指出，美国国会信函要求FDA限制中国临床数据使用，以及药明康德被列入1260H清单，更多是“政治信号”而非“可执行的壁垒”。在电话会上，他们反复使用了一个词：noise。这不是轻描淡写，而是基于对产业逻辑的深刻理解。

关键在于，全球创新药研发的底层逻辑没有改变：科学在哪里，创新就流向哪里。如果中国科学家和临床团队能够产出差异化的临床数据，大型药企没有理由因为政治噪音而放弃获得这些资产的机会。分析师特别强调，大型药企管理层普遍认为，即便面临政治压力，“实用主义解决方案”仍然可以找到——比如共同开发、美国之外桥接试验、共享商业化权益等。这些结构安排并非无法实现，只是增加了交易复杂度，而非阻断交易本身。

对于药明康德等CXO企业，美国分析师的判断同样务实：大型药企将继续寻找“中国+1”的供应商策略，但这不是“去中国化”，而是供应链多元化。药明康德在效率、规模和成本上的优势，短期内没有替代者可以完全复制。供应链调整是一个渐进过程，而非一刀切的撤离。

这一节的so what非常清晰：中国生物医药企业的估值中，政策风险溢价已经被过度定价。如果这个风险溢价开始消退，资产价格的重估空间将是可观的。

JPM分析师在电话会上引用了一项调查数据：超过60\%的投资者预计今年的并购交易将多于去年。大型药企被描述为“弹药充足”，聚焦领域包括肿瘤、免疫、心血管代谢/肥胖和神经科学。

但这里有一个关键变量，也是报告中最具洞察力的判断之一：买家越来越不愿意为一个西方资产支付“大额支票”，如果六个月后有一个同样优秀的中国资产出现的话。这句话的含义极其深刻。它意味着，中国创新资产的存在本身，正在改变全球并购市场的定价结构。大型药企在评估一个西方生物技术公司的收购标的时，现在必须考虑中国同类资产的竞争。快速跟随者正在侵蚀先行者的护城河，从而压缩收购回报率。

这对中国生物医药公司而言，既是机会也是挑战。机会在于，高质量的中国资产仍然是可货币化的——全球大型药企有强烈的意愿通过out-licensing获得这些资产。挑战在于，美国买家在估值上会更加严格，交易结构会更倾向于风险共担和分阶段对价支付。JPM分析师称之为“更精细的BD策略”：优先选择风险分担的合作结构，有选择性地进行中国之外地区的对外授权，并在做出大规模交易承诺之前，对全球竞争格局（包括中国管线）进行更深入的尽职调查。

这意味着，中国生物医药公司不能简单地将“licen

[... middle omitted ...]

这一药物类别已经基本被去风险化。虽然康方生物的ivonescimab能否在全球III期试验中成功仍不确定，但大型药企普遍遵循“创新在哪里就去哪里”的原则。分析师特别指出，如果一款药物在肺癌等适应症中明显改善生存率，仅仅因为部分试验在中国进行或由中国公司发现，就拒绝美国患者使用，这在伦理和临床实践上都是不可持续的。

这一判断的支撑点在于：美国国会信函虽然要求FDA限制中国临床数据的使用，但产业界认为这更多是“噪音”。实用主义的解决方案——共同开发、桥接试验、共享商业化——足以确保临床差异化的PD-1/VEGF资产仍然能够进入美国市场。

这对中国PD-1/VEGF领域的领先企业是明确的利好信号。只要数据能够证明与现有疗法相比的差异化优势，全球out-licensing和共同开发的需求就会持续存在。JPM分析师没有说的是，这实际上为中国创新药定价设定了一个新的参照系：如果中国资产能够在美国市场获得批准，其价值将不再仅仅由国内医保谈判定价，而是可以参照全球创新药定价体系。

\subsection{[R008] JPM：全球电动车市场的真正分歧不在渗透率，而在区域利润结构}
{\small\color{refgray}归类：电动车与汽车：区域利润结构重构与燃油车定价崩塌}

JPM：全球电动车市场的真正分歧不在渗透率，而在区域利润结构

五月的全球电动车销量数据，读起来像一份分裂的诊断报告。欧洲五国EV+PHEV渗透率同比跳升8个百分点至32\%，中国乘用车市场整体萎缩22\%但新能源渗透率却逆势突破63\%，而美国市场仍在原地踏步，渗透率甚至同比下滑至7\%。这三条曲线放在一起，最容易被忽略的判断是：全球电动车投资的主题正在从“总量增长”切换到“区域利润结构”的重构。

JPM这份五月全球电动车与电池月度更新，提供了几个值得产业决策者重新审视的信号。中国新能源车渗透率突破60\%之后，整车销量总量下滑反而成为结构性利好——它意味着车企的盈利驱动因素正在从国内规模转向海外溢价和新能源车占比。欧洲的补贴政策和廉价车型正在真实地拉动渗透率，而美国市场的持续疲软则暗示着消费者选择不足和OEM撤退的负反馈循环。与此同时，储能系统作为整份报告中最明确的亮点，正在成为电池产业链中定价权最集中的环节。

以下是我们从这份报告中提炼出的四个关键洞察，以及一个报告尚未完全回答的关键问题。

1. 中国新能源车渗透率突破60\%后，总量下滑反而成为盈利结构优化的催化剂

五月中国新能源乘用车零售销量为95万辆，同比下滑7\%，但渗透率却达到63\%，同比提升10个百分点。这个数字放在全球背景下来看，意义不亚于欧洲的复苏。JPM的中国汽车团队已经将2026年国内乘用车零售预测从此前的同比下滑4\%大幅下调至下滑15\%，但报告明确指出：“新能源车占比和海外敞口正在成为比行业总量增长更重要的盈利驱动因素。”

这意味着什么？当渗透率超过60\%，新能源车已经不再是“增量市场”中的增长故事，而是“存量市场”中的结构性替代。对于比亚迪和吉利这样的头部OEM，国内销量的绝对数字增长不再是估值锚点，取而代之的是两个更关键的变量：一是新能源车在整体销量中的占比能否继续提升，二是海外市场能否提供更高的单车利润和更稳定的定价环境。

JPM在参加完欧洲汽车大会后，将比亚迪和吉利2026年的海外销量预测分别上调了约15\%和30\%，同时下调了国内预测。这个调整本身就是一个清晰的信号：中国车企的利润增量正在从国内转向海外，而海外市场的竞争逻辑与国内完全不同——它考验的是品牌建设、渠道铺设和本地化生产能力，而非单纯的性价比。

2. 欧洲的复苏不是偶然，而是补贴、廉价车型和OEM策略调整的三重叠加

五月欧洲五国EV+PHEV销量同比增长39\%，渗透率达到32\%。法国和意大利的渗透率同比提升最为显著，分别达到13和9个百分点。报告将此归因于国家层面的补贴计划和平价电动车型的渗透。本周欧盟还推出了15亿欧元的电池助推设施，用于支持欧洲经济区内的电动车电池生产。

但真正值得关注的不是渗透率本身，而是欧洲市场的结构性变化。JPM在报告中提到，中国OEM在欧洲新能源车市场的份额已经达到18\%，同比翻了三倍。驱动因素不是简单的低价策略，而是“零售执行护城河”和“标配化的强功能配置”。这意味着中国车企在欧洲的竞争已经从“价格战”进入了“价值战”阶段——它们正在用国内打磨成熟的供应链效率和产品定义能力，在欧洲市场建立新的竞争维度。

对于欧洲本土OEM而言，这个趋势既是威胁也是验证。威胁在于中国品牌的份额增长正在加速，验证在于欧洲市场的电动车需求确实存在且正在被释放。关键的分水岭在于：欧洲OEM能否在2027年之前推出有竞争力的平价电动车型，并重建消费者信心。

美国五月EV+PHEV销量同比下降21\%，渗透率仅7\%。报告给出的解释很直接：消费者选择减少，因为OEM取消了新电动车发布并收缩了营销投入。这是一个典型的负反馈循环——需求不足导致OEM减少供给，供给减少进一步压制需求。

这与欧洲形成鲜明对比。欧洲的补贴政策在创造需求，中国的规模效应在降低成本，而美

[... middle omitted ...]

JPM在报告中用了一个明确的表述：“储能仍然是最清晰的亮点。”全球储能出货量在2026年前四个月同比增长超过一倍，驱动因素是中国国内需求和面向非美地区的强劲出口。更关键的是，AI数据中心储能需求正在形成新的增长极。

报告提到，西门子与英伟达和Fluence Energy合作开发AI数据中心参考电源与控制架构，Fluence已经与两家超大规模云服务商签署了主供应协议，首批订单预计在第三季度落地。JPM认为，这些进展应该能够缓解市场对“AI数据中心是否需要储能”的疑虑，并推动 年的需求预期上升。

这里有一个容易被忽略的细节：AI数据中心储能电池与公用事业级储能电池是两种不同的产品类别。AI数据中心储能要求快速响应、高功率密度和频繁的微循环，这天然倾向于那些拥有强产品质量的头部厂商。这意味着，储能市场的增长并非均匀分布，而是会进一步加剧电池厂商之间的分化。CATL作为全球最大的储能电池制造商，以及LGES和SDI作为美国市场的关键受益者，是JPM在这个领域最看好的标的。

\subsection{[R009] 美国银行：欧元弱势远未结束，市场低估了技术面与基本面的共振}
{\small\color{refgray}归类：区域市场：欧元弱势与人民币汇率信号切换}

这份报告的核心判断并不复杂，但它的分量在于时机。美国银行在5月20日通过一个 的看跌价差组合建立了欧元空头头寸，并在6月8日欧元测试1.15心理关口后，明确表示将继续持有空头进入夏季。这不是一个战术性的短线交易，而是一个基于结构性逻辑的持仓决策。

在多数市场参与者仍在猜测欧元何时触底、中东停火协议能否带来反转的时候，美国银行提供的分析框架给出了一个更清晰的判断：欧元当前面临的不是单一风险，而是技术面破位、基本面背离和历史模式重复三重压力的叠加。这三者单独看都不足以构成决定性信号，但合在一起，指向的是一轮可能被低估的美元上行周期。

1. 技术面已经给出了明确的破位信号，反弹是卖出机会而非反转起点

美国银行的技术分析师指出了三个关键的图表信号。第一个是欧元在6月5日美国非农数据公布后跌破了上升趋势线支撑，这是自4月以来的最低水平。第二个信号更加值得关注——周线图上，欧元正在构建一个头肩顶形态的右肩，颈线支撑位在 区间。一旦这个区间被有效跌破，头肩顶形态确认，下行目标将指向1.11甚至200周移动均线附近的1.0980。

第三个信号来自美元指数。美国银行的分析师发现，美元指数在特朗普第二任期的走势与第一任期高度相似。 年的模式是：大选后美元突破多年区间上涨约6.4\%，随后形成头肩顶并在2017年下跌，2018年一季度触底，4月底突破后继续上涨约6.6\%至8月。 年的走势几乎在重复这一路径。如果这个类比继续成立，美元指数在当前底部确认突破后，可能在未来几个月向103-105区间上行。

对于读者而言，这意味着什么？欧元当前1.15附近的反弹，在技术分析框架下不是反转的开始，而是空头重新入场的机会。美国银行明确表示，他们倾向于在反弹中卖出欧元，尤其是如果欧元无法在 区域（50日和200日移动均线附近）实现日线级别突破。

2. 基本面背离正在加剧，但市场因地缘政治幻想而低估了美元支撑

技术面提供了交易的时机和点位，但基本面的逻辑才是持仓信心的来源。美国银行的核心论点是：美国和欧元区的经济增长差异已经非常显著，但利率市场对这一差异的定价仍然不足。

具体来看，美国劳动力市场持续表现出韧性甚至走强，6月5日的非农数据是最新例证。与此同时，欧洲面临的是滞胀压力——虽然天然气价格远低于2022年危机水平，但霍尔木兹海峡的持续关闭使得能源价格保持高位，这对欧洲的贸易条件和经济增长构成持续拖累。

但这里有一个关键矛盾：市场对中东和平协议的期待正在阻止基本面因素完全定价。美国银行的分析师明确指出，这种“地缘政治愿望投射”是当前欧元未能充分反映基本面弱势的主要原因。换句话说，市场在押注一个尚未实现的和平协议，而忽略了协议即使达成，能源库存问题也将持续存在。

对于投资者的含义是：如果和平协议未能迅速达成，或者达成后的效果不及预期，当前被压制的美元上行压力将集中释放。更关键的是，即使协议达成，欧洲面临的能源成本结构性上升和增长动力不足的问题不会立即消失。

利率市场的定价变化正在为美元提供额外的支撑。美国银行指出，战争初期市场迅速定价欧洲央行（以及其它G10央行）将加息以应对全球通胀冲击，同时预期美联储将保持结构性鸽派立场。但最近的数据开始改变这一预期。

市场目前仍然定价欧洲央行将有更多加息，这在短期内限制了欧元的下行空间。但美国银行的分析师提出了一个关键质疑：考虑到欧洲更疲弱的经济增长背景，欧洲央行的加息可能难以持续，甚至可能成为政策错误。如果欧洲央行被迫在经济增长疲软的情况下加息，这对于欧元来说并非利好信号——它可能意味着更糟糕的经济前景而非更强的货币。

与此同时，美联储面临的加息压力正在增加。即使新任主席沃什倾向于抵制加息，持续的通胀压力、劳动力市场改善和金融条件宽松（如股市韧性）可能迫使美联储采取行动。市场已经反映

[... middle omitted ...]

开，这些恰恰是投资者需要自行判断的风险点。

第一个问题是中东和平协议的时间窗口。报告承认“和平协议的希望阻止了更大主题的形成”，但没有给出协议可能的时间框架或概率判断。如果协议在未来几周内意外达成，能源价格可能快速回落，欧洲的滞胀压力将暂时缓解，欧元可能获得短期反弹动力。这种情景下，美国银行的技术面信号可能暂时失效，但基本面的结构性差异不会因此消失。

第二个问题是欧洲央行政策错误的可能性。报告指出欧洲央行的加息可能“适得其反”，但没有量化这种风险的概率或幅度。如果欧洲央行在经济增长放缓的背景下加息，可能引发欧洲资产价格更广泛的重新定价，这种二阶效应对欧元的影响可能比直接利率差异更为复杂。

第三个问题是美元指数 模式重复的可靠性。技术分析中的历史类比从来不是精确预测，而是概率框架。当前的地缘政治环境和全球经济结构已经与 年有显著不同，包括更高的债务水平、更复杂的供应链调整和更分散的全球贸易格局。这个类比的成立需要更多确认信号，而不仅仅是日线收盘价突破100.51。

\subsection{[R010] JPM：中国通胀的真正信号不是PPI回升，而是AI定价权开始向消费端传导}
{\small\color{refgray}归类：AI/算力：从模型竞赛到成本信任，供应链紧俏重塑定价权}

JPM：中国通胀的真正信号不是PPI回升，而是AI定价权开始向消费端传导

中国5月通胀数据公布后，市场目光大多停留在PPI同比3.9\%的温和超预期上。但这份JPM研报传递的核心判断，远比一个数字深刻：中国正在经历一场由产业升级和AI需求驱动的结构性定价重估，而这股力量正在从工业端向消费端扩散。市场如果只看到“通胀温和、需求疲软”的旧叙事，就会错过供给端正在发生的根本性变化。

这份报告的价值不在于确认通胀符合预期，而在于揭示了一个正在成型的传导链条——AI驱动的半导体和电子元器件涨价，已经从上游的封装测试、存储设备，蔓延到下游的通讯工具、平板电脑等消费电子终端价格。这不再是短暂的成本冲击，而是与全球科技上行周期和中国“AI+”战略深度耦合的持久性通胀尾风。

1. 食品通缩掩盖了非食品价格的韧性，核心通胀的“假疲软”需要重新审视

5月CPI同比1.2\%，略低于市场预期的1.3\%，表面看延续了温和态势。但拆解结构后会发现，拖累主要来自食品端——鲜菜和猪肉价格继续下行，食品CPI同比降至-1.7\%。这恰恰是数据中最具误导性的部分。

非食品CPI同比已升至1.9\%，环比仍录得0.2\%的正增长。其中交通运输和通信类价格环比上涨0.2\%，虽然较前两个月月均2.0\%的环比增速有所放缓，但通信工具价格环比跳升1.5\%（经季调后JPM估算为1.8\%），统计局明确提及手机和平板电脑分别上涨1.6\%和1.1\%。背后驱动因素正是AI相关存储芯片的紧张正在向终端传导。

更值得关注的是核心CPI——剔除食品和能源后，同比虽回落至1.1\%，但环比持平。这个“平”字背后，是服务业CPI同比放缓至0.8\%的隐忧。但JPM的分析暗示，核心CPI的疲软更多反映的是消费信心不足和收入预期偏弱，而非通缩螺旋的加深。服务价格中，教育文化娱乐在劳动节旅游支撑下环比仍有0.2\%的正增长。

这意味着，中国通胀的结构正在发生分化：一端是食品通缩和部分消费品价格承压，另一端是科技相关产品和服务的价格正在上行。市场如果仅凭CPI整体读数就得出“通缩风险仍在”的结论，会严重低估供给端结构性变化对定价体系的重塑。

2. PPI上行的真正驱动力不是能源，而是工业升级和AI需求的叠加效应

5月PPI同比3.9\%，环比0.8\%，延续了自3月以来的加速态势。但驱动因素正在发生重要转变。JPM明确指出，本月的环比涨幅主要来自工业升级和AI驱动的算力需求，而非单纯的能源成本推动。

有色金属冶炼和压延加工是最大亮点。锡冶炼价格环比上涨4.8\%，铜冶炼上涨3.1\%，这背后是AI数据中心建设和电力基础设施升级对铜、锡等关键金属的持续拉动。更值得关注的是电子产业链——集成电路封装测试价格环比上涨2.9\%，外部存储设备和组件上涨1.9\%。这些数据清晰地勾勒出一条由AI需求驱动的价格传导路径：从上游的半导体制造，到中游的封装测试，再到下游的终端产品。

季节性因素也提供了额外支撑。夏季用电高峰推高了煤炭、制冷家电和电力价格，但这类因素具有明显的季节性，不会构成持久推力。全球油价波动则对上游石油加工和化工品价格形成压制，部分抵消了其他板块的涨幅。

生产者资料价格同比5.2\%，环比1.1\%，继续强势扩张；而消费者资料价格同比仍为-0.8\%，环比仅微涨0.2\%。这个剪刀差的存在，恰恰说明当前PPI上行的驱动力来自供给端的结构性变化，而非需求端的全面复苏。对于投资者而言，这意味着通胀的传导机制正在重塑，传统的“PPI-CPI传导”框架需要加入AI和产业升级这个新变量。

3. AI相关通胀的传导正在从“上游成本推升”转向“下游定价扩散”

这是JPM报告中最重要的洞察之一。研报指出，AI相关的存储通胀正在从数据中心的建设阶段，向更广泛的产业链扩散。超大规模云服务商的资本支出已经超出

[... middle omitted ...]

体现。

关键在于，这一通胀尾风可能比能源冲击更具持久性。JPM的分析暗示，即使中东局势缓和、霍尔木兹海峡重新开放，能源价格的修复也需要时间，且可能高于冲突前水平。但AI相关的通胀动力来自全球科技上行周期和中国产业升级的长期规划，这些都不是短期因素。

这意味着，市场需要重新评估中国通胀的中期路径。传统的通缩叙事可能正在被AI驱动的结构性通胀所取代，至少是在工业品和部分科技消费品领域。对于科技产业链上的公司而言，定价权的转移意味着利润结构可能发生根本性变化。

4. 报告尚未完全回答的关键问题：AI通胀的持续性究竟有多强，以及消费端何时会感受到压力

尽管JPM对AI相关通胀的传导路径给出了清晰的描述，但报告在几个关键问题上留出了讨论空间。

首先，AI相关通胀的持续性高度依赖于全球半导体周期的走势。如果全球科技资本支出在某个时间点出现回调，或者半导体产能瓶颈得到缓解，这一通胀尾风可能迅速消退。报告没有给出具体的持续时间判断，只是暗示“更持久”，但这需要后续数据验证。

\subsection{[R011] JPM：市场真正低估的不是通胀读数，而是AI成本传导正在重塑中国的价格结构}
{\small\color{refgray}归类：AI/算力：从模型竞赛到成本信任，供应链紧俏重塑定价权}

JPM：市场真正低估的不是通胀读数，而是AI成本传导正在重塑中国的价格结构

中国5月通胀数据“大致符合预期”——这是JPM最新研报的开场判断。CPI同比1.2\%，PPI同比3.9\%，与市场共识几乎一致。如果只看这个表面数字，你可能会认为这是一份平淡无奇的月度数据更新。

但这份报告真正值得关注的，不是数字本身，而是数字背后的结构性变化。JPM在看似温和的通胀读数中，捕捉到了一个正在加速的趋势：AI驱动的成本传导已经从半导体上游蔓延到消费电子终端，并在中国的工业升级周期中找到了独特的放大机制。

这不是一次性的价格脉冲。这是中国价格体系正在经历的一次供给侧再定价。

1. PPI上行中隐藏着一个比能源冲击更持久的驱动力：AI成本传导正在拓宽

5月PPI环比上涨0.8\%（季调后），同比加速至3.9\%。乍看之下，这似乎延续了过去两个月由能源价格和季节性需求推动的上行趋势。但JPM的分析揭示了一个更重要的细节：推动PPI上行的力量正在切换。

报告明确指出，本月的价格支撑来自“工业升级和AI驱动的计算需求”，具体体现在金属、电气机械和电子设备价格上。国家统计局重点提及了锡冶炼环比上涨4.8\%、铜冶炼上涨3.1\%，以及集成电路封装测试和外部存储设备分别上涨2.9\%和1.9\%。

这些数据点共同指向一个判断：AI相关的内存紧张正在从数据中心建设端向消费电子终端传导。 手机和平板电脑价格环比上涨1.6\%和1.1\%，通信工具价格环比跳升1.5\%（JPM估算季调后为1.8\%），这是AI成本传导进入终端消费品的明确信号。

与能源冲击不同——中东局势缓和后油价可能回落——AI驱动的通胀脉冲“可能更具持久性”。全球科技上行周期仍然完整，而中国的工业升级和AI+战略是多年规划。

这意味着什么？对于关注中国PPI走势的投资者而言，需要重新评估驱动因素的权重。过去两年主导PPI的能源和原材料周期正在让位于结构性更强的技术升级周期。这不是一个季度内会消退的力量。

2. CPI的温和掩盖了一个结构性缺口：消费端仍然缺乏定价权

与PPI的活跃形成鲜明对比，CPI保持温和。5月环比仅微涨0.1\%（季调后），核心CPI同比回落至1.1\%，服务CPI进一步放缓至0.8\%。

第一，食品通缩持续。鲜菜和猪肉价格继续下跌，这不仅是季节性因素。猪肉价格在经历了前几年的剧烈波动后，产能调整尚未完成，价格仍在寻底。食品项对CPI的拖累短期内难以消除。

第二，消费品PPI仍然处于通缩区间。5月消费品PPI同比为-0.8\%，尽管环比微涨0.2\%，但“需求疲弱和传导受限”的格局没有改变。这意味着上游PPI的上涨尚未有效传导至下游消费品价格。

更值得关注的是，JPM的数据显示，消费者信心指数中的收入和就业分项仍然处于低位。 年的数据点显示，消费者信心甚至低于2024年水平。在信心修复之前，消费端很难获得定价权。

*CPI和PPI的分化本质上是一个分配问题：上游涨价能否传导到下游，取决于终端需求能否承接。** 当前的数据表明，中国经济的价格传导机制仍然存在明显的断裂。

报告中的一个关键判断值得单独讨论：GDP平减指数可能在二季度转正，结束持续三年的通缩周期。一季度平减指数为-0.06\%，二季度有望转为正值。

这对市场意味着什么？平减指数转正将被视为价格复苏的宏观信号，可能支撑再通胀预期，也为央行的价格回升目标提供了数据支持。

但JPM同时给出了一个重要的限定条件：核心CPI仍然在1\%左右的低位徘徊，能源冲击对近期经济活动构成压力，关税不确定性重新浮现，美联储偏鹰派。在这些因素叠加下，央行短期内可能按兵不动。

*平减指数转正是一个技术性信号，不是趋势性拐点。** 它更多反映了PPI上行的拉动效应，而非消费端的全面复苏。如果下半年经济增长下行风险持续

[... middle omitted ...]

告提到“超大规模云服务商的资本支出正在从数据中心建设扩展到上游赋能者和较小规模的参与者”，这推动了全球半导体供应进一步收紧。但半导体供应链的紧张是周期性的还是结构性的？如果全球芯片产能投资在下半年开始释放，AI相关的价格压力可能边际减弱。

第二，中国的AI+战略能否转化为持续的终端需求。中国的工业升级和AI+计划是多年蓝图，但具体到价格传导路径，需要区分“政府投资驱动的采购”和“市场化需求驱动的消费”。前者可能带来阶段性价格脉冲，后者才能支撑持续的价格上升。

第三，消费品通缩的底部在哪里。消费品PPI仍然在-0.8\%的同比负增长区间，报告认为“产能过剩和需求疲弱限制了传导”。但产能过剩的程度、需求修复的时点，报告没有给出明确的量化判断。

这些未解问题决定了AI成本传导的实际影响。如果上述假设被验证为正确，那么中国PPI可能在更长时间内维持高于市场预期的水平，而CPI则可能滞后但最终跟随。如果假设不成立，那么当前PPI的上行可能只是能源和季节性因素叠加的短期现象。

\subsection{[R012] Bernstein：AI模型商品化进程比市场想象的更快，竞争格局正从智能竞赛转向成本与信任竞赛}
{\small\color{refgray}归类：AI/算力：从模型竞赛到成本信任，供应链紧俏重塑定价权}

Bernstein：AI模型商品化进程比市场想象的更快，竞争格局正从智能竞赛转向成本与信任竞赛

过去一周，AI行业发生了三件看似独立、实则指向同一方向的事件：Anthropic的Claude Fable 5定价高到让开发者开始理性配给token用量；OpenAI被报道正在降价以争夺市场份额；腾讯、微软等大型AI用户开始重新审视“token最大化”策略的经济性。这些信号叠加在一起，指向一个Bernstein在最新研报中明确提出的判断：AI模型的商品化进程，正在被人类用户的感知边界加速，而非由底层模型智能的收敛速度决定。

这份由Bernstein全球互联网团队发布的报告，提供了一套全新的分析框架——不是从模型能力本身出发，而是从“什么情况下人类用户认为AI已经足够好”这一视角，重新定义AI市场的竞争逻辑。报告的核心主张是：一旦某个垂直场景的AI任务完成质量达到“可靠可用”的阈值，竞争的决定性因素将从推理能力转向成本、可用性和开发者信任。这将对全球AI实验室的商业模式、中国AI厂商的市场空间、以及投资者评估AI资产的方式产生深远影响。

1. AI商品化的真正驱动力不是技术收敛，而是人类用户的“够用”感知

Bernstein报告提出的核心框架，与市场上主流讨论的“模型智能趋同”理论有本质区别。当前流行的观点认为，随着模型能力提升，不同AI模型之间的智力差距会逐渐缩小，最终导致商品化。但Bernstein认为，这种技术收敛的假设既无法验证，也不重要。

真正驱动商品化的，是人类用户在不同使用场景中对AI表现的主观感知。当一个AI系统能够可靠地完成特定任务——比如订奶茶、预订酒店、生成营销文案——用户就不再关心底层的模型架构是否与前沿模型一致，而只关心它是否好用、便宜、稳定。

报告用了一个非常直观的比喻：一旦某个场景被“解决”，继续在该场景上投入研发的边际回报急剧下降。这意味着AI实验室可以将研发资源重新分配到更复杂的前沿任务上。从经济学角度看，这正是“从指数增长曲线中走出来”的关键机制——值得投入指数级成本的场景范围会逐渐收窄，训练成本增速放缓，AI实验室才有机会展现经营杠杆。

2. 消费者端AI将最先商品化，企业端和前沿科学场景的竞争逻辑截然不同

Bernstein对AI市场进行了清晰的三层分割，每一层的竞争逻辑和商品化时间表都不同。

最底层是消费者端的AI应用。报告明确指出，通过微信智能助手订奶茶、在淘宝上用AI购买机票、通过AI生成短视频内容——这些场景正在快速达到“可靠可用”的阈值。Bernstein判断，消费者端AI将是第一个被广泛商品化的领域。在这个层面，用户对模型差异的感知最弱，对价格和响应速度最敏感。这意味着，谁能在消费者场景中提供成本最低、体验最流畅的AI服务，谁就能占据主导。

中间层是企业级通用场景。从Excel建模、代码生成到营销策略的头脑风暴，这些场景对输出的确定性要求更高，商品化速度会慢于消费者端。但报告认为，一旦企业用户建立了对特定AI工具的信任，切换成本会很高，因此先发优势在中间层尤为重要。

最顶层是前沿科学场景——药物发现、核聚变、太空探索、国防等。这些领域的用户对AI能力有近乎无限的支持意愿，且需求高度专业化、主权化。Bernstein认为，这个市场将长期由美国的前沿实验室主导，并且可能永远不会完全商品化。

这个分层框架的意义在于：投资者不能再用单一的“AI赛道”视角来评估公司。不同层级的竞争壁垒、定价权和增长前景截然不同。

3. 中国AI厂商的全球可及市场约为35-40\%，但竞争逻辑是“够用+低价”

报告对中国AI厂商在全球市场中的定位给出了一个既务实又不乏想象空间的判断。Bernstein认为，中国AI模型在推理能力上与美国前沿模型存在6-12个月的差距，但

[... middle omitted ...]

恰是中国模型“够用”优势最明显的领域。

4. 报告尚未完全回答的关键问题：AI实验室何时能实现独立盈利？

Bernstein对AI实验室盈利前景的判断是“谨慎乐观”，但这份乐观建立在一个尚未验证的假设之上：一旦商品化进程启动，研发投入的增速会自然放缓，从而让经营杠杆得以显现。

这个逻辑链条存在几个需要进一步验证的环节。首先，商品化是否会真正降低研发投入？还是说，竞争会迫使实验室继续在多个前沿方向同时投入，导致总研发开支不降反升？其次，即使某个场景被“解决”，实验室是否有动力主动降低该场景的定价？还是说，他们会利用品牌溢价和用户粘性维持高利润率？最后，开源模型的竞争是否会压缩所有厂商的利润率，使得“足够好”本身成为一个无法盈利的定位？

报告没有给出这些问题的明确答案。这恰恰是投资者最需要关注的风险点——AI资产的估值，很大程度上取决于市场对AI实验室长期盈利能力的预期。如果商品化进程不能带来盈利改善，而是演变成一场持续的价格战，那么当前的高估值可能需要重新审视。

\subsection{[R013] GS：铝市场真正超预期的不是中国产量，而是海外供给的加速释放}
{\small\color{refgray}归类：能源与大宗：需求破坏对冲供给冲击，铝与原油的结构性变局}

GS：铝市场真正超预期的不是中国产量，而是海外供给的加速释放

全球铝市场正在经历一个被多数投资者低估的结构性变化。多数市场参与者仍在争论中国4500万吨产能红线能否守住，但GS最新一场专家电话会揭示了一个更关键的信号：海外铝冶炼产能的扩张速度正在显著加快，其体量和节奏远超此前的市场预期。

这份报告的核心判断是：未来三年全球铝市场面临的主要矛盾，将从“中国供给约束”转向“海外产能加速释放”。印尼和中亚正在成为中国铝企海外扩张的核心战场，2026年至2028年间，仅这两个地区的铝产能增量就可能达到近600万吨。这一数字意味着什么？它相当于全球铝市场规模的5\%以上，足以改写未来几年的供需平衡表。

为什么这个判断现在重要？因为市场对铝的定价逻辑仍停留在“中国减产-供给偏紧”的旧框架中。但GS专家在实地调研后的结论是：中国产量已持续超过4500万吨，而海外产能的扩张正在形成新的供给曲线。如果投资者只盯着中国的环保督察和产能置换，可能会错过更大的结构性变量。

1. 印尼正在成为全球铝供给的新变量，产能释放节奏远超年初预期

GS专家在近期访问印尼后指出，印尼铝冶炼产能的扩张条件正在显著改善。最核心的变化是电力瓶颈的解除。在青山工业园，部分暂停的镍冶炼产能被重新调配给铝冶炼使用，释放了大量电力资源。与此同时，当地电厂建设的进度好于预期，包括710MW超临界火电机组的顺利出口，都为铝冶炼提供了稳定的能源保障。

基于实地调研，专家将印尼2026年的铝产能增量预期从年初的142万吨大幅上调至215万吨，2027年进一步上调至197.5万吨。这意味着仅印尼一地，两年内的产能增量就将超过400万吨。更重要的是，这些项目的投产节奏将从2026年四季度开始明显加速。

具体来看，几个关键项目的进展值得关注：项目B的一期400kt产能已于2026年5月投产，二期65kt也已启动，三期扩建正在等待电力配套；项目D的一期500kt预计2027年下半年投产；项目C的800kt产能计划在2026年底投产。这些项目不再是纸面规划，而是已经进入实质建设阶段。

这里的关键含义是：印尼铝产能的扩张不再是“可能”或“预期”，而是正在发生的现实。对于全球铝的定价者而言，这意味着未来两年的供给曲线需要被重新绘制。

2. 中亚正在成为第二个海外产能聚集地，中国企业的布局比想象中更大

如果说印尼是中国铝企海外扩张的第一站，那么中亚正在迅速成为第二极。GS专家观察到，中国生产商在哈萨克斯坦的铝产能建设正在加速推进，规划总产能达到540万吨。

两个项目构成了这一轮扩张的主体。项目G规划了三期共240万吨产能，一期80万吨预计2028年投产，目前已完成选址并开始打桩。项目H的规划更为宏大，总产能300万吨，同样分三期建设，一期100万吨也计划在2028年投产。这两个项目合计，仅2028年就有180万吨产能计划落地，相当于全球市场规模的2.4\%。

值得思考的是，为什么是中国企业选择在哈萨克斯坦大规模投建铝产能？答案可能在于能源成本和地缘优势。中亚地区拥有丰富的天然气和煤炭资源，电力成本具有竞争力，同时地理位置靠近欧洲和亚洲两大消费市场。对于中国企业而言，在海外布局产能不仅能规避国内的产能天花板，还能更灵活地服务全球客户。

但这里有一个需要验证的问题：这些项目的建设周期是否可靠？从当前进度看，项目G已进入实质性施工阶段，但项目H仍处于早期规划。报告并未完全展开这些项目面临的融资、政策和技术风险。投资者需要持续跟踪项目进展，而不是简单线性外推。

3. 中国实际产量已持续高于4500万吨，环保督察并未改变这一趋势

市场一直在争论中国4500万吨的产能红线是否会被打破。GS专家的判断是：这一红线实质上已经被突破。根据专家的估算，2026年中国原铝产量将达到

[... middle omitted ...]

的可能性是存在的。

这个判断对市场定价的含义是什么？如果中国产量持续高于4500万吨，那么全球铝市场的供给过剩压力将比预期更大。投资者需要重新评估“中国供给约束”这一长期叙事是否仍然成立。

4. 出口订单的强劲增长正在改变国内库存节奏，但可持续性存疑

GS专家注意到，2026年4-5月以来，中国铝产品出口订单出现了显著增长，尤其集中在未锻轧铝、铝合金、铝板带和铝型材等产品。这一变化将推动中国铝出口实现两位数增长，并可能从6月开始带动国内社会库存的去化。

出口订单的强劲增长，背后是海外市场对中国铝产品的旺盛需求。在全球铝价高企的背景下，中国铝产品的价格竞争力凸显，出口套利空间持续存在。但这里有一个关键问题：这种出口动能能否持续？

报告没有完全展开的是，海外铝产能加速释放后，中国铝出口的高增长可能面临挑战。当印尼和哈萨克斯坦的新产能逐步投产，这些地区的铝产品可能优先满足当地和周边市场需求，从而挤压中国产品的出口空间。此外，贸易摩擦和关税政策的变化也是不确定因素。

\subsection{[R014] 摩根斯坦利：中国原油进口的骤降并非需求崩塌，而是库存周期的主动调节}
{\small\color{refgray}归类：能源与大宗：需求破坏对冲供给冲击，铝与原油的结构性变局}

摩根斯坦利：中国原油进口的骤降并非需求崩塌，而是库存周期的主动调节

当全球市场都在关注霍尔木兹海峡的供应冲击时，一个关键变量正在被误读。中国在3月至5月期间，将原油进口量从2月的约12.8百万桶/日骤降至5月的约7.8百万桶/日，降幅接近500万桶/日。这一数字足以让任何大宗商品分析师警觉：全球最大的原油进口国是否正在经历需求端的断崖式下跌？

摩根斯坦利在最新发布的《石油数据摘要：欧洲》专题报告中，给出了一个反直觉但逻辑严密的判断。该机构大宗商品策略师Charlotte Firkins与Martijn Rats通过拆解中国的原油贸易、炼厂运行以及终端需求数据，揭示了一个被市场普遍忽略的事实：中国进口的锐减，并非因为经济突然失速或需求崩溃，而是在一个持续了数月的库存过剩周期中，主动进行的调节性收缩。

这份报告的价值不在于它告诉市场中国需求在下降——这已经是共识——而在于它精确地刻画了下降的性质、结构以及背后的驱动力。对于试图理解全球油市再平衡路径的投资者而言，这组数据提供了远比“需求疲软”四个字更丰富的分析框架。

1. 进口暴跌与库存反增的悖论，指向的是此前的过度进口而非当前的需求坍塌

市场最容易犯的错误，是将进口量的变化直接等同于消费量的变化。摩根斯坦利的数据清晰地展示了这一逻辑陷阱。5月中国原油进口降至约7.8百万桶/日，同比下降3.2百万桶/日，较2月危机前水平更是低了近500万桶/日。单看这个数字，很容易得出“中国需求大幅萎缩”的结论。

但报告提供了一个关键的反证：根据Vortexa的可观测原油库存数据，中国的商业和战略原油库存在3月和4月竟然还在持续累积，平均建库速度约为1.1百万桶/日。直到5月，库存才开始转为去化，平均速度约为75万桶/日。

进口大幅下降的同时库存仍在增加，这一悖论只有一个解释：在冲突爆发前的2025年大部分时间以及2026年初，中国一直在超额进口原油，大规模建立商业和战略库存。报告明确指出，在2026年2月，中国的原油平衡表已经处于约2.7百万桶/日的过剩状态，因为进口量远超炼厂需求。

因此，3月至5月的进口削减，是在一个已经严重过剩的库存背景下发生的。这不是需求端的主动收缩，而是供给端的被动调整。4月的完整政府数据进一步印证了这一点：当月国内产量（4.4百万桶/日）加上净进口（9.3百万桶/日），仍然略微超过了炼厂需求（13.4百万桶/日）。库存微增的数据与这一计算方向一致。

这意味着，市场如果仅凭进口数据就下调对中国需求的预期，很可能高估了实际需求的降幅。中国正在消化的是此前积累的库存，而非需求已经永久性失速。

2. 炼厂利润深度转负是进口削减的直接推手，国有企业与独立炼厂的行为出现结构性分化

如果说库存过剩是进口削减的背景，那么炼厂利润的崩塌则是直接触发机制。摩根斯坦利的报告揭示了这一传导链条的清晰逻辑。

到4月，中国炼厂的综合利润已经全面转负。原因有三重叠加：高昂的原油采购成本、政府对成品油出口的限制，以及国内成品油价格调整机制未能完全传导成本上涨。这三股力量共同挤压，使得炼厂每加工一桶原油都在亏损。

面对亏损，两类炼厂做出了截然不同的反应。国有企业（State-owned refiners）选择了激进地削减开工率。4月中国炼厂总开工量较2月下降了1.8百万桶/日，同比下降0.8百万桶/日。国有炼厂是此次减产的主力。它们有更强的财务纪律和更少的政治压力，能够迅速根据利润信号调整运营。

而独立炼厂（Independent refiners）则扮演了截然不同的角色。报告指出，这些炼厂被政府要求不得削减开工率，因此承担了更多供应国内市场的责任。这一结构性分化意味着，中国的原油进口削减并非均匀分布。独立炼厂由于被强制维持生产，其原料需求（尤其是燃料油等替

[... middle omitted ...]

）显示，4月总需求同比下降1.2百万桶/日，降幅为8\%。但这一数字需要仔细拆解。

需求下降的主要驱动力来自LPG和燃料油，这两项贡献了约70\%的同比降幅。这背后的逻辑是清晰的：石化装置（petchem facilities）削减了开工率，而独立炼厂则减少了对燃料油作为原料的使用。这是工业端的调整，而非消费端的萎缩。

真正值得关注的是运输燃料市场。报告提供了两个看似矛盾的数据点：一方面，中国的出行活动数据（如百度拥堵指数和航班数量）依然强劲，甚至同比略有增长；另一方面，柴油和汽油的表观需求却呈现疲软态势。

答案在于替代效应的加速。报告引用了Argus的估算数据：2026年一季度，新能源和LNG重卡对柴油的替代量同比增加了23万桶/日。LNG重卡在4月的销售渗透率已达到26\%，而电动重卡的销售在5月同比增长超过90\%，渗透率从1月的约30\%跃升至5月的约40\%。在乘用车领域，新能源汽车渗透率在3月达到51\%，4月进一步攀升至61\%，而传统燃油车销量同比萎缩了37\%。

\subsection{[R015] GS：CLO套利空间的真正故事不在利差本身，而在资本结构如何重塑定价权}
{\small\color{refgray}归类：固定收益与信用：CLO套利空间收窄与信用分化中的alpha机会}

GS：CLO套利空间的真正故事不在利差本身，而在资本结构如何重塑定价权

CLO（担保贷款凭证）市场正在经历一个被多数投资者低估的结构性转变。传统上，市场参与者衡量CLO股权层吸引力时，习惯于计算新发杠杆贷款利差与CLO债务加权平均成本之间的差值。这个指标直观、易得，但GS最新发布的研报指出，它遗漏了一个关键维度：交易层面的实际超额利差。

GS使用Intex的交易级和分层级数据，构建了一套“交易匹配套利利差”指标，覆盖美元和欧元CLO。结论清晰且反直觉：美元CLO套利空间并非只是周期性波动，而是出现了更显著的趋势性收窄。这背后，一个重要的结构性驱动力是“自有股权资本”的崛起。这种资本改变了CLO股权投资者的行为模式，使得市场对即期套利条件的敏感度下降。对于决策者而言，理解这一变化，远比盯着利差数字本身更重要。

这份报告的价值不仅在于提出了一个更精确的测量工具，更在于它揭示了CLO市场定价权的转移——从依赖市场波动获取超额收益，转向依赖规模、品牌和持续部署能力。以下是我们从报告中提炼出的四个核心洞察。

1. 美元CLO套利空间的结构性收窄，根源在于“自有股权资本”改变了行为逻辑

GS构建的交易匹配套利指标显示，美元CLO套利空间在过去两年收窄幅度远超传统指标所反映的程度。传统指标可能显示利差在某个区间波动，但新指标清晰地画出了一条向下的趋势线。报告认为，这并非简单的周期性现象，而是一个结构性的变化。

驱动这一变化的关键变量是“自有股权资本”的增长。这类资本通常承诺跨年份部署，并将管理权和投资决策权交给发起机构。它们追求的不是单笔交易的最大化入场利差，而是连续部署和年份分散化。这意味着，即使即期套利空间收窄，这些资本也不会轻易退场。这实际上降低了股权需求对利差变化的弹性。

对于独立股权投资者而言，这意味着一个更严峻的竞争环境。非回报导向的参与者增多，会系统性压低股权层回报。GS明确指出，“更低的起始套利空间减少了股权投资者可用于吸收信用损失的安全垫。”这不是一个短期现象，而是市场成熟化后的新常态。决策者需要重新评估CLO股权投资的回报预期，并意识到，过去依赖高利差获取超额收益的模式正在变得不可持续。

2. 欧元CLO看起来更便宜，但“免费午餐”背后是更弱的基本面

与美元市场形成对比的是，欧元CLO的套利空间仍相对较宽。这容易让投资者产生“欧元股权更便宜、更有吸引力”的错觉。GS的分析对此提出了审慎的警告。

报告指出，欧元套利空间更宽，部分原因是“自有股权资本”尚未在欧元市场形成同等规模的影响力。这确实在估值层面提供了吸引力。然而，关键矛盾在于欧元区贷款市场的基本面更弱。过去两年，欧元区违约率整体上升，而美元违约率已从2024年的高点回落。GS预测，由于政策和宏观环境组合不利，欧元区违约率可能继续小幅攀升。与此同时，欧元区贷款价格相对于美元同类资产已经偏贵。

这意味着，欧元CLO表面上的超额利差，有相当一部分需要用于补偿潜在的信用损失。报告特别提到，尽管软件等轻资产企业在美元贷款市场中占比更高，但欧元市场同样存在显著的轻资产敞口，且2028年将迎来再融资到期高峰。因此，投资者不应将更宽的套利空间简单视为机会，而应将其视为对更高信用风险的一种定价。对于CLO债务投资者，GS的分析则相对积极，认为欧元投资级分层的利差相对于美元同类资产仍具吸引力，这为负债端提供了不同的机会。

3. 重置交易的经济效益远非“负债成本下降”那么简单，美元和欧元市场分化明显

CLO管理者经常通过重置或再融资来降低负债成本。传统观点认为，这能直接改善交易的经济效益。但GS的交易匹配框架揭示了一个更复杂的图景。

在美元市场，自2023年以来，贷款市场的重定价浪潮极为猛烈。这意味着，当管理者重置CLO时，不仅负债成本在下降，

[... middle omitted ...]

二，贷款供给对再投资更有利；第三，自有资本更倾向于重置CLO以延长再投资期并继续赚取管理费，即使在更紧的交易条件下。

这一发现对于CLO管理者和投资者至关重要。它意味着，在美元市场，单纯依赖重置来改善交易经济性的策略已基本失效。管理者需要寻找其他价值创造来源，例如更精准的资产选择、更高效的资本结构设计，或者转向更有利可图的细分市场。

4. 管理者分层揭示了一个反直觉的真相：更低的负债成本并不必然牺牲超额利差

GS根据市场定价对CLO管理者进行了分层。一个常见的假设是，负债成本更低的管理者（通常是大型、声誉卓著的机构）必然在资产端接受了更低的利差，从而使得不同层级管理者的套利空间大致相同。但数据否定了这一假设。

在美元市场，高等级管理者（Tier 1）实际上实现了略宽的套利空间。这意味着，他们能够凭借更强的品牌和发行能力，在获得更低负债成本的同时，并未牺牲资产端的超额利差。这背后是美元市场管理者在负债端的巨大分化：顶级管理者的AAA级分层利差远低于低等级管理者。

\subsection{[R016] GS：市场低估的不是地缘冲突的烈度，而是需求破坏的持久性}
{\small\color{refgray}归类：能源与大宗：需求破坏对冲供给冲击，铝与原油的结构性变局}

当霍尔木兹海峡的通行量降至战前水平的不足10\%，布伦特原油价格却从三月末的高点回落了约25\%。这一看似矛盾的价格走势，正在挑战一个根深蒂固的市场直觉：地缘供给中断必然导致油价飙升。GS最新发布的这份研报提供了一个关键判断——市场真正需要重新定价的，不是冲突升级的风险溢价，而是在供给冲击持续期间，需求端正在发生的结构性坍塌。这份报告的核心主张是，一个更持久的霍尔木兹中断，已被一个更小的供需赤字所抵消，而后者主要源于需求破坏的规模远超预期。

这份研报维持了对2026年第四季度布伦特原油每桶90美元的预测。维持预测本身并不惊人，真正值得关注的是维持预测的理由发生了根本性置换：此前市场担忧的是供给中断的持续时间，而GS现在指出，一个更小的赤字（5-6百万桶/日）抵消了一个更长的中断（恢复时间从六月底推迟至八月底）。这意味着，油价的地板不再由供给侧的物理瓶颈决定，而是由需求端的脆弱性重新锚定。

对于产业决策者和资产配置者而言，理解这种定价逻辑的转移，比预测下一个地缘事件更为紧迫。

GS估算，中东液体燃料生产在二季度受到的冲击高达14-15百万桶/日。然而，全球石油市场同期出现的供需赤字仅为5-6百万桶/日。两者之间近9百万桶/日的差距，揭示了市场缓冲机制的深度。

这一差距由三个因素构成。首先，战前市场本就存在约3百万桶/日的供给过剩。其次，中东以外地区的供给在冲突期间反而上调了约1.5百万桶/日，主要来自美洲（美国、巴西、圭亚那、委内瑞拉）的增产。但最大的变量，是需求侧。GS估算，全球石油需求因冲突而损失的规模接近5百万桶/日。这意味着，供给中断的冲击，有超过三分之一被需求端的同步萎缩所吸收。

需求破坏的集中度值得关注。报告指出，中国原油进口需求近期出现了接近4百万桶/日的同比下滑。这并非简单的经济放缓，而是中国向电动车等替代能源转型加速的结构性信号。如果这一趋势在冲突结束后仍有一部分（报告假设略高于10\%）持续存在，那么供给恢复后的市场将面临一个永久性的需求缺口。

这里的含义是清晰的：传统的“供给中断等于油价暴涨”的线性模型已经失效。市场正在学习一个更复杂的函数，其中需求弹性在危机期间可能远高于正常时期。

GS将海湾国家石油出口恢复正常的时间预期从六月底推迟至八月底。这一调整基于对霍尔木兹海峡通行量的最新估算：即使考虑通过延布、富查伊拉、杰伊汉等替代路线的分流，要使总出口恢复到战前23百万桶/日的水平，霍尔木兹的通行量仍需从当前水平回升至战前水平的约70\%。

但报告同时提供了一个重要的乐观视角：一旦海峡重新开放，产量恢复的速度可能快于市场共识。GS的预测显示，在重新开放后的几个月内，海湾国家的原油产量将迅速反弹，其恢复曲线比国际能源署（IEA）和美国能源信息署（EIA）的预测更为陡峭。支撑这一判断的关键证据是，沙特阿拉伯等国的活跃钻机数量在冲突期间不降反升。这意味着生产商并未停止准备，而是在为重启进行预投资。

报告将管道运力视为恢复速度的主要约束，而非材料、劳动力或油井流量。这一判断与许多市场参与者的直觉不同，后者往往更关注物理基础设施的损坏。如果恢复速度确实快于共识，那么当前期货曲线中隐含的远期溢价可能被高估。

3. 2027年的价格韧性并非来自基本面，而是来自库存结构和风险溢价

报告将2027年布伦特原油均价预测下调5美元至每桶80美元，但同时指出，这一价格水平仍比2025年均价高出约10美元，尽管同期市场将出现超过3百万桶/日的供给过剩。

这种看似矛盾的价格韧性来自两个非基本面因素。第一，库存结构。由于2026年的大规模库存消耗，经合组织（OECD）商业石油库存不太可能回到很高水平。叠加2027年全球超过1百万桶/日的战略储备补库需求，库存缓冲的缺乏将为价格提供一个底部支撑。第二，安全溢价

[... middle omitted ...]

报中最引人深思但也最不确定的部分。GS假设，冲突结束后，仅有略高于10\%的需求损失会持续存在。但这一假设的论证基础在报告中并未充分展开。

需求破坏的性质是复杂的。一部分来自经济活动放缓，这部分在冲突结束后可能迅速恢复。另一部分来自行为改变和替代加速，例如中国电动车的渗透率提升、中东地区工业活动的调整、以及全球供应链的重构。后者具有更强的粘性。报告提到了中国原油进口的急剧下滑，但未能深入拆解其中有多少是临时性的库存去化，有多少是永久性的结构替代。

此外，报告对需求破坏的估算（接近5百万桶/日）本身带有高度不确定性。在冲突环境下，数据收集和质量都受到限制。如果实际需求破坏小于估算，那么当前的赤字可能被低估，价格存在上行风险。反之，如果需求破坏的持续性超出预期，即使供给恢复，市场也可能面临一个长期的需求不足格局。

这份报告为读者提供了一个更实用的分析框架，而非简单的价格预测。核心在于，不要将注意力过度集中在霍尔木兹何时重新开放的单一变量上，而应建立一个两维度观测体系。

\subsection{[R017] 摩根斯坦利：北美油气高管薪酬设计已传递出比产量更重要的信号}
{\small\color{refgray}归类：未被主题摘要引用}

2026年北美油气公司的委托书（proxy filing）刚刚披露完毕。对于关注这个行业的投资者和决策者来说，这些文件通常被视为例行公事——一堆关于薪酬、治理和投票事项的合规性披露。但摩根斯坦利这份最新研报揭示了一个被市场低估的信号：高管薪酬框架的稳定本身就是一个重要的战略声明。

这份报告的核心判断是：北美油气公司董事会正在通过薪酬设计，不动声色地锁定行业纪律，而不是为下一轮增长周期做准备。如果你还在用产量增长或资本开支计划来衡量这些公司的战略意图，你可能错过了真正重要的结构性变化。

为什么现在需要关注这个信号？过去两年，市场对油气板块的定价逻辑经历了剧烈摇摆。从2022年的暴利到 年的回归均值，投资者的核心焦虑已经从“这些公司能不能赚钱”转变为“这些公司会不会把钱花掉”。2026年的委托书披露，恰恰提供了检验这一焦虑的窗口——董事会是否真的站在投资者这边。

摩根斯坦利团队分析了覆盖范围内的北美油气公司2026年短期激励（STI）计划，发现了一个值得深思的结论：整体框架几乎没有变化。这种稳定不是惰性，而是一种刻意的选择。

2026年与2025年相比，北美油气公司短期激励指标的加权平均权重变化幅度极小。美国生产商的核心指标——EHS（环境、健康与安全）、现金流、资本效率——权重分别为18\%、18\%和15\%，与上年几乎完全一致。加拿大同行则更侧重EHS（24\%）、生产/交付（16\%）和运营效率（14\%），同样保持稳定。

这种稳定性的含义需要仔细拆解。在一个能源价格波动、能源转型压力持续、投资者对资本纪律高度敏感的行业环境中，董事会选择不调整薪酬框架，意味着他们相信现有的激励结构已经足够有效地引导管理层行为。这不是一个“没有变化”的信号，而是一个“我们没有理由改变”的信号——而这恰恰是资本纪律得以持续的制度性保障。

更值得注意的是，自2019年以来，行业整体的STI指标权重已经发生了根本性迁移：生产/交付指标的权重从18\%降至5.5\%，战略/自由裁量指标从14\%降至9\%，而EHS从6\%升至17.5\%，资本效率从11\%升至14.5\%。2026年的稳定，是在这一结构性转变基础上的巩固，而非停滞。

对于投资者而言，这意味着：如果你持有的公司正在悄悄增加产量激励权重，那可能是一个需要警惕的早期信号。但2026年的数据显示，这个风险整体上并未出现。

这份研报最有趣的发现之一，是美加两国生产商在短期激励设计上的系统性差异。美国公司更看重现金流（18\%）、资本效率（15\%）和回报指标（12\%）；加拿大公司则更强调EHS（24\%）、生产/交付（16\%）和运营效率（14\%）。

这种差异并非偶然。它反映了两个市场在股东结构、监管环境和行业文化上的深层不同。美国独立生产商的股东基础更加注重短期资本回报和自由现金流生成，因此激励设计直接挂钩这些指标。加拿大生产商则面临更严格的ESG监管环境和更集中的机构股东结构，因此EHS和运营效率获得了更高的权重。

对于跨市场配置的投资者，这是一个重要的观察框架：美国公司的薪酬设计天然更倾向于保护股东回报，而加拿大公司则可能在运营纪律和ESG表现上更有优势。但两者在同一个核心方向上是一致的——都不再鼓励盲目的产量增长。

3. 长期激励占比75\%意味着CEO的注意力已经被锁死在长期价值创造上

研报指出，高管薪酬中平均约75\%来自长期激励（LTI），部分公司如EQT和PR的CEO薪酬甚至100\%为浮动、基于绩效的薪酬。这是一个经常被低估的数字。

当一家公司的CEO超过四分之三的薪酬取决于长期绩效而非年度奖金时，其行为逻辑会发生根本性变化。短期产量冲刺、资本开支突击、季度盈利管理——这些行为的激励基础被大幅削弱。取而代之的是，CEO有更强的动机去关注资源禀赋的优化、资产负债表的结构性改

[... middle omitted ...]

倍。

但这个平均值掩盖了巨大的分散性。从最低的0.8倍到最高的5.7倍，不同公司的CIC条款差异显著。这种差异直接影响了并购交易的定价逻辑：一家CIC倍数较高的公司，其CEO在谈判中可能更倾向于接受收购要约，因为个人财务利益与交易完成高度绑定；反之，CIC倍数较低的公司，CEO可能更倾向于维持独立运营。

对于正在评估潜在并购标的的投资者，这个数据点提供了一种超越传统估值指标的分析维度。一家公司的CIC条款设计，实际上透露出董事会和管理层对独立运营价值的真实判断。

5. 报告尚未完全回答的关键问题：薪酬纪律能否在下一轮油价上涨中保持韧性

这份研报提供了2026年的静态快照，但它也留下了一个未解的问题：当前的薪酬纪律能否经受住下一轮油价大幅上涨的考验？

历史经验表明，油气行业的资本纪律通常在价格低迷时期最容易维持，而在价格高涨时期最容易崩溃。 年后的行业整合和成本削减，在 年的油价反弹中一度面临考验。虽然整体纪律得以维持，但部分公司确实出现了资本开支上升的迹象。

\subsection{[R018] 美国银行：市场低估了房地产信用分化中的结构性机会，而非整体复苏}
{\small\color{refgray}归类：地产与消费：分化中的结构性机会与成本端出清 / 固定收益与信用：CLO套利空间收窄与信用分化中的alpha机会}

美国银行：市场低估了房地产信用分化中的结构性机会，而非整体复苏

这份美国银行最新发布的大中华区房地产研报，表面上呈现的是6月香港一手房成交放缓、高收益债收益率小幅走阔等短期波动信号。但如果把报告中的三个独立模块——香港住宅销售、中央财政支持城市更新、以及信用研究员对个别债券的明确超配建议——放在一起看，真正的核心判断并非“市场在变差”，而是“市场正在从总量博弈转向个券和个票的alpha挖掘”。

报告的真正价值不在于描述了6月的数据有多弱，而在于它指出了一个正在发生的结构性转变：政策子弹的边际效用正在递减，但信用主体的分化正在加速。对于高净值投资者和产业决策者来说，这意味着过去两年“买整个板块”的交易策略正在失效，取而代之的是一种更精细、更依赖主体信用判断的配置逻辑。

1. 6月香港一手房成交放缓不是需求坍塌，而是开发商主动的供给节奏调整

报告显示，截至6月，香港一手房成交仅约240宗，而过去五个月的月均成交量超过2000宗。单看这个数字，很容易得出“市场急冻”的结论。但美国银行的分析师将原因归结为三个方面：本地股市下跌带来的负财富效应、对美国加息的重新担忧、以及中国内地新出台的离岸投资规则带来的不确定性。

这里的关键在于，报告明确指出开发商正在“更谨慎地控制推盘节奏”。这不是需求端的全面崩溃，而是供给端在多重不确定性下的主动收缩。从投资角度看，这意味着香港楼市的短期成交量数据可能失真，不能简单线性外推为价格趋势。真正需要关注的是，当开发商重新加速推盘时，市场能否消化——而这取决于上述三个不确定性因素的演变。

报告也坦承，离岸投资规则的不确定性“短期内将持续存在，直到更清晰的实施细则出台”。这是一个尚未落地的核心变量。

2. 中央2570亿城市更新资金看似庞大，但对行业总盘子的增量影响“相对较小”

报告披露了一个重要数字：2026年中央预算将拨付970亿元人民币，加上1600亿元超长期特别国债，总计2570亿元用于支持地方政府的城市更新。但紧接着，报告给出了一个关键判断：相对于每年超过3万亿元的城市更新总投资，这笔资金的增量影响“相对较小”。

这个判断背后的逻辑链条是：大头资金仍然需要地方政府自己解决，包括专项债和政策性银行贷款。中央资金的定位更像是“撬动杠杆”而非“兜底”。对于投资者而言，这意味着不能把这条新闻解读为房地产行业的重大利好。真正需要跟踪的是地方政府专项债的发行节奏和政策性银行的配套贷款规模，而非中央预算的绝对数字。

报告没有展开的是：这2570亿的分配机制和项目筛选标准是什么？哪些城市或哪些类型的项目会优先受益？这些细节才是决定投资标的的关键。

3. 高收益债收益率YTD收窄162bp，但过去一周走阔18bp，市场正在重新定价尾部风险

报告数据显示，剔除万科后，中国高收益地产债的YTD平均收益率已收窄162个基点至9.58\%，但过去一周走阔了18个基点。这个数据组合透露的信息是：年初以来的信用修复趋势没有逆转，但边际上出现了新的压力。

更值得关注的是，报告提到“过去一周该板块没有新发行”，YTD净赎回维持在14亿美元。这意味着一级市场的大门对于高收益地产发行人仍然是关闭的。再融资风险并没有消失，只是被年初以来的二级市场上涨暂时掩盖了。

美国银行信用研究员在这一背景下做出的超配建议——对LASUDE‘26——恰恰说明了个券选择的重要性。报告给出的超配理由包括：基于债务合约的条款仍有可观的上行空间、利息覆盖倍数约1倍、以及公司通过资产出售主动去杠杆。这不是一个板块性的推荐，而是一个高度针对性的个券判断。

4. 万科海外资产违约与国内债务展期同时发生，信用分化正在从“能不能还”变为“愿意用什么还”

报告提到了两个关于万科的独立信息：一是万科获得银行批准展期总额11.84亿

[... middle omitted ...]

以保全核心。这个判断框架，比任何财务比率都更能预测信用事件的走向。

5. 南向资金对内地房企和香港本地股的持仓出现微妙背离，资金流向正在发出信号

报告中的Exhibit 1是一张南向资金持仓变化表，数据截至5月底。这张表揭示了一个被市场忽视的细节：内地房企的平均南向持股比例从4月的14.13\%微降至5月的14.08\%，基本持平；但香港本地地产股的平均南向持股比例从4月的2.67\%进一步降至2.50\%，延续了年初以来的下降趋势。

具体来看，内地房企中，龙湖、越秀、绿城等优质主体的南向持仓仍在增加，而保利置业、深圳控股的持仓下降明显。香港本地股中，几乎所有的标的——新世界、嘉里、长实、新鸿基、恒隆等——都出现了南向减持。

这个背离的隐含含义是：南向资金正在从香港本地地产股撤出，但对内地房企采取的是“结构性调仓”而非系统性减持。优质国企和民企龙头仍在获得边际增持。这印证了我们在前文提到的观点：资金正在从“买香港”转向“买中国”，且在内地板块内部也在进行更精细的筛选。

\subsection{[R019] BARC：校园网络市场正在回归正常增长，真正的竞争焦点是份额再分配}
{\small\color{refgray}归类：未被主题摘要引用}

BARC：校园网络市场正在回归正常增长，真正的竞争焦点是份额再分配

市场已经消化了疫情后混合办公带来的网络设备需求爆发，也消化了2024年的库存修正。但BARC这份最新发布的校园网络模型更新报告揭示了一个更值得关注的信号：未来两年市场增速将温和回落至7\%-8\%，而真正拉开公司间估值差距的，不再是行业贝塔，而是结构性份额迁移。

这份报告的核心判断可以简化为一句：校园网络市场正在从“周期修复”转向“常态增长”，而在这个阶段，赢家不是靠市场涨潮，而是靠从竞争对手手中夺取份额。

对于关注Arista Networks、Cisco、HPE等网络设备巨头的投资者而言，现在需要回答的问题不再是“市场会不会恢复”，而是“在恢复之后，谁会留下，谁会被挤出”。

BARC的模型显示，2024年全球校园网络市场经历了16\%的同比下滑。这个数字本身并不令人意外——2021年至2023年，市场分别增长了20\%、15\%和17\%，远高于疫情前多年的中个位数增速。混合办公带来的超额采购需求在三年内被集中释放，随后必然是库存消化和订单回落。

但真正重要的不是2024年的修正幅度，而是修正之后的增长中枢。BARC将2026年和2027年的增速预期分别下调至8\%和7\%，2021年至2026年的复合年增长率从7\%降至6\%。这意味着，即使市场走出低谷，也不会回到 年的高位，而是回到一个更接近历史均值、但略高于疫情前水平的“新常态”。

这个判断的背后逻辑是：混合办公带来的结构性需求增量已经基本被定价，未来增长将更多依赖正常的设备更新周期和AI推理在边缘侧的网络准备。对于产业决策者而言，这意味着不能再以“市场会回到2023年水平”来制定产能和库存计划；对于投资者而言，这意味着收入增速的预期需要相应下调，估值倍数需要重新锚定。

2. 有线交换仍是最大基本盘，但AI推理带来的边缘升级才刚刚开始

在校园网络的各个子市场中，有线交换始终占据最大份额。BARC预计，2026年有线交换市场增速约为8\%，2027年约为7\%，与整体市场增速基本同步。但报告特别指出，企业级数据中心交换市场的增速远高于校园交换——2026年和2027年分别达到47\%和24\%。

这个差异值得深挖。数据中心交换的高速增长显然与AI训练集群的部署直接相关，但BARC同时提到，企业正在为AI推理准备网络基础设施。AI推理发生在边缘侧，这意味着校园网络中的交换设备需要具备更高的带宽、更低的延迟和更强的可管理性。如果这一判断成立，那么未来几年校园交换的更新周期可能会比历史均值更短，设备单价也可能更高。

对于Arista和Cisco这类同时覆盖数据中心和校园市场的厂商，这意味着一个潜在的交叉销售优势：当客户在数据中心部署了Arista的交换设备后，在校园端也倾向于选择同一品牌，以简化运维和降低兼容性风险。BARC的份额数据也印证了这一点——Arista在校园市场的份额从2020年的0.3\%预计升至2027年的3.4\%，而同期其在数据中心交换市场的份额从11.8\%升至23.1\%。

3. Wi-Fi 7是WLAN市场的明确催化剂，但云管理平台才是真正的竞争壁垒

WLAN市场是BARC报告中增速最为乐观的子领域。预计2026年增长12\%，2027年增长8\%，均高于整体市场。其中，云管理WLAN的增长更为突出，2021年至2026年的复合年增长率被上调至17\%，高于此前预期的15\%。

Wi-Fi 7的商用化无疑是这一增长的核心驱动力。但报告隐含的一个更重要的判断是：云管理平台正在成为WLAN市场的竞争壁垒。BARC估计，云管理WLAN目前约占WLAN总收入的15\%左右，但增速远高于传统WLAN。这意味着，那些能够提供端到端云管理解决方案的厂商——如HPE旗下的Aruba、Cisco的Mer

[... middle omitted ...]

的数据清晰地展示了校园网络市场的竞争格局：Cisco以约50\%的份额占据绝对主导地位，HPE维持在11\%-12\%左右，Arista虽然目前只有低个位数份额，但增长轨迹最为陡峭。

Cisco在2024财年经历了严重的库存修正，校园网络收入从2023年的266亿美元降至204亿美元。但进入2025财年后，订单恢复强劲——第三季度校园网络订单同比增长超过25\%，无线订单增长超过40\%。BARC预计Cisco的校园网络收入将在2026年和2027年继续增长，但增速低于整体市场，这意味着其份额将继续被Arista等竞争对手蚕食。

Arista的增长逻辑非常清晰：从数据中心交换切入校园市场，利用VeloCloud的SD-WAN能力补齐产品组合，加上前Cisco高管Todd Nightingale负责校园产品线的战略推进。BARC预计Arista的校园市场份额将从2020年的0.3\%升至2027年的3.4\%。虽然绝对值仍小，但考虑到校园市场的总规模，这对应着数亿美元的收入增量。

\subsection{[R020] 美国银行：市场低估的不是AI算力需求，而是CPU在Agentic时代的重新定价}
{\small\color{refgray}归类：AI/算力：从模型竞赛到成本信任，供应链紧俏重塑定价权}

美国银行：市场低估的不是AI算力需求，而是CPU在Agentic时代的重新定价

当整个半导体行业的目光都聚焦在GPU算力军备竞赛时，一份来自美国银行的最新研报提出了一个截然不同的核心判断：到2030年，服务器CPU的市场总规模将从当前的350亿美元跃升至1700亿美元以上，实现近5倍的增长。这个数字远高于该机构此前预估的1250亿美元。

这个修正并非简单的乐观情绪。它源自一个正在发生的结构性转变：AI的工作负载正在从“一问一答”的生成式模式，转向“自主规划、调用工具、执行代码”的Agentic模式。在这种新范式下，CPU的角色从“后台管家”升级为“前线总指挥”，其价值也随之被重新定义。

这份报告的核心洞察在于：CPU TAM的扩张，不是对GPU市场的零和博弈，而是整个AI系统复杂化带来的增量。对于投资者和产业决策者而言，理解这个“增量”从何而来，比争论“CPU vs. GPU”更有价值。

1. Agentic AI让CPU从“配角”变为“总指挥”，这才是TAM扩张的真正引擎

传统AI推理中，GPU承担了绝大部分计算，CPU主要负责数据搬运和调度。但在Agentic AI系统中，工作流发生了根本变化：一个AI Agent需要自主理解任务、分解步骤、检索知识库、调用外部API、执行代码，并在多步骤循环中保持状态记忆。

这些任务的特点是：延迟敏感、顺序执行、I/O密集。这正是CPU的专长领域，而非GPU的强项。美国银行的报告将CPU在这一新架构中的角色定义为“中央思维引擎与任务协调器”。

该机构据此将CPU市场拆分为三个相互独立且不重叠的细分领域： 传统/本地/多云CPU：约300亿美元，基本盘，增长平稳。 AI集群计算/头节点CPU：约700亿美元，负责管理GPU集群的调度、数据流和监控。 AI Agentic独立节点CPU：约700亿美元，这是全新的增量，专门为运行Agent逻辑而设计的服务器节点。

这种拆分方式清晰地回答了“为什么是5倍增长”的问题：传统市场（300亿）是存量，AI头节点（700亿）是伴随GPU集群扩张的衍生需求，而Agentic节点（700亿）则是全新的、由应用范式驱动的增量市场。后两者共同构成了1400亿美元的AI CPU市场，这才是TAM跃升的关键。

2. 英特尔的双倍升级与AMD的守擂：谁在Agentic时代拥有结构性优势？

基于新的TAM假设，美国银行对主要CPU厂商的评级和估值进行了大幅调整，其背后的逻辑值得深究。

最引人注目的是对英特尔的双倍升级，从“跑输大盘”直接上调至“买入”，目标价从96美元大幅提升至135美元。该机构认为，市场此前低估了英特尔在解决整个行业面临的先进制程和封装产能瓶颈中的价值。报告特别指出，Cadence的14A节点IP签约和Terafab的合作是支撑其信心的新数据点。到2030年，英特尔在CPU市场的价值份额预计将稳定在24\%左右，虽然份额下降，但来自ASP提升和传统企业市场的收入仍能实现20\%以上的年复合增长。

对于AMD，该机构维持“买入”评级，并将其目标价从500美元上调至560美元，视其为CPU领域的首选股。理由在于AMD在核心数量和高频率上的领先优势，使其在AI头节点和Agentic节点这两个高价值领域都能占据有利位置。预计AMD到2030年将维持约25-26\%的稳定价值份额。

这里的关键洞察在于：英特尔和AMD的份额变化路径不同。英特尔是“份额下降但收入增长”，受益于行业供需紧张和ASP提升；AMD则是“份额稳定”，依靠性能优势在高价值市场扎根。两者都从TAM扩张中获益，但逻辑各异。

3. ARM的崛起与英伟达的“跨界”：格局重塑中的赢家与悬念

ARM架构是这份报告中最激进的变量。该机构预计，到2030年，ARM将

[... middle omitted ...]

MD解决方案的75\%，反映了云厂商自研芯片的成本优势。

英伟达的角色则更为特殊。它既是GPU霸主，也是ARM商用CPU的重要玩家。报告认为，英伟达受益于CPU-GPU-网络的紧密集成，在Agentic CPU时代同样占据有利位置。但一个尚未完全回答的问题是：当英伟达的CPU（如Vera）在自家集群中占据主导时，它与其他CPU供应商（如AMD、英特尔）在头节点和Agentic节点上的竞争边界在哪里？这份报告将英伟达的CPU机会视为“头节点与Agentic节点各占一半”，但这背后的市场份额动态仍需进一步观察。

4. 一个尚未完全回答的关键问题：TAM扩张的假设是否过于乐观？

首先，1700亿美元TAM的核心驱动力是Agentic AI的普及。这个假设成立的前提是，Agentic AI能像当前的生成式AI一样，快速渗透到企业核心业务流程中，并产生可量化的商业价值。如果Agentic AI的应用落地慢于预期，那么700亿美元的Agentic节点市场就将面临显著下修风险。

\subsection{[R021] Citi：光伏市场的真正拐点不在国内价格战，而在海外框架协议的规模与结构}
{\small\color{refgray}归类：工业与材料：供给约束与新兴需求驱动的再定价}

Citi：光伏市场的真正拐点不在国内价格战，而在海外框架协议的规模与结构

过去一年，中国光伏行业在“内卷”与“出清”之间反复拉扯。市场参与者习惯于将注意力锁定在国内产业链价格的每一次波动上——多晶硅又跌了多少，电池片是否跌破现金成本，组件价格何时见底。

但这份Citi在2026年6月10日发布的周度更新，提供了一个截然不同的观察信号。报告的核心判断并非关于价格底部，而是关于需求结构的切换。这份报告最值得看的结论是：中国光伏产品的海外销售正在以框架协议的形式大规模落地，而国内市场的价格信号已经不再是行业景气度的唯一风向标。

对于产业决策者和投资者而言，理解这一结构性变化，远比预测下周多晶硅价格是否再跌1\%更为重要。

截至6月10日当周，n型多晶硅价格下跌1.3\%至每公斤33.3元人民币，太阳能电池片价格下跌4.6\%至每瓦0.31元，组件价格微跌0.3\%至每瓦0.721元。这些数字本身并不令人意外——它们延续了过去两年的下行趋势，且Citi明确指出，当前市场价格已接近可变生产成本。

真正值得关注的不是跌幅本身，而是价格波动的形态。从Citi提供的长期价格走势图来看，多晶硅、硅片、电池片、组件的价格曲线在2024年至2025年间经历了最陡峭的下滑，而进入2026年后，曲线斜率明显放缓。这意味着，即便行业产能过剩的格局尚未根本改变，但价格下行的边际冲击正在减弱。

这背后有两个支撑因素：一是部分高成本产能已经退出市场，二是产业链库存虽然仍处高位，但增速已明显放缓。以多晶硅为例，截至6月4日，生产商库存为29.5万吨，环比下降1.0\%。虽然绝对水平依然很高，但方向性的变化值得留意。

对于投资者来说，这意味着继续押注价格进一步大幅下跌的风险回报比正在恶化。市场可能需要寻找新的定价逻辑。

2. SNEC展会揭示的真正信号：海外框架协议正在重塑竞争格局

Citi报告中最具洞察力的部分，并非价格数据，而是对近期在上海举办的SNEC展会的观察。报告指出，展会上现场签约以吉瓦级的组件框架谅解备忘录为主，隆基、晶科、天合、晶澳等一线组件厂商与主要来自中东、非洲和东南亚的海外买家签署了新的供应协议。

首先，框架协议的性质不同于现货交易。它意味着海外买家对中国光伏产品的需求正在从“按需采购”转向“长期锁定”。这种转变通常发生在下游需求确定性增强、且买家对供应链稳定性要求提高的阶段。中东和非洲市场正在经历能源转型的加速期，这些地区的政府和企业正在为未来数年的光伏项目做规划，而非短期补库。

其次，Citi报告特别提到，多晶硅和硅片的现场签约很少，这些产品的交易主要通过年度长协在场外进行。这进一步印证了产业链上游的定价权正在向能够提供“一体化解决方案”的组件厂商集中。组件厂商通过框架协议锁定了海外需求，进而向上游施压，这一逻辑链条正在变得清晰。

对行业格局的含义是：未来中国光伏企业的竞争力，将越来越取决于其海外渠道的深度和框架协议的规模，而非国内产能的大小。 那些能够在中东、非洲、东南亚等新兴市场建立长期合作关系的企业，将获得比同行更稳定的利润曲线。

3. 政府“反内卷”力度可能弱于预期，行业出清更多依赖市场力量

Citi报告提出了一个容易被市场忽视的判断：2026年政府针对光伏行业的“反内卷”措施，力度可能弱于2025年。原因是中国的PPI和CPI已部分回升，部分受到中东冲突推高能源价格的影响。这意味着政策制定者对于行业亏损的容忍度可能有所提高，不再像2025年那样急迫地推动行政性减产。

如果行政干预减弱，行业产能过剩的消化将更多依赖市场力量——即持续的低价迫使高成本产能退出。这一过程可能比市场预期的更漫长，但也更彻底。那些依赖补贴或政策保护才能存活的企业，将面临真正的生存考验。

对于投资者而言，这意味着在选择标

[... middle omitted ...]

进速度、融资条件、地缘政治风险、以及当地政策变动，都可能影响协议的实际转化率。Citi在德业股份和阳光电源的风险提示中提到了“新兴市场储能需求低于预期”和“海外贸易摩擦加剧”，但并未对这些框架协议的具体执行条款和违约概率进行量化分析。

这是读者在消化这份报告时需要保持审慎的地方。海外框架协议的规模确实令人振奋，但它们是否能够转化为实际的收入和利润，取决于一系列超出光伏行业本身的因素——包括国际关系、汇率波动、当地电网基础设施的配套能力等。

Citi报告在展会观察中还提及了一个值得关注的细节：由协鑫和晶澳分别牵头，两个空间光伏产业联盟已正式启动，领先制造商正在开发用于卫星的轻质叠层电池，以开拓全新的应用市场。

虽然这一领域目前规模极小，对主流光伏企业的短期业绩几乎没有影响，但它代表了一个重要的方向性变化。光伏技术的应用边界正在从地面电站、屋顶分布式，向航空航天等高端领域延伸。如果空间光伏能够实现商业化突破，将为中国光伏产业打开一个技术溢价远高于地面电站的全新市场。

\subsection{[R022] Citi：AI供应链的紧俏不是短期噪音，而是新一轮定价权转移的信号}
{\small\color{refgray}归类：AI/算力：从模型竞赛到成本信任，供应链紧俏重塑定价权}

Citi：AI供应链的紧俏不是短期噪音，而是新一轮定价权转移的信号

市场对AI硬件的讨论，近期似乎出现了一些杂音。800V高压直流（HDVC）的采用被质疑进展缓慢，共封装光学（CPO）的落地时间表被推后，VR NVL72的内存配置低于预期。这些声音让部分投资者开始怀疑，AI供应链的景气周期是否已经见顶。

但Citi刚刚发布的这份台湾科技月度追踪报告，给出了一个与市场情绪截然不同的结论。报告的数据和实地调研指向一个清晰的判断：供应链的紧张不仅没有缓解，反而正在从产能层面蔓延到定价层面，而这一轮结构性变化，真正受益的或许不是市场最关注的那些名字。

这份报告最值得关注的，不是某个公司的月度营收超预期，而是一个正在发生的系统级变化——整个台湾科技供应链，从先进制程晶圆代工到封装测试，再到基板与铜箔基板（CCL），正在经历一轮由成本推动、由产能紧缺支撑的全面涨价周期。这不仅是短期的供需错配，更是AI基础设施建设从“抢产能”进入“抢议价权”阶段的标志。

1. 先进制程与封测的紧俏不是短期现象，而是AI芯片迭代加速的结构性结果

Citi报告中最硬的信号来自台积电。5月营收4170亿新台币，同比增长30\%，前5个月累计增长同样达到30\%。这个数字本身并不令人意外，但报告明确指出，随着下半年N3制程的AI芯片和N2制程的智能手机SoC开始放量，台积电有极大可能再次上调全年营收增长目标。这意味着，市场对AI芯片需求的测算可能仍然偏保守。

真正值得关注的是，Citi提到了一个增量因素：用于“代理型AI”（agentic AI）的CPU，将在明年为台积电提供额外的增长动力。这暗示AI芯片的需求结构正在从单一的GPU训练，向更多元的推理与边缘计算场景扩展。如果这一判断成立，那么先进制程的产能紧张就不是一个季度或两个季度的问题，而是一个贯穿未来2-3年的主题。

封测端的信号同样清晰。ASEH和KYEC的订单流在2H26将明显改善，驱动力来自LEAP（先进封装技术）和新的AI芯片测试业务。Citi特别提到，测试接口领域的MPI、Hon Precision和Chroma，二季度营收有望超出市场预期，原因是来自GPU和AI ASIC客户的订单强劲，同时规格还在持续升级。这揭示了一个容易被忽视的事实：AI芯片的迭代不仅拉动晶圆制造，更在同步拉高测试环节的技术门槛和单价。

2. CCL的涨价不是成本转嫁，而是供应链议价权正在从下游向上游转移

在众多子行业中，铜箔基板（CCL）板块的表现最值得深究。报告指出，CCL厂商的4月和5月营收连续超预期，驱动力来自“持续有效的逐客户涨价”。这种表述意味着，涨价不是一次性的，而是分批次、分客户逐步推进的，且执行效果良好。

与此形成对比的是，基板（substrate）板块虽然也实现了环比增长，但涨价的幅度和速度低于多头预期。CCL和基板的差异，本质上反映的是不同环节的议价能力差异。 CCL行业集中度更高，且下游客户（PCB厂商）相对分散，因此涨价更容易落地。而基板行业竞争格局更复杂，客户议价能力更强，导致涨价传导更慢。

这对投资者的含义是：在AI供应链中，并非所有环节都能平等受益。那些行业集中度高、技术壁垒明确、客户转换成本高的上游材料环节，可能比中游的组装或模组环节拥有更强的定价权。Citi的观察暗示，CCL的涨价可能只是一个开始，如果成本压力持续，PCB乃至更下游的环节都可能被迫跟进。

3. AI服务器ODM的分化：鸿海领先，但纬颖的月营收暗示AWS的AI ASIC进度可能低于预期

ODM/OEM板块的5月营收环比微增，但内部出现了显著分化。鸿海5月营收8594亿新台币，同比增长40\%，环比增长3\%，领先Citi对二季度环比增长15\%的预估。广达和纬创也略好于预期。但纬颖的5月营收表

[... middle omitted ...]

判断，但同时也指出，VR NVL72的出货将在3Q26才开始逐步爬坡，二季度出货量预估仅为1.3-1.5万台。

这意味着，AI服务器的增长驱动力，在2026年下半年之前，可能仍然高度依赖于NVIDIA的GPU平台，而非ASIC。任何对ASIC替代GPU的过度乐观预期，都需要对照纬颖的实际出货数据重新校准。

4. 消费电子需求疲软但未恶化，真正的风险在于成本上涨对终端价格的传导

非AI领域的表现，Citi的描述是“比担心的要好”。联发科、瑞昱和联咏的营收虽然受消费需求低迷影响，但凭借产品组合优化，表现相对稳健。Citi甚至预期联发科可能在2H26涨价，以应对成本上升。

但NB（笔记本电脑）板块的观察更值得警惕。5月NB出货量环比增长2\%，原因是客户担心元器件成本持续上涨而提前拉货。Citi预估二季度NB出货量环比增长2\%，同比下降10\%。报告明确写道，虽然市场对AI PC渗透率加速有期待，但“整个生态系统的准备仍需时间”，且“零售价格上涨正在压制消费者购买力”。

\subsection{[R023] JPM：中国汽车市场真正的拐点不是电动化渗透率，而是燃油车库存的定价崩塌}
{\small\color{refgray}归类：电动车与汽车：区域利润结构重构与燃油车定价崩塌}

JPM：中国汽车市场真正的拐点不是电动化渗透率，而是燃油车库存的定价崩塌

这份来自JPM日本股票研究团队的五月中国汽车行业报告，表面上看是在更新月度销量数据，但真正值得产业决策者注意的，不是新能源渗透率突破63\%这个数字本身，而是报告中一个被轻描淡写、却可能改变整个市场定价逻辑的信号——零里程二手燃油车数量正在上升。

这意味着什么？意味着燃油车的新车价格和销量之间，正在形成一个自我强化的下行螺旋。这个夏季，燃油车市场将面临的不只是季节性低迷，而是一场由库存结构引发的定价机制重塑。电动化的故事我们已经听了太多次，但燃油车资产贬值速度的加速，才是当前最被低估的结构性变量。

这份报告的价值在于，它用数据揭示了一个正在发生但尚未被充分讨论的转折点：消费者在新能源车比同级别燃油车更贵的情况下，仍然在加速转向新能源。这不是价格驱动的替代，而是认知驱动的替代。当燃油车的新车折价和二手残值同步下降，整个燃油车生态的盈利模型都需要重新审视。

以下是我们从这份JPM报告中提炼出的五个核心洞察，以及一个尚未被充分回答的关键问题。

1. 新能源渗透率突破63\%的真正含义：消费者已不再用价格做选择

五月中国乘用车零售销量152.8万辆，同比下降22\%，连续第八个月同比下滑。但在这个总量萎缩的背景下，结构性的分化才是关键。燃油车零售57.9万辆，同比暴跌38\%；新能源车零售84.9万辆，仅下降7\%。新能源渗透率从四月的61\%进一步攀升至63\%。

更值得关注的是价格信号。JPM指出，在10万元人民币（约合230万日元）以下的入门级市场，主流新能源车型的价格已经高于同级别燃油车。换言之，消费者选择新能源车，不是因为它们更便宜，而是因为它们在智能化体验、使用成本和使用场景上提供了燃油车无法替代的价值。

这对传统车企意味着什么？过去几年，传统车企和投资者普遍认为，新能源车的优势主要来自补贴和政策驱动，一旦补贴退坡或价格优势消失，消费者会回归燃油车。但五月的数据否定了这个假设。即便新能源车涨价、燃油车打折，消费者的迁移仍在加速。这不是一个价格弹性问题，而是一个产品价值再定义的问题。

对于丰田、本田等日系厂商而言，这个信号尤其严峻。五月数据显示，丰田中国零售10.2万辆，同比下降29\%；本田2.8万辆，同比下降49\%。它们不是没有降价，而是降价本身已经无法阻止市场份额的流失。

JPM在报告中做出了一个具体的季节性判断：燃油车销量将在今年夏季进一步下滑。这个判断本身并不意外，但背后的支撑逻辑值得深究。

报告提到，当前零里程二手燃油车的数量正在上升。所谓零里程二手车，指的是已经上牌但从未被使用过的车辆。这类车辆的出现通常有两个原因：一是经销商为了完成月度销售任务而自行上牌，二是部分企业或租赁公司批量采购后未实际投入使用。无论哪种原因，这些车辆一旦进入二手市场，就会对新车价格形成直接压制。

更关键的是，这类库存的上升会引发连锁反应。当市场上出现大量零里程二手车，新车经销商不得不进一步降价以维持竞争力；新车降价又会压低正常二手车的残值；残值下降反过来削弱消费者购买燃油车的意愿，进一步压缩新车销量。这是一个典型的负反馈循环。

JPM判断，这个夏季燃油车面临的销售环境将比往年更加严峻。对于持有大量燃油车库存的经销商和依赖燃油车利润的传统车企，这可能是过去几年中最难熬的一个夏天。

与国内市场的疲软形成鲜明对比，中国乘用车出口在五月达到80.9万辆，同比增长73\%，连续第三个月创历史新高。其中新能源车出口43.5万辆，同比增长110\%；燃油车出口37.4万辆，同比增长42\%。

这里有一个容易被忽略的细节：JPM特别指出，比亚迪等主流中国品牌的产能正在接近极限，交付周期在中国国内和海外市场都在延长。这意味着，出口需求的强劲不仅是一个量的故

[... middle omitted ...]

系车企的困境：不是电动化转型慢，而是燃油车资产贬值速度超预期

JPM对日本车企在中国市场的表现着墨不多，但数据本身已经说明问题。丰田、本田、日产三家日系车企五月在中国市场的销量均出现同比两位数下滑，其中本田的跌幅接近50\%。

市场通常将日系车企的困境归因于电动化转型缓慢。这个判断没错，但不完整。JPM的报告揭示了另一个维度：即便日系车企加速电动化，它们在中国市场面临的真正挑战可能是燃油车资产的快速贬值。

报告提到，过去几个月燃油车的二手残值已经下降，相对于新能源车的优势正在消失。对于日系车企而言，它们在中国市场的核心资产——庞大的燃油车产能、经销商网络、零部件供应链——都建立在燃油车的高利润和高残值基础之上。一旦这个基础动摇，调整的代价将远高于推出几款电动车型。

更值得关注的是，日系车企在中国市场的困境可能只是一个开始。如果零里程二手燃油车库存继续上升，新车价格和二手残值的下行螺旋会加速，届时不仅是日系车企，所有依赖燃油车利润的传统车企都将面临类似的压力测试。

\subsection{[R024] 摩根斯坦利：市场低估了供给侧的“结构性断裂”，而非仅仅需求回暖}
{\small\color{refgray}归类：能源与大宗：需求破坏对冲供给冲击，铝与原油的结构性变局 / 工业与材料：供给约束与新兴需求驱动的再定价}

摩根斯坦利：市场低估了供给侧的“结构性断裂”，而非仅仅需求回暖

这份摩根斯坦利2026年6月发布的“亚洲暑期学校：中国材料”报告，传递了一个与市场普遍叙事不同的信号。多数投资者仍在纠结于中国房地产何时触底、基建投资能否加码，但该机构认为，真正值得关注的并非需求端的线性复苏，而是供给端正在发生的、不可逆的结构性变化。

报告的核心判断是：中国材料行业正经历一场由供给约束与新兴需求（储能系统ESS与人工智能AI）共同驱动的再定价。摩根斯坦利明确表示，在这一环境下，他们更偏好铝、铜、锂、铀、金和玻璃纤维等领域的权益资产。这意味着，投资者需要将分析框架从“需求预测”转向“供给弹性评估”。

这份报告的独特价值在于，它不仅给出了自上而下的行业观点，还提供了详尽的消费指数模型，将宏观叙事拆解为可追踪的微观指标。以下是我们从这份报告中提炼出的几个关键洞察。

市场对大宗商品的分析往往过度聚焦于需求端，尤其是中国房地产的拖累。但摩根斯坦利在报告中反复强调的，是供给侧的长期紧缩。这并非简单的矿山停产或冶炼厂检修，而是一种更深层次的、由资本开支不足、ESG合规成本上升以及地缘政治风险共同塑造的供给刚性。

以铝为例，中国电解铝产能已经逼近4500万吨的政策天花板，新增产能极为有限。同时，海外能源成本的上升也抑制了非中国地区的复产意愿。摩根斯坦利对2026年铝价的预测为3386美元/吨，较市场共识高出11\%，对2027年的预测更是高出13\%。这种显著的偏离，恰恰说明市场尚未充分定价供给侧的“硬约束”。

对于投资者而言，这意味着那些拥有稳定、低成本产能且符合环保标准的龙头企业，其定价权将比历史周期更强。传统的“需求下行-价格下跌”逻辑链，可能被供给侧的刚性所打断。

2. 铜的“双引擎”正在切换：房地产拖累减弱，电网与AI成为新锚点

铜是市场分歧最大的品种之一。摩根斯坦利对2026年铜价的预测（12292美元/吨）与市场共识基本持平，但对2027年的预测高出2\%。这看似温和的差异背后，是需求结构的根本性变化。

报告中的中国铜消费指数（CCI）拆解显示，电力（占比47\%）已经成为绝对主导力量。更重要的是，报告明确提到了“来自ESS与AI的改善需求”。虽然报告没有详细展开AI如何具体拉动铜消费，但我们可以合理推断：数据中心、算力基础设施的扩张，以及配套的电力传输与冷却系统，都是铜的增量消费场景。与此同时，房地产在铜消费中的占比已降至9\%，其边际影响正在快速衰减。

这意味着，铜价的下一个驱动因素不再是中国的“旧经济”复苏，而是全球能源转型与科技基础设施建设的共振。投资者需要关注的宏观指标，应从“中国房屋新开工面积”转向“中国电网投资计划”和“全球数据中心资本开支”。

3. 锂的“预期差”最大：摩根斯坦利对2026年锂价的预测远超共识

在报告提供的所有价格预测中，锂的预期差最为突出。摩根斯坦利对2026年碳酸锂价格的预测为22840美元/吨，较市场共识的19559美元/吨高出17\%。这是所有品种中偏离度最大的之一。

这似乎与市场的主流叙事相悖。目前，市场普遍认为锂行业正面临严重的供给过剩，价格中枢将持续下移。但摩根斯坦利的观点暗示，供给侧的出清速度可能快于预期，或者需求端（尤其是储能）的增长正在超出模型的假设。报告中“改善需求来自ESS与AI”的判断，对锂而言同样适用——储能电池是锂的绝对消费主力。

这个预期差值得深度关注。如果供应端的减产、停产计划执行到位，而储能装机量持续超预期，那么锂价可能会在当前市场极度悲观的情绪中迎来反弹。但报告没有详细解释其预测的假设条件，这恰恰是我们在完整报告中需要重点验证的部分。

摩根斯坦利对中国钢铁和水泥行业的态度显然更为谨慎。报告中的钢铁消费指数显示，2026年钢铁表观消费量预计同比下降

[... middle omitted ...]

统性机会。

这提醒我们，即便供给侧有约束，但如果需求端萎缩的斜率更陡，企业的盈利仍将承压。对于钢铁和水泥，更值得关注的是行业内部的结构性整合，而非整体性的贝塔行情。

5. 报告尚未完全回答的关键问题：AI与ESS的“二阶效应”如何量化？

这是整份报告最具启发性、同时也最需要深入探讨的部分。摩根斯坦利提出了“供给约束+ESS与AI需求改善”这一框架，但对于AI对材料需求的拉动机制，报告并未给出详细的量化模型。

例如，AI对铜的需求，是体现在数据中心内部的电力布线，还是更上游的电网扩容？对锂的需求，是体现在大型储能电站，还是AI终端的便携式电源？这些不同场景下的单位材料消耗强度差异巨大。此外，铀作为核能发电的原料，其需求逻辑也高度依赖于AI带来的电力需求增长预期。报告将铀列入偏好清单，但未详细展开其与AI的关联路径。

这些“二阶效应”的量化，正是专业投资者可以比市场共识走得更远的地方。理解这些传导机制，才能对摩根斯坦利给出的价格预测做出独立判断，而不是盲目跟随。

\subsection{[R025] 摩根斯坦利：中国聚乙烯净进口量才是美国合同价格真正的定价锚}
{\small\color{refgray}归类：工业与材料：供给约束与新兴需求驱动的再定价}

市场参与者习惯将目光聚焦于中东霍尔木兹海峡的供应中断，并据此线性外推美国聚乙烯价格的上涨路径。摩根斯坦利最新发布的研报却提供了一个截然不同的分析框架：真正决定美国合同价格走向的关键变量，并不在中东，而在中国。这份报告的核心判断是，中国聚乙烯净进口量的变化幅度与可持续性，将成为未来数月美国合同价格走势最核心的KPI，其影响力甚至超过了供应中断事件本身。

4月份的数据已经验证了这一逻辑的威力。当霍尔木兹海峡关闭导致全球供应减少约15\%时，中国净进口量的急剧收缩，从冲突前每月约110万吨骤降至4月的约18万吨，相当于向除中国以外的全球市场释放了约9\%的额外供应。这一行为几乎完全对冲了中东供应中断带来的紧张局面，使得4月美国出口市场并未如市场普遍担忧的那样出现极端紧缺。市场低估的，从来不是地缘政治事件本身，而是中国作为全球最大聚乙烯买方，其库存行为与贸易流调整对全球供需平衡的再定价能力。

报告的核心发现之一，是中国在冲突爆发前积累了远超市场认知的过剩库存。摩根斯坦利通过咨询行业专家推断，在伊朗冲突爆发前的6至12个月内，中国贸易商可能积累了约100万吨的额外库存，将中国的过剩库存总量推高至约230万至350万吨，相当于约30天的国内需求量。这一判断直接解释了4月中国为何有能力将净进口依赖度从2025年约31\%的水平骤降至约5\%。

这一行为意味着，中国并非被动接受供应中断，而是主动利用其庞大的库存作为缓冲，不仅减少了进口（尤其是从美国的进口减少了约40\%，约50万吨），还增加了出口（约40万吨，主要流向东南亚以填补中东出口缺口）。这种主动的贸易流调整，使得中国在4月从一个巨大的净进口国，几乎转变为自给自足状态。对于全球聚乙烯市场而言，中国这一角色的切换，相当于一个巨大的需求缓冲器被激活，其短期内的平抑效果远超任何单一供应端的扰动。

报告的分析表明，中国这一行为可能只是暂时的。情景分析显示，上述过剩库存可能在6月或7月消耗殆尽。摩根斯坦利给出的关键假设是：一旦中国恢复至冲突前的净进口水平，将为美国出口带来每月约25万吨的增量需求，这相当于美国年供应量的约12\%。这是一个巨大的数字。

更值得关注的是，中国国内聚烯烃库存数据已经开始验证这一去化趋势。来自中石化和中石油的库存数据显示，目前库存已较近期高点下降约19\%，回落至与2024年和2025年相当的水平。这一去库存过程在季节性上也是异常的——2024年和2025年4月的库存高点通常低于2月水平，而2026年4月的库存却逆势走高，随后在5月大幅下降。这意味着，中国库存的“超额部分”正在迅速被消化，市场即将面临一个关键节点：当库存降至正常水平后，中国是否会如报告所预期的那样，重新大规模回归国际市场采购。

报告提供了另一个关键背景：美国聚乙烯的有效开工率在5月已接近极值，约98.8\%。这意味着，美国生产商几乎没有多余的闲置产能来应对突然增加的出口需求。如果中国在6月或7月恢复进口，美国必须从现有产量中挤出额外的出口量，唯一的途径就是消耗国内库存。

摩根斯坦利测算，25万吨的月均增量出口将使美国库存水平在8月前下降约8天。考虑到当前美国生产商库存已从4月的42.5天上升至5月的约46天（因出口不及预期），这一库存去化将迅速扭转当前的宽松局面。报告中的情景分析清晰地展示了两种路径：在增量出口推动下，合同价格能够维持甚至走高；而在均值回归加速的情景下，价格将更快回落。当前市场对6月合约的定价（平盘）可能低估了库存去化一旦启动后的价格弹性。

4. 报告尚未完全回答的关键问题：中国库存数据的透明度和行为逻辑

尽管报告提供了极具洞察力的分析框架，但一个核心问题并未完全展开：中国库存数据的透明度和可信度。报告引用的中石化与中石油库存数据，能否完全代表中国整体贸易商和下游用户的

[... middle omitted ...]

观察框架：从关注“中东供应”转向追踪“中国净进口”

基于这份报告，投资者需要重新校准自己的观察框架。过去几个月，市场过度聚焦于霍尔木兹海峡的重新开放时间表，而忽略了需求端的结构性变化。摩根斯坦利的研究表明，中国净进口数据才是更具前瞻性的价格领先指标。

一个实用的观察框架是：首先，密切追踪中国海关每月发布的聚乙烯净进口数据，尤其是6月和7月的数据，这将直接验证库存去化假设是否成立。其次，关注美国生产商的库存天数变化，如果库存从当前的46天开始持续下降，将是出口需求启动的先验信号。最后，关注中国国内聚烯烃库存的周度数据，如果库存去化速度放缓甚至反弹，可能意味着中国仍在通过其他渠道补充库存，而非回归正常的进口节奏。

这份报告的价值，在于它提供了一个从“买方视角”理解全球大宗商品定价的新范式。在全球供应链日益碎片化的背景下，拥有庞大库存和贸易调节能力的买方，其行为对价格的影响力可能不亚于任何供应端的冲击。理解中国作为“定价锚”的角色，远比猜测中东的炮火何时停歇更为重要。

\subsection{[R026] BARC：通胀数据揭示的真正风险不是物价上涨，而是定价权断裂}
{\small\color{refgray}归类：宏观与利率：通胀结构分裂与政策信号再校准}

BARC：通胀数据揭示的真正风险不是物价上涨，而是定价权断裂

中国5月通胀数据已经发布。CPI同比1.2\%，PPI同比3.9\%。从数字看，物价似乎在温和回升。但真正值得关注的，不是这两个数字本身，而是它们背后的结构性断裂：上游成本正在加速传导，但下游企业几乎完全丧失了转嫁能力。PPI上游原材料同比上涨9.2\%，采掘业同比上涨15.8\%，而下游消费品PPI却仍在收缩，同比负增长0.8\%。这组数据描绘的不是一个通胀周期，而是一个利润分配极端失衡的图景。

这份BARC研报的核心价值，不在于它预测了通胀数字，而在于它揭示了当下中国经济中最关键的价格传导机制正在失效。当上游成本无法向下游有效传递时，制造业企业的利润率将面临持续挤压，而终端需求的疲弱又会反过来抑制就业和收入预期。这不是一个短期波动，而是一个需要重新评估资产定价的结构性信号。

BARC的报告指出，5月核心CPI从1.2\%回落至1.1\%，服务CPI从0.9\%降至0.8\%。在能源价格推高整体通胀的同时，扣除能源和食品的核心通胀却在走弱。这意味着，所谓的“通胀”更多是输入性的成本推动，而非需求拉动的内生复苏。对于投资者和产业决策者而言，真正需要关注的不是CPI和PPI的绝对值，而是这两者之间的剪刀差所反映的定价权分布。

这份研报最清晰的信号来自PPI的分项数据。5月PPI同比增速从4月的2.8\%跳升至3.9\%，连续第二个月加速。但这一增长高度集中于上游。采掘业PPI同比上涨15.8\%，原材料PPI上涨9.2\%，而制造业PPI仅增长2.3\%。这意味着，利润正在向上游资源品环节集中，而大量中下游制造企业面临成本上升与售价低迷的双重挤压。

BARC特别指出，有色金属矿采选和加工业是5月PPI上涨的最大拉动项，贡献了2.56个百分点。能源和化工品是另一大驱动因素，油气开采、炼油和化工产品合计贡献了近2个百分点。这两大板块的强势，与AI和绿色技术需求带动的铜价上涨，以及中东地缘冲突引发的能源供给担忧直接相关。

但问题的另一面是，消费品PPI仍在收缩，5月为负0.8\%。虽然较4月的负1.0\%略有收窄，但连续多月处于通缩区间。这意味着，下游企业——从家电、纺织到日常消费品——根本无法将上游涨价的压力向下游消费者传递。在需求疲弱的环境下，任何涨价行为都可能导致市场份额的流失。

对于投资者而言，这意味着需要重新审视产业链上不同环节的定价能力。上游资源品企业可能在当前环境下受益，但中下游制造企业的利润率面临系统性压力。那些缺乏品牌溢价、产品差异化不足的企业，将首当其冲。

CPI整体数字维持在1.2\%，但结构分化严重。能源价格是唯一的支撑项，汽油价格同比上涨23.5\%，拉动CPI约0.66个百分点。而食品价格连续第二个月下跌，同比负增长1.7\%，其中猪肉价格跌幅加深。扣除食品和能源的核心CPI从1.2\%回落至1.1\%，低于1-2月冲突前的平均水平1.3\%。

BARC的研报特别强调，服务CPI从0.9\%降至0.8\%，其中旅游相关服务价格涨幅收窄了0.9个百分点至2.8\%。住房租金持续下降，同比负增长0.6\%，为2023年1月以来最快降速。更值得关注的是，中国指数研究院的数据显示，主要城市的房租实际同比跌幅高达3.2\%，远高于官方统计的0.6\%。

汽车价格继续下跌，同比降1.1\%，环比降0.4\%。这表明，在激烈的市场竞争和产能过剩的背景下，车企仍然无法通过涨价来转嫁成本。家电和服装价格虽然小幅上涨，但合计对CPI的拉动仅约0.12个百分点，难以对冲其他项目的下行压力。

这些数据合在一起意味着：当前的通胀格局是“成本推动型”而非“需求拉动型”。能源价格上涨推高了整体数字，但居民消费的内生动力依然疲弱。对于政策制定者而言，这增加了决策的复杂性——如果因为整体CPI回升而收

[... middle omitted ...]

跌0.4\%。在新能源车渗透率不断提升、传统车企面临转型压力的背景下，价格竞争只会更加激烈。汽车作为大件消费品，其价格走势直接反映了居民消费意愿和购买力。

这两个指标背后，是同一个逻辑：当居民对未来收入预期不确定时，他们会推迟大宗消费、压缩居住成本。这种行为的累积效应，就是终端需求的持续疲弱。而终端需求疲弱，又是下游企业无法转嫁成本的根本原因。

4. 报告尚未完全回答的关键问题：这轮价格传导断裂会持续多久

BARC的研报提供了清晰的数据和判断，但有一个关键问题并没有完全展开：这种上游涨价、下游通缩的格局会持续多久？或者说，在什么条件下，价格传导机制才能重新恢复？

从历史经验看，中国PPI和CPI之间的剪刀差通常会经历几个月的滞后才逐步收敛。但当前的环境与以往不同。以往的价格传导断裂，往往是因为需求端暂时性疲弱，一旦政策刺激到位，传导机制就会恢复。但这一次，需求疲弱的背后是结构性因素：房地产行业调整仍在持续，居民杠杆率已处于高位，地方财政压力限制了基建投资的空间。

\subsection{[R027] NOM：人民币中间价模型正在揭示一个被低估的政策信号}
{\small\color{refgray}归类：宏观与利率：通胀结构分裂与政策信号再校准 / 区域市场：欧元弱势与人民币汇率信号切换}

人民币汇率中间价的每日设定，长期以来被市场视为中国央行汇率管理意图的“温度计”。但多数观察者只盯着最终公布的数值本身，而忽略了模型内部各因子的结构性变化。NOM最新发布的USD/CNY定盘价模型报告，提供了一个更锐利的观察视角：模型隐含的“逆周期因子”调整力度正在发生质变，这背后可能不是简单的维稳，而是对人民币定价逻辑的一次系统性再校准。

这份报告的核心价值不在于给出一个预测点位，而在于它拆解了中间价形成的“黑箱”。当模型投影值与实际公布值之间的偏差持续扩大时，市场真正应该关注的不是汇率会贬到多少，而是政策制定者正在用什么样的新框架来管理预期。本文将从NOM模型的三个关键维度出发，解读这一信号对投资者意味着什么。

NOM模型显示，最新的USD/CNY定盘价模型投影值为6.7816，较前一次的6.8130下降了314个基点。这个幅度在近期并不常见。但更值得关注的是，模型投影的下降并非由单一货币驱动，而是来自一个多元化的因子组合。

根据报告中的贡献度分析，欧元和澳元分别贡献了15.6和8.5个基点的正向变动，而韩元则反向贡献了13.5个基点。这意味着，人民币中间价的模型投影变化，本质上是全球外汇市场在特定时间窗口内对主要货币相对价值重估的映射。欧元和澳元的走强，部分抵消了韩元等亚洲货币带来的贬值压力。

这里的关键洞察是：中间价模型并非被动反应，而是一个动态加权系统。当欧元和澳元的权重上升时，美元指数的波动对人民币的影响就会被部分对冲。NOM模型揭示的，正是一个正在从“盯住美元”转向“盯住一篮子货币”的精细化操作过程。对于投资者而言，这意味着不能再简单用美元指数的涨跌来预判次日中间价，而需要跟踪一篮子货币的日内变动。

模型投影值为6.7816，但计入逆周期因子后的投影值为6.7957。两者之间173个基点的差异，恰恰是报告中最有信息量的部分。逆周期因子作为央行调节中间价的可选工具，其使用力度和方向，往往比模型本身的输出更能反映政策层的真实态度。

NOM报告没有直接给出逆周期因子的具体数值，而是通过“模型投影与含逆周期因子投影”的对比，让市场自行推导。173个基点的正向调整，意味着政策层面在模型本身指向更强的人民币时，选择了适度“压住”升值幅度。这透露出的信号是：决策者并不希望人民币在短期内出现单边快速升值，即便基本面因素支持这一方向。

这种“反向对冲”的操作逻辑，与 年期间通过逆周期因子抑制贬值预期的做法，在方法论上如出一辙，但方向截然相反。它表明政策工具箱仍然有效，且使用逻辑高度一致——核心目标是“平滑”，而非“引导”。对于外汇交易者而言，这意味着中间价的实际公布值可能持续低于模型隐含的“无干预”水平，从而为即期汇率提供一个“软天花板”。

报告中的模型误差图（Fig. 2）提供了一个历史视角。从2025年初到2026年中，模型误差从-1800个基点逐步收窄，并在2025年7月附近趋近于零，随后又出现正向波动。这种误差的“归零”过程，通常意味着模型对现实定价的解释力在提升，但误差的再次扩大，则暗示有新的非模型变量正在影响实际定盘价。

值得警惕的是，2026年3月19日前后，USD/CNY定盘价出现了约150个基点的单日跳升。NOM报告没有对此事件做归因分析，但结合模型误差的历史模式，这种“异常点”往往对应重大政策表态或外部冲击。当前模型误差处于相对低位，但并不意味着未来不会再次出现类似2025年初的大幅偏离。

这里的核心问题是：误差的收窄是否可持续？如果逆周期因子的使用频率增加，模型误差可能被人为压低，但这反而增加了未来出现“补跌”或“补涨”的风险。投资者不能因为近期模型预测的准确性提高，就放松对尾部事件的警惕。NOM报告没有给出答案，但它的数据框架为每个读者提供了自行推演的基础。

4. 报告尚

[... middle omitted ...]

没有提供任何关于政策层内部决策规则的线索，而这也正是机构投资者最需要自行研判的部分。它可能取决于中美关系、出口数据、资本流动压力等一系列外部变量。对于读者而言，与其猜测一个具体的点位，不如建立一个跟踪框架：持续监测模型投影与含逆周期因子投影的差值，当这个差值出现趋势性扩大时，就意味着政策逻辑可能正在发生转变。

基于NOM报告的分析框架，投资者可以构建一个更系统的观察方法。核心不再是预测明天的人民币中间价是多少，而是理解模型内部各因子的相对权重变化。

首先，关注前四大贡献货币的每日变动。欧元、澳元、韩元、卢布的权重并非固定不变，它们的相对贡献度变化，反映了全球资本流动的底层逻辑。其次，跟踪模型误差的滚动标准差。如果误差在短期内迅速扩大，通常意味着有未被模型捕捉的“政策冲击”或“事件驱动”正在发生。最后，将逆周期因子的调整幅度与同期中国国债收益率、A股外资流入等数据交叉验证。当逆周期因子与资本流动方向一致时，政策的可信度较高；当两者出现背离时，则需警惕市场预期的分化。

\subsection{[R028] NOM：人民币中间价模型正在释放一个被市场低估的政策信号}
{\small\color{refgray}归类：宏观与利率：通胀结构分裂与政策信号再校准 / 区域市场：欧元弱势与人民币汇率信号切换}

市场对人民币汇率的讨论，往往集中在两个维度：一是美元指数的强弱，二是中国央行的政策意图。但NOM这份发布于2026年6月12日的USD/CNY定盘价模型报告，提供了一个更精细、更值得关注的视角——中间价定价机制本身正在发生结构性变化，而市场对此的定价可能远远不够。

这份报告的模型投影显示，在没有逆周期因子的情况下，USD/CNY中间价模型预测值为6.7607，较前一日的6.8150大幅下调543个基点。即使计入逆周期因子，预测值也降至6.7759，较前一日定盘价低391个基点。这不是一次简单的技术调整。它意味着，在当前的汇率形成机制下，人民币中间价的内在压力方向已经发生了根本性扭转。

真正重要的不是这个数字本身，而是数字背后的逻辑：NOM的模型正在捕捉到一种从“贬值压力”向“升值压力”的切换信号。这种切换，如果持续，将深刻影响所有持有人民币资产、或对中国宏观经济有敞口的投资者的定价框架。

1. 模型投影的543个基点下调，不是噪音，是趋势拐点的信号

理解这份报告的价值，首先要理解NOMUSD/CNY定盘价模型的工作原理。该模型并非预测市场交易价格，而是基于一篮子货币的隔夜变动、以及其他技术性因子，来推算央行可能设定的每日中间价。其核心价值在于，它剥离了市场情绪和短期投机因素，呈现的是“如果央行完全遵循规则，中间价应该在哪里”。

6月12日的模型投影显示，无逆周期因子调整的中间价应为6.7607，较前一日官方定盘价低543个基点。这是一个巨大的单日变动。回顾报告中的历史模型误差图（Fig. 2），在2025年的大部分时间里，模型误差持续为负，意味着实际中间价持续高于模型预测——央行在用逆周期因子对冲贬值压力。但从2025年底开始，误差开始收窄并转正，到2026年4月，误差已转为正值约400个基点。

这个趋势变化意味着什么？它意味着，央行此前需要持续“托住”中间价、防止其过快贬值的工作，已经基本完成。现在，模型本身正在推动中间价向升值方向移动。如果这一趋势延续，央行甚至可能面临“是否需要允许中间价更快升值”的新问题。

对于投资者而言，这是一个需要重新校准汇率预期的重要信号。过去两年，市场习惯性地将人民币视为有贬值压力的货币，并据此定价所有中国相关资产。如果中间价机制本身正在转向升值方向，那么所有基于“人民币贬值”假设的定价逻辑都需要被重新审视。

2. 逆周期因子的角色正在从“防御”转向“调节”，其信号意义截然不同

报告特别区分了“有逆周期因子”和“无逆周期因子”两个模型投影。有逆周期因子的预测值为6.7759，虽然也较前一日下调391个基点，但幅度小于无因子模型。这意味着，央行仍然在使用逆周期因子进行一定程度的“缓冲”，但方向与之前完全不同。

在2025年的大部分时间里，逆周期因子被用于“对抗”贬值压力——央行通过该因子将中间价设定在比模型预测更低的水平（即更偏向人民币升值），以稳定市场预期。而现在，模型本身已经转向升值，逆周期因子的作用变成了“减缓”升值的速度，而不是“阻止”贬值。

这是一个根本性的角色转换。当逆周期因子从“防御性工具”变为“调节性工具”时，它释放的信号是：央行对人民币汇率的评估已经不再紧迫地需要干预，而是可以更从容地管理节奏。这往往是一轮汇率趋势变化的早期信号。

从历史经验看，逆周期因子的大幅使用往往伴随着单边预期的强化，而其使用幅度的收窄乃至方向反转，则预示着市场预期正在走向均衡。当前的情况，更像是央行在“允许升值”和“防止过快升值”之间寻找平衡。这种平衡一旦被打破，就可能引发汇率定价的快速重估。

3. 韩元、欧元、澳元、墨西哥比索的隔夜贡献，揭示了人民币定价的“全球锚”正在变化

报告中的Fig. 1列出了对模型投影贡献最大的四个货币：韩元（-36个基点）、欧元（-

[... middle omitted ...]

元和澳元的贡献紧随其后，说明欧洲和大宗商品市场的变动正在通过人民币汇率定价机制传导至中国市场。墨西哥比索的贡献则进一步印证了全球贸易链的复杂性——北美、欧洲、亚洲的货币变动，最终都会在人民币中间价上留下痕迹。

对于投资者而言，这意味着单纯跟踪美元指数已经不足以判断人民币走向。需要建立一个更复杂的全球货币联动分析框架，特别是关注那些与人民币有间接竞争或互补关系的货币。

4. 报告尚未完全回答的关键问题：模型误差的收窄是暂时的，还是结构性的？

尽管NOM的模型释放了清晰的升值信号，但报告本身也留下了几个尚未完全回答的关键问题。最核心的一个是：模型误差的收窄和转正，是短期现象还是结构性变化？

从Fig. 2的模型误差历史数据看，误差波动幅度巨大——2025年初曾达到-1800个基点，2025年第三季度收窄至接近零，随后又再次扩大。这种高波动性意味着，单次模型投影的变动，并不足以确认趋势。如果未来几周模型误差再次扩大为负值，那么当前的升值信号就可能只是一次短暂的噪音。

\subsection{[R029] GS：市场真正低估的不是2万亿数据中心投资，而是“六张网”的宏观含义}
{\small\color{refgray}归类：中国政策与投资：‘六网’落地与资本管控信号再解读}

GS：市场真正低估的不是2万亿数据中心投资，而是“六张网”的宏观含义

一份关于中国政府计划在未来五年投入2万亿元建设数据中心的外媒报道，上周引发了市场对AI基建主题的再度关注。然而，这份来自GS中国宏观团队的研报给出了一个与市场直觉相反的判断：这个数字本身并非新闻，真正值得关注的，是它背后所揭示的中国财政扩张逻辑正在发生一次结构性切换。

市场往往将此类消息解读为“AI概念股利好”，但GS的分析框架提醒我们，如果只看到数据中心本身，就错过了这份报告最核心的洞察——中国正在通过“六张网”的投资框架，重新定义政府主导投资的投向、融资方式与节奏。理解这个框架，比猜测2万亿元是否超预期重要得多。

1. 2万亿数据中心投资不是增量信号，但“六张网”的加速落地才是

GS明确指出，所谓“2万亿元数据中心投资计划”并非新消息。早在今年3月“两会”期间发布的“十五五”规划中，数据中心已被纳入“六张网”的算力网络范畴。4月政治局会议和6月央视的再度强调，更多是对既有规划的重复确认，而非新增承诺。

那么市场为什么会有反应？GS的解释是，市场对政策信号的定价存在错位。投资者习惯于将单一项目的规模数字视为催化剂，而忽略了政策执行节奏的变化才是更关键的变量。真正的新信息并非2万亿元本身，而是近期政策沟通频率和项目准备工作的明显提速。

GS观察到，5月底以来，包括政策性银行新融资工具在内的多种财政工具已开始加速落地，部分地方政府在6月初即加快了项目推进。这与一季度GDP超预期后、二季度财政支出节奏明显放缓形成了对比。换言之，政策制定者正在从“等待数据”转向“提前准备”。

对这一判断的深层含义在于：如果投资者只盯着2万亿元这个数字是否被市场消化，就会错过“六张网”整体投资加速这一更大的叙事。GS估算，今年“六张网”总投资可能超过7万亿元，相当于固定资产投资的约14\%。数据中心的2万亿元在 年仅占FAI的0.8\%，但“六张网”的规模效应才是宏观层面的真正冲击。

2. 融资结构的悄然转变：从地方隐性债务到多元化“阳光融资”

市场对财政刺激的一个长期担忧是：钱从哪里来？会不会重蹈 年地方隐性债务扩张的覆辙？GS在报告中给出了一个值得深思的回答：这一轮“六张网”投资的融资结构正在发生根本性变化。

报告指出，资金来源将包括超长期特别国债（专门用于“两重”项目）、地方专项债、政策性银行新融资工具以及商业银行贷款。这个组合的核心意图是：在撬动私人资本的同时，避免上一轮周期中高风险、表外的地方融资平台借贷模式。

这里仍需验证的一个关键问题是：私人资本是否愿意入场？GS的报告没有完全展开这一层。从逻辑上推导，数据中心、新型电网等领域的投资回报周期较长，且涉及国家安全和技术自主，民营企业能否获得公平的参与机会、退出机制是否清晰，都会影响实际撬动效果。但至少从融资工具设计上看，政策层已经意识到了“阳光化”的必要性。

这一判断对资产定价的启示是：如果融资结构确实得到改善，那么政府主导投资的可持续性将高于市场预期，相关产业链的现金流可见度也会提升。反之，如果融资工具再次沦为地方债务的变相扩张，那么市场对“六张网”的定价就需要重新折价。

GS提出的一个反直觉判断是：二季度财政支出放缓并非因为“没钱了”，而是政策层的主动选择。报告用了两个数据来支撑这一观点：截至5月底，未使用的政府债券额度高达7.7万亿元（全年总额11.9万亿元），同时财政存款同比增加了7000亿元。

这意味着政策层手中有充足的“弹药”，但选择在二季度暂缓扣动扳机。GS的解释是，在当前“反应式（而非预防式）”的政策框架下，政策制定者倾向于在数据走弱后再出手，而非提前预判。一季度GDP超预期后，二季度刻意踩刹车，是为了避免经济过热或政策浪费。

但这一策略也带来了风险。报告提到，全

[... middle omitted ...]

半年政府主导投资的增速可能显著回升，而市场目前对这一节奏变化的定价仍不充分。

4. 报告尚未完全回答的关键问题：AI相关产业的10万亿元经济规模从何而来？

GS在报告中引用了发改委主任郑栅洁的一个预测：到2030年，AI相关产业的经济规模可能超过10万亿元。但报告也坦承，“这些产业的定义仍不清楚”。

这是一个需要审慎对待的伏笔。10万亿元的规模相当于当前中国GDP的约8\%，如果仅靠数据中心建设和算力基础设施，显然无法支撑这一数字。那么，这10万亿元究竟包含哪些领域？是AI芯片、大模型应用、智能制造，还是更广泛的数字化转型？不同的定义将导向完全不同的投资逻辑。

GS没有在这一问题上深入展开。从逻辑上推测，如果10万亿元主要来自硬件基础设施（如数据中心、芯片制造），那么相关公司的收入可见度较高，但竞争格局和利润率将是关键变量。如果更多来自下游应用和软件服务，那么市场空间虽大，但兑现周期更长、不确定性更高。这一问题的答案，将直接影响投资者对“六张网”主题的估值框架。

\subsection{[R030] 摩根斯坦利：市场误读了中国的资本管控信号，真正焦点在实体经济出海而非个人资金外逃}
{\small\color{refgray}归类：中国政策与投资：‘六网’落地与资本管控信号再解读}

摩根斯坦利：市场误读了中国的资本管控信号，真正焦点在实体经济出海而非个人资金外逃

过去几周，关于中国收紧境外投资监管的消息在市场上引发了一轮不小的焦虑。AIA、HSBC、渣打等香港金融股出现波动，一些投资者开始担忧中国正在逆转资本账户开放的方向。但摩根斯坦利最新发布的一份研报给出了一个与市场情绪截然不同的判断——这些新规的核心目标不是堵截个人资金外流，而是为中国企业日益增长的海外实体投资建立一个系统性的法律保护框架。

这份报告的价值不在于它提供了多少新的政策细节，而在于它拆解了一个关键问题：当市场把注意力集中在“资本外逃”这个叙事上时，真正重要的结构性变化是什么。如果我们只看新闻标题，很容易得出“中国在收紧”的结论。但如果仔细审视政策制定的参与者、文件的法律分类、以及具体的执行细则，会发现一个完全不同的故事。

这份摩根斯坦利研报最值得关注的判断是：当前的政策调整本质上是制度升级，而非方向逆转。中国正在从“逐案审批”的碎片化管理模式，转向“立法框架”的系统性保护模式。对于长期关注中国金融开放和香港市场地位的投资者来说，这意味着需要重新校准对政策风险的定价。

1. 国务院新规的核心信号：只有实业主管部门出席，没有金融监管机构参与

判断一项政策的真实意图，最简单的方法就是看谁在主导。摩根斯坦利在报告中指出，国务院新规的新闻发布会上，出席的只有司法部、发改委和商务部——没有任何金融监管机构参与。这个细节本身就是最直接的答案。

如果这是一项针对个人金融投资的收紧政策，银保监会、证监会、外汇局不可能缺席。这三家机构的缺席意味着，政策的目标群体是进行海外实业投资的企业，而不是进行证券投资的个人。

从文件的法律分类来看，新规被归入国务院文件体系中的“商贸、海关、旅游、对外经济合作”类别，而不是“金融监管”类别。这进一步确认了政策的性质：它是一部关于“走出去”的实业投资管理法，而不是关于“资本管制”的金融监管法。

报告进一步指出，新规的核心架构围绕对外直接投资（ODI）展开。政策明确支持市场化境外投资、投资者自主决策，并强调要建立覆盖法律、财税、金融、贸易、物流、移民、海关和贸易促进资源的海外服务体系。这些内容指向的是一个清晰的目标：帮助中国企业在海外更安全、更有组织地运营。

对于投资者来说，这个判断的含义是：新规对香港金融股的直接影响被市场夸大了。AIA、HSBC、渣打等公司的业务主要服务于合法的个人金融投资和保险需求，而不是企业ODI。摩根斯坦利明确认为，市场对这些公司的担忧“看起来过头了”。

2. 政策出台的真正驱动力：地缘政治风险与供应链重组背景下的企业保护机制

如果新规不是为了控制资本外逃，那它真正的政策动机是什么？摩根斯坦利给出了一个更宏观的解释框架：中国国际收支结构已经发生了根本性变化，政策制定者需要一套新的制度来管理这个新格局。

报告引用了摩根斯坦利宏观团队的判断——中国的国际收支已经进入“镜像结构”：大规模的经常账户盈余越来越多地被非官方主体转化为海外资产，而不是积累为央行外汇储备。这意味着，中国企业的海外投资规模正在快速扩大，而原有的碎片化监管框架已经无法满足需求。

新规第四条明确提到了“对接高标准国际经贸规则”、“推动高质量共建‘一带一路’”、“参与国际投资规则制定”、“支持产业链供应链合作”。这些表述指向的是一个更深层的政策目标：在复杂的地缘政治环境中，为中国企业的海外资产和运营提供法律保护。

报告指出，新规强调国家安全审查、海外安全风险预警、争端解决机制、对歧视性限制的反制措施。这些条款的核心逻辑是：单个企业在面对外国政府的审查或限制时，谈判能力有限；而通过中央立法框架，政府可以以国家力量为后盾，为企业在海外争取更有利的条件。

这个判断对投资者的含义是：新规的长期效果可能是

[... middle omitted ...]

着：个人境外投资被承认存在，但它不是这份法规的核心议题。摩根斯坦利明确认为，“将这份法规定性为由资本外流压力驱动是不恰当的”。

但这并不意味着个人境外投资领域没有变化。摩根斯坦利同时分析了证监会近期对跨境股票交易的清理行动，并给出了一个精准的区分：监管关注的是“发生在境内的活动”，而不是“发生在境外的账户”。

具体来说，受影响的是“境内投资者直接在境内交易境外股票”——这不再被允许。但不受影响的是：在境外开立投资账户、银行账户、交易或购买保险。证监会的要求是：境外券商不能在中国境内提供证券经纪服务，包括通过APP、社交媒体提供营销、开户指导或交易服务。

这个区分非常关键。摩根斯坦利强调，香港证监会6月10日的澄清明确表示，持牌公司仍可为来自中国大陆的合格投资者开立账户，前提是资金全部来自境外合法来源，并且严格遵守KYC和尽职调查要求。同时，香港证监会也指示券商在服务境内居民时必须遵守当地规则——这意味着即使账户在境外开立，券商也不能在客户身处境内时提供交易服务。

\subsection{[R031] Bernstein：市场低估的不是电力需求，而是非日照时段的供给刚性}
{\small\color{refgray}归类：新兴市场专题：印度电力时段性短缺与日本半导体设备份额反转}

Bernstein：市场低估的不是电力需求，而是非日照时段的供给刚性

印度电力板块在过去二十年里，几乎被市场当作“空调股”来交易——夏季来临前买入，季风来临前卖出。这种季节性博弈在2024年中达到顶峰，随后需求增速骤降至同比约1\%，持续了近20个月。今年，天气、能源安全与数据中心三大因素重新点燃了市场热情。但Bernstein这份研报的核心判断是：市场正在低估一个结构性变化——印度电力短缺的真正瓶颈，已经从总量转向时段，而非日照时段的供给刚性在未来三年内几乎无法缓解。

这不仅仅是印度电力股的投资信号。对于关注全球制造业转移、数据中心布局与新兴市场基础设施投资的决策者而言，这份报告提供了一个难得的观察窗口：当一个经济体的电力系统从“总量过剩”切换到“时段性短缺”时，资产定价的逻辑正在被重写。

1. 表面数据在说“宽松”，但拆解时段后看到的却是持续恶化的晚间缺口

如果只看印度电力现货市场的总量数据，结论很可能是“供需正在趋于平衡”。Bernstein指出，今年4-5月现货市场的平均电力价格约为3.8卢比/千瓦时，显著低于2023年的高点；买卖报价比也低于1，意味着卖盘多于买盘。这些指标似乎在暗示紧张局势正在缓解。

但问题出在平均数据掩盖了结构性分化。当Bernstein将数据按“一天中的时段”进行拆解后，一个截然不同的图景浮现出来：早晨与晚间的电力价差已经扩大到约7卢比/千瓦时，并且仍在扩大。而在过去十年里，这个价差几乎为零。晚间时段的买卖报价比在今年夏季达到2.2倍，而2020年仅为0.5倍。

这意味着什么？总量层面的宽松，完全是由日照时段的太阳能电力供给充裕造成的。而真正的压力——非日照时段——正在以肉眼可见的速度恶化。对于任何依赖晚间连续生产或高负荷运行的工业用户来说，这构成了一个真实且不断上升的运营成本风险。

2. 需求端并非没有放缓，但结构性增长动力正在取代周期性波动

Bernstein承认，过去一段时间印度电力需求增速确实令人失望。FY26年度需求仅增长约1\%，这让市场对其此前“电力需求增速为实际GDP增速1倍”的长期假设产生了质疑。但报告认为，在低基数和支持性天气因素（降雨预报低于正常水平）的共同作用下，今年6\%的需求增长假设现在看来甚至可能偏保守。

更关键的是，需求增长的性质正在发生变化。Bernstein明确指出，电气化与能源安全是全球性主题，而印度作为缺乏油气资源的国家，这个主题的重要性远超其他经济体。这意味着，电力需求不再仅仅是“夏天开空调”的季节性脉冲，而是由制造业扩张、数据中心建设与电气化进程驱动的结构性需求。

对于投资者而言，这意味着判断电力股价值的框架需要从“需求增速是否超预期”转向“供给能否跟上结构性需求的节奏”。而恰恰在供给端，报告给出了一个令人不安的结论。

3. 供给侧的真相：不是“缺电”，而是“缺晚间电”，且未来三年看不到缓解

这是整份报告最具洞察力的部分。Bernstein将本轮周期与上一轮进行了对比：上一轮周期（2012年左右），印度有约100吉瓦的燃煤电厂正在建设中，相当于当时已安装燃煤装机容量的约75\%。而本轮，正在建设的燃煤容量仅为约40吉瓦，仅占当前装机容量的约16\%。

更令人担忧的是未来三年的新增供给。Bernstein预计，未来三年印度仅能新增约15吉瓦的燃煤容量。而即使在6\%的保守需求增长假设下，晚间电力需求累计将增长约50吉瓦。缺口是巨大的。

理论上，电池储能系统（BESS）可以填补这个缺口。但Bernstein对此持谨慎态度：过去一年确实出现了大量BESS招标，但多数由小型玩家以激进报价中标。此后，BESS价格与卢比/美元汇率都变得不利。报告明确表示，“我们预计其中许多项目将无法落地”。

这一判断意味着什么？晚间电力短缺在未来

[... middle omitted ...]

压力。Bernstein在报告中明确表示，其偏好排序是“非太阳能电力（燃煤-核电-电池）优于太阳能”。

对于产业决策者而言，这意味着在选址新产能或数据中心时，不仅要考虑总电力可得性，更要评估“晚间电力可得性”这个变量。那些能够与提供晚间电力的电厂签订长期购电协议的区域，将获得显著的竞争优势。

5. 报告尚未完全回答的关键问题：电池储能的真实经济性拐点何时到来？

Bernstein对BESS项目的落地前景持怀疑态度，这一判断基于BESS价格与汇率的不利变化。但报告并未深入讨论一个关键问题：如果晚间电价持续走高（目前早晨与晚间的价差已达7卢比/千瓦时），BESS的经济性拐点是否会被提前触发？

从逻辑上推理，晚间电价每上升1卢比，BESS项目的内部收益率就会改善一个可观的幅度。如果晚间电价在 年继续攀升，那些此前被认为不可行的BESS项目可能会重新变得有吸引力。但这里存在一个二阶效应：如果大量BESS项目同时上马，晚间电价可能被迅速压低，从而压缩后续项目的回报空间。

\subsection{[R032] Bernstein：日本半导体设备在中国市场的份额流失并非结构性，反转正在酝酿}
{\small\color{refgray}归类：新兴市场专题：印度电力时段性短缺与日本半导体设备份额反转}

Bernstein：日本半导体设备在中国市场的份额流失并非结构性，反转正在酝酿

市场对日本半导体设备商在中国市场份额持续下滑的担忧，可能夸大了问题的严重性。Bernstein最新发布的一份深度研报指出，2024至2025年间日本设备商在中国进口市场的名义份额从26\%降至23\%，表面上看是连续三年的下滑趋势延续。但报告的核心判断是：这一流失主要由汇率和客户下单节奏等暂时性因素驱动，而非竞争力被中国本土或海外对手结构性侵蚀。随着日元贬值红利转化为定价空间、中国客户投资节奏恢复正常化，日本设备商有望从2026年下半年起重新收复失地。

这份报告之所以值得关注，不仅因为它挑战了市场对“日本设备在中国节节败退”的简单叙事，更因为它提供了一个拆解复杂市场份额变化的分析框架——在汇率波动超过30\%、中国自主化进程加速、客户投资周期分化的背景下，名义份额数据本身可能产生严重误导。对于持有或关注东京电子、科休半导体、斯库林等日本前道设备龙头的投资者而言，报告的判断直接关系到未来12至18个月的估值叙事能否逆转。

1. 日元贬值31\%带来的名义份额流失，掩盖了日本设备实际竞争力的稳定

Bernstein报告首先做了一件必要但容易被市场忽略的工作：剔除汇率影响重新计算日本设备商在中国的真实市场份额。过去五年日元对美元贬值约31.5\%，而日本设备商普遍以日元报价，这意味着以美元计价的进口统计会系统性地低估日本设备的实际出货量。

调整后的数据令人意外。2024年日本设备商在中国进口市场中的名义份额为26\%，但经汇率调整后实际份额为31\%；2025年名义份额降至23\%，调整后为26\%。如果将2024和2025年合并平均，调整后的份额为28.6\%，与2021至2023年间的平均水平28.1\%几乎持平。换言之，过去两年日本设备商在中国市场的“实际”份额并未出现结构性下滑，而是围绕一个稳定中枢小幅波动。

这个发现的意义在于：市场此前将名义份额下降解读为日本设备竞争力下降，但Bernstein的分析表明，汇率变动解释了一大部分波动。对于投资者而言，这意味着如果日元趋于稳定甚至升值，日本设备商在中国市场的“名义”份额将自动获得修复，无需等待任何竞争格局变化。这是一个被低估的顺风因素。

2. 客户集中度与投资节奏的错位，是2025年份额下滑的更直接原因

除了汇率，Bernstein将另一部分份额流失归因于中国客户投资节奏的阶段性错位。报告特别点名了长鑫存储——这家对日本设备依赖度极高的DRAM厂商，在2025年大幅缩减了资本开支，直接拖累了日本设备商在中国市场的出货量。

具体来看，科休半导体2024年来自长鑫存储的收入约为470亿日元，2025年骤降至约110亿日元；斯库林同样在过去两年几乎没有从长鑫存储获得新订单。东京电子的情况略有不同，其份额下滑更多来自客户结构和工艺步骤组合的变化，而非单一客户依赖。

Bernstein的分析揭示了一个关键点：日本设备商在中国市场的表现高度受制于少数几家核心客户的资本开支周期。当这些客户进入投资空窗期或转向国产设备测试时，日本设备商的出货数据就会明显走弱。但这种走弱是“客户集中度风险”的体现，而非“产品竞争力下降”的证据。

对于投资者而言，这意味着判断日本设备商在中国的前景，不能只看笼统的市场份额数据，而需要拆解到具体客户和工艺段。报告指出，长鑫存储的投资有望在2026年恢复，同时斯库林已确认来自晋华集成的可观订单，YMTC的采购量也将同比增长50\%。如果这些客户层面的积极信号兑现，日本设备商的出货数据将在2026年下半年出现明显改善。

3. 日元贬值创造定价空间，日本设备商正从“被动降价”转向“主动提价”

Bernstein报告中最具洞察力的判断之一，是日元大幅贬值实际上为日本设备商创造

[... middle omitted ...]

断的隐含含义值得深思：日元贬值在初期侵蚀了日本设备商的名义市场份额，但经过一段时间的传导，反而可能转化为定价权和利润率提升的催化剂。对于投资者而言，这意味着日本设备商在中国市场的收入可能不再只是量的故事，还可能叠加价的故事。如果提价成功，即使出货量持平，收入和利润也能实现增长。这是市场当前定价中尚未充分反映的因素。

4. 前道设备相对商品半导体的落后表现，为反转提供了估值安全垫

报告还观察到一个重要的市场现象：过去2至3个月，半导体设备类股明显跑输商品半导体（存储、模拟、衬底、硅片等）。东京电子、斯库林、科休半导体的股价表现落后于行业平均水平。

Bernstein认为这种落后是不合理的。随着投资者逐步认识到未来几年资本开支将增加，以及日本设备商在中国市场份额有望恢复，设备类股的相对表现应该会改善。报告在估值表中给出了具体判断：东京电子目标价59,200日元，评级跑赢大盘；科休半导体目标价8,240日元，评级跑赢大盘；斯库林目标价12,600日元，评级与大市持平。

\subsection{[R033] Citi：中国新能源车市真正的拐点不是渗透率破60\%，而是龙头集中度开始松动}
{\small\color{refgray}归类：电动车与汽车：区域利润结构重构与燃油车定价崩塌}

Citi：中国新能源车市真正的拐点不是渗透率破60\%，而是龙头集中度开始松动

中国新能源汽车市场五月销量数据出炉，渗透率首次稳稳站上61\%。这个数字本身并不令人惊讶——市场早已预期2026年渗透率将突破60\%大关。真正值得关注的信号，藏在Citi这份五月数据库更新中：CR5（前五大厂商集中度）在五月为57.0\%，环比微增0.4个百分点，但同比下降了2.0个百分点。

这不是一个简单的数字波动。它意味着，在渗透率快速攀升的“上半场”结束后，中国新能源车市正在进入一个没有历史参照的新阶段——头部阵营的份额不再单向扩张，而是出现了结构性松动。对于投资者和产业决策者而言，过去几年“买龙头就能分享行业增长”的简单逻辑，可能需要重新审视。

Citi分析师Jeff Chung和Kyle Wu在报告中呈现的，远不止一份月度销量统计。它是一张正在变化的竞争版图，其中既有比亚迪的份额压力、新势力的分化加速，也有“第二梯队”玩家以超出预期的速度蚕食头部空间。这些信号叠加在一起，指向一个更深层的判断：市场真正低估的不是需求，而是供给侧竞争格局的再定价。

1. 渗透率突破六成之后，增长引擎从“增量红利”切换为“存量争夺”

五月新能源乘用车批发销量135.2万辆，同比增长11\%，环比增长10\%。纯电动（BEV）批发88.6万辆，同比增速达到17\%，环比14\%；插电混动（PHEV）46.6万辆，同比仅增长1\%，环比增长4\%。前五个月累计新能源乘用车批发531万辆，同比增长1\%。

渗透率61.1\%，同比提升8.4个百分点，环比提升3.0个百分点。这个数字放在全球任何主要汽车市场都是里程碑式的——中国新能源车已经不再是“替代者”，而是绝对主流。但Citi报告的数据也揭示了一个容易被忽视的事实：在渗透率快速提升的同时，整体市场的同比增速正在放缓。五月11\%的同比增速，虽然高于前五个月1\%的累计增速，但远低于 年动辄50\%以上的增速区间。

这意味着什么？行业正在从“所有人都在抢增量”的阶段，过渡到“从对手盘口中抢份额”的阶段。增量红利在收窄，竞争烈度在上升。对于企业而言，谁能在存量市场中建立可持续的竞争优势，谁就能在下一轮洗牌中存活下来。对于投资者而言，过去那种“只要赛道正确，龙头自然受益”的配置逻辑，正在被“谁能守住份额、谁能提升议价权”的选股逻辑所取代。

Citi这份报告最值得细读的部分，不是渗透率数字本身，而是渗透率背后各家厂商份额的此消彼长——这才是判断下一阶段投资方向的关键线索。

2. 比亚迪份额跌破28\%：规模领先者的“舒适区”正在被侵蚀

五月比亚迪批发销量37.7万辆，同比零增长，环比增长20\%。市场份额27.9\%，同比下降3.0个百分点，环比仅回升2.2个百分点。前五个月累计市场份额26.0\%，同比下降7.1个百分点。

这是这份报告中最值得警惕的信号之一。作为连续多年稳坐中国新能源车头把交椅的比亚迪，其市场份额在2026年出现了明显松动。五月同比下滑3个百分点，累计同比下滑超过7个百分点，这个幅度在比亚迪的历史数据中并不多见。

比亚迪的问题不在于绝对销量——37.7万辆的单月批发量仍然遥遥领先第二名吉利（13.1万辆）近三倍。问题在于增速。在行业整体增长11\%的背景下，比亚迪同比零增长，意味着它的增长完全落后于市场。从产品结构看，主力车型宋DM环比增长15\%至6.2万辆，海豚环比增长2\%至2.2万辆，但海鸥环比暴增47\%至3.99万辆——这说明比亚迪的增长主要依靠低端车型放量，高端化进程似乎遇到瓶颈。

比亚迪的规模优势依然巨大，但“大”不等于“强”。当竞争对手在产品定义、渠道效率、品牌溢价等维度上不断逼近，规模本身带来的成本优势正在被稀释。这份报告暗示了一个关键问题：比亚迪能否在未来12个月内重新证明其定价权

[... middle omitted ...]

上汽集团（含合资）新能源批发4.6万辆，同比增长123\%，市场份额3.4\%，同比提升1.7个百分点。

这些数字背后是一个清晰的趋势：在头部厂商增长放缓的同时，一批原本被认为“追赶者”的玩家正在以远超行业平均的速度崛起。奇瑞依靠的是其传统燃油车渠道的转化能力和成本控制；零跑则凭借性价比定位和持续的产品迭代，成为这一轮增长中最亮眼的黑马之一。

值得注意的是，零跑五月市场份额已经达到6.0\%，超过了蔚来（2.8\%）、理想（2.5\%）、小鹏（2.4\%）等长期被市场关注的新势力。这种“闷声发大财”的态势，说明市场对新能源车竞争的认知可能仍停留在“蔚小理”的叙事框架中，而忽略了真正在份额上快速扩张的玩家。

对于投资者而言，这意味着选股视野需要从“头部三强”扩展到“第二梯队中的结构性赢家”。Citi报告虽然没有直接给出投资建议，但数据本身已经提供了清晰的观察框架：谁的增速高于行业、谁的份额在持续提升、谁的产品矩阵覆盖了最主流的细分市场——这些才是判断下一阶段竞争力的核心指标。

\subsection{[R034] GS：鲜零食零售的兴起不是威胁，而是价值零售龙头品类升级的催化剂}
{\small\color{refgray}归类：地产与消费：分化中的结构性机会与成本端出清}

GS：鲜零食零售的兴起不是威胁，而是价值零售龙头品类升级的催化剂

中国消费市场从来不缺新概念。2025年至今，鲜零食零售连锁以200多家门店、十余个新品牌的速度集中涌现，日销超过10万元、毛利率30\%以上、EBITDA利润率达到两位数——这些数字让市场再次兴奋，也带来了一个自然的问题：这对已经跑通模型的食品饮料价值零售商（如鸣鸣很忙、万辰）意味着什么？

GS最新发布的研报给出了一个清晰且反直觉的判断：鲜零食零售的崛起，非但不是对现有价值零售龙头的威胁，反而为它们提供了品类扩张、店型升级和竞争壁垒加固的三重机遇。市场如果只看到“新业态抢流量”，就低估了规模龙头在供应链密度、运营系统和会员资产上的结构性优势。

这份报告的价值不在于罗列几个新品牌的数据，而在于它提供了一个分析框架：当消费零售行业出现新物种时，如何判断它是颠覆者还是催化剂。以下是我们从报告中提炼出的三个核心层次。

1. 鲜零食零售证明了零食消费的韧性，但真正的变量是“高频复购”和“单店天花板”

GS报告的第一个洞察是：鲜零食零售的繁荣，首先验证了零食消费的基本盘依然强劲。这不是零和博弈，而是品类边界的扩展。

从数据看，鲜零食零售的店效模型确实令人印象深刻。以头部品牌金粒门为例，其旗舰店年化GMV可达2400万至4800万元，日均订单量在1300至2600单之间，客单价约50元。相比之下，鸣鸣很忙和万辰的典型门店年化GMV约在500万元左右，日均订单约400至500单，客单价约31元。鲜零食门店的坪效（每平米销售额）达到8万至12万元，是价值零售门店（约3.6万元）的两到三倍。

但更值得关注的不是绝对值，而是驱动这些数字的底层逻辑。GS指出，鲜零食零售的核心优势在于两个维度：

第一，高频复购。鲜零食的短保质期属性（多数SKU保质期在1至5天）天然创造了更高的购买频次。消费者不是为了囤货进店，而是为了即时满足——早餐、下午茶、夜宵、佐餐。这种消费场景的扩展，使得门店不再是“路过买一包”的随机交易，而成为日常生活的固定节点。

第二，单店天花板提升。鲜零食门店通过品类组合（坚果、卤味、酥饼、现制茶饮）覆盖了更多消费时段。GS报告提到，鸣鸣很忙已经在2025年下半年尝试蛋挞和热狗等鲜食产品，并且验证了这些品类对单店GMV的拉动效果。

这意味着什么？对于价值零售龙头而言，鲜零食不是要替代它们，而是提供了一个经过市场验证的“品类升级路线图”。如果鸣鸣很忙和万辰能够将鲜零食的逻辑嫁接到自身庞大的门店网络中，它们将不再是“卖包装零食的折扣店”，而是“社区食品零售平台”。

但这里有一个关键前提：品类扩张的成功取决于供应链和周转效率。短保质期意味着更高的损耗风险，这不是所有玩家都能驾驭的。

2. 规模龙头在供应链密度和运营系统上的优势，正在从“成本护城河”升级为“品类扩张权”

市场对新进入者的担忧往往集中在“它们会不会抢走份额”。GS的分析给出了一个冷静的判断：鲜零食零售在短期内不会对价值零售龙头构成实质性竞争威胁，原因是这个模式在规模化时会遇到三个结构性瓶颈。

第一个瓶颈是SKU管理。鲜零食门店的SKU通常在数百到一千个之间，远少于价值零售的数千个SKU。但鲜零食的难点在于，每个SKU的保质期极短，需要精准的销售预测和快速的库存周转。这意味着运营复杂度呈指数级上升，而不是线性增长。

第二个瓶颈是选址和租金。鲜零食门店依赖高客流、高坪效，因此必须占据购物中心B1层等核心位置。GS的UE模型显示，金粒门旗舰店的月租金可达8万至20万元，是价值零售门店（约1.6万元）的5到12倍。这种成本结构意味着，鲜零食模式在下沉市场的渗透难度远高于价值零售。

第三个瓶颈是供应链和冷链。鲜零食的自制比例极高（部分品牌达到99\%），需要中央厨房、冷链物流和区域工

[... middle omitted ...]

“损耗率失控”拖垮。

第三，会员资产和复购率。价值零售龙头拥有数千万会员基础，获客成本远低于新品牌。在鲜零食的“高频复购”逻辑下，会员资产的杠杆效应会更加明显。

所以，这份报告真正想表达的是：鲜零食不是价值零售的终结者，而是价值零售龙头的“能力放大器”。那些在供应链、运营和会员资产上已经建立起壁垒的玩家，将更有能力把鲜零食的高频、高客单逻辑整合进自己的模型。

3. 报告尚未完全回答的关键问题：鲜零食的UE模型能否在规模化后保持吸引力？

GS的报告做了一件事：它用UE模型（单位经济模型）对比了鲜零食和价值零售的财务结构。但报告也坦诚地指出，鲜零食零售仍处于早期阶段，门店数量极少，因此这些UE数据是“标杆门店”而非“平均门店”。

第一，鲜零食门店在扩张到100家以上时，单店UE是否会显著下降？目前头部品牌的门店数在20到100家之间，且集中在高线城市的核心商圈。当它们进入次核心商圈或低线城市时，客流量和客单价的下滑是否会被租金下降所抵消？报告没有给出明确的答案。

\subsection{[R035] GS：IP零售的“高基数疲劳”正在重塑市场对增长的定价逻辑}
{\small\color{refgray}归类：地产与消费：分化中的结构性机会与成本端出清}

当市场还在为潮玩IP的爆发力欢呼时，一份来自GS的月度追踪报告揭示了一个更微妙的信号：中国IP零售的核心玩家正在经历“高基数疲劳”。这不是简单的增速放缓，而是增长引擎从“流量红利”切换到“IP运营能力”的关键转折点。

这份5月更新的报告覆盖了泡泡玛特、名创优品和布鲁可三家代表性公司。数据显示，泡泡玛特国内线上销售额同比增长29\%，但相比4月的95\%明显减速；名创优品线上增速从一季度的43\%降至5月的21\%；布鲁可虽然在加速推新，但2季度至今的线上增速仍低于一季度。这些数字单独看并不差，但放在一起，它们指向一个共同问题：IP零售的增长故事正在进入新阶段，市场需要更精细的观察框架。

更值得关注的是，GS在海外市场也捕捉到了类似信号。泡泡玛特美国市场5月信用卡销售同比下降36\%，虽然较4月的43\%降幅有所收窄，但2季度至今仍维持约40\%的同比下滑。名创优品美国市场增速虽从4月的19\%回升至5月的31\%，但2季度累计增速仍低于一季度。这些数据意味着，即使在中国IP零售商的海外扩张故事中，基数效应和消费者情绪波动正在成为不可忽视的变量。

对于决策者而言，现在需要回答的核心问题不是“IP零售是否还有增长空间”，而是“下一阶段的增长来自哪里，以及哪些公司具备穿越周期的能力”。

1. 泡泡玛特的减速不是需求问题，而是供给节奏与基数效应的双重挤压

GS报告中最引人注目的数据是泡泡玛特国内线上增速从4月的95\%骤降至5月的29\%。这个数字本身已经足够说明问题，但背后的结构性原因更值得深究。

从供给端看，5月的新品表现明显弱于此前。GS追踪显示，5月推出的“Dear Birds”多IP系列在抖音和天猫的首月销量分别超过3.2万和1万件，但截至6月8日仍未售罄。对比此前“Have a good run”系列首月在抖音超20万件的销量，差距显著。Crybaby和Hacipupu的毛绒新品同样未能在首日售罄，只有Zsiga在抖音实现了售罄。新品拉动力的下降，直接反映在线上增速上。

从基数效应看，2025年6月Labubu的“Big Into Energy”系列曾带来供给高峰，这意味着即将到来的6月同比基数将更为严峻。GS指出，截至6月6日，抖音旗舰店的6月日销已落后于5月同期。如果这一趋势延续，6月线上增速可能进一步承压。

更关键的是，GS追踪的二手市场价格数据并未显示IP热度在5月有明显改善。Labubu的样本SKU折价率从4月的30\%-45\%进一步扩大；Molly和Dimmo的折价同样在扩大，只有Skullpanda和Crybaby出现企稳或小幅回升。Lisa世界杯MV中Labubu的露出虽然获得了1360万次观看，但Google搜索指数并未出现有意义的上扬。

这些数据合在一起意味着：泡泡玛特目前面临的不是需求断崖，而是“供给节奏未能匹配市场预期”与“高基数”的组合压力。对于投资者而言，需要关注的是6月世界杯能否成为Labubu IP的催化剂，以及公司能否在7-8月推出具有更强拉动力的新品系列。

2. 名创优品的IP加速策略正在见效，但增速放缓暴露了规模扩张的边界

名创优品5月的表现相对乐观，但也存在隐忧。GS报告显示，公司5月线上销售额同比增长21\%，高于4月的10\%，但低于一季度43\%的增速。管理层披露4-5月同店销售正增长，五一期间销售额同比增长高双位数，创历史新高。

名创优品5月加速了IP联名节奏，推出了自有IP Yoyo、星球大战、玩具总动员等电影IP，以及Chiikawa x Sanrio、Loopy等热门IP。Yoyo甚至出现在Met Gala红毯上，作为模特Paloma Elsesser的包饰。这些动作表明，名创优品正在将IP联名从“偶尔爆款”升级为“常态化运营”。

但2季度至今的增速放缓值得

[... middle omitted ...]

多的SKU”。

GS报告显示，布鲁可5月加速了新品发布节奏，2季度至今推出的系列和SKU数量同比和环比均有所增加。在天猫旗舰店，除宝可梦外，大多数成人向IP的5月销量环比均有所上升。整体来看，天猫+京东+抖音的线上销售额同比增速较4月有所加快。

具体数据方面，假面骑士Legend系列在5月贡献了约7000件销量和70万元销售额，是布鲁可当前最核心的IP驱动引擎。EVA Legend系列虽然销量稳定在1500件左右，但销售额贡献已从峰值的175万元降至15万元，表明该IP已进入成熟期。初音未来、名侦探柯南、火影忍者等IP的销量和销售额均在低位徘徊。

这种分化意味着，布鲁可的IP矩阵中真正具备持续变现能力的IP仍然有限。假面骑士的成功能否复制到其他IP，是公司需要回答的关键问题。GS报告没有完全展开这一点，但数据已经暗示，IP授权生意的本质是“少数头部IP贡献大部分收入”，而非均匀分布。

4. 海外市场：美国消费情绪恶化是最大变量，但泡泡玛特和名创优品面临不同处境

\subsection{[R036] JPM：AI驱动的半导体超级周期才刚刚进入第二阶段}
{\small\color{refgray}归类：AI/算力：从模型竞赛到成本信任，供应链紧俏重塑定价权}

2026年4月的全球半导体销售数据，让市场上大多数关于“周期见顶”的讨论显得为时过早。JPM最新发布的亚洲半导体行业报告显示，当月整体半导体营收同比增长106\%，这是自1994年以来的最高增速，也是连续第八个月加速。这个数字本身已经足够震撼，但真正值得决策者关注的，不是增速本身，而是支撑这一增速的结构性力量正在从“单一引擎”切换为“多引擎驱动”。

这份报告的核心判断是：市场对AI半导体的定价逻辑仍停留在“需求爆发”的第一阶段，但供给侧的硬约束——尤其是先进制程和HBM的产能瓶颈——正在成为未来2-3年更关键的定价因子。这意味着，当前的营收增长不是周期性的脉冲，而是结构性供需失衡的早期信号。

整体半导体营收同比增长106\%，其中逻辑半导体增长33\%，存储器增长359\%。这些数字的含金量在于它们的“广度”。JPM指出，4月半导体出货量同比增长15\%，较3月的10\%明显加速，说明增长不再仅仅依赖价格提升，而是有实实在在的出货量支撑。ASP同比上涨79\%，其中存储器ASP飙升198\%，逻辑ASP上涨18\%。

更深层的信号在于增长的“去集中化”。虽然AI加速器仍然是核心驱动力，但报告明确提到，逻辑半导体的增长已经扩展到服务器CPU、网络芯片，以及苹果iPhone/Mac产品线的持续需求。这意味着，AI对半导体的拉动正在从数据中心内部向外溢出，进入终端设备和基础设施层面。

对于产业观察者而言，这意味着一个关键判断：如果AI半导体的需求来源变得更加分散，那么单一客户或单一应用的订单波动对整体产业链的冲击力将显著下降。这是行业进入更稳定增长阶段的前提条件。

JPM在这份报告中给出了一个值得反复推敲的数字：2026年N3制程的供需缺口预计约为60万片晶圆，对应150亿至180亿美元的代工收入。这个缺口的背后，是几乎所有主流AI加速器——包括NVIDIA Rubin、Google TPU v8t/i、AWS Trn 3——都计划采用N3制程。

报告认为，N3将是AI加速器的主要瓶颈。这个判断的逻辑链条清晰：需求端，美国四大云服务商的资本支出在2026年仍将保持80\%以上的同比增长，且这一趋势正在变得更加“结构性”；供给端，N3的产能扩张受限于设备交付周期、技术良率爬坡以及物理空间的限制。JPM预计，先进制程的紧张状态将持续到2027年甚至2028年初。

这对产业链的启示是：在N3产能真正释放之前，拥有先进制程产能的代工厂将获得持续的超额定价权。而那些依赖先进制程的芯片设计公司，可能面临“有订单无产能”的交付风险，这将加速他们向更早期介入产能锁定的策略转变。

3. HBM的供需失衡正在从“短期紧张”演变为“长期结构性缺口”

存储器部分的增长是最极端的——4月营收同比增长359\%，ASP上涨198\%。JPM的团队已经将2026年至2028年的存储器TAM分别上调了37\%和53\%。但真正值得关注的不是这些数字本身，而是HBM对存储器产业结构的重塑。

报告预测，未来三年HBM的DRAM晶圆分配比例将从2026年的24\%上升到2028年的31\%，而HBM的bit安装需求复合年增长率将达到85\%。这意味着，HBM正在从存储器的一个细分品类，转变为整个DRAM产业的“需求锚点”。HBM的定价逻辑将溢出到整个DRAM市场，因为产能向HBM倾斜必然挤压传统DRAM的供给。

JPM还预计，存储器厂商将在明年推动HBM ASP上调10\%（同类比较）或30\%（混合比较）。这个判断背后是一个尚未被市场充分定价的博弈：当HBM与传统DRAM的每bit价格差距持续扩大时，存储器厂商有强烈的动力去重新定价HBM，以反映其更高的技术附加值和稀缺性。

对于投资者而言，这意味着存储器周期的传统框架——基于供需的周期性波动——可能已经失效。HBM

[... middle omitted ...]

这一趋势持续，成熟制程的代工厂和IDM厂商可能会迎来一轮意外的营收增长。

但报告也明确指出了下行风险：消费电子需求仍然疲软，尤其是中低端智能手机，在存储器成本上升的背景下，2026年下半年的压力可能进一步加大。苹果iPhone 17的需求相对强劲（出货量同比增长9\%），但这更多是结构性份额提升的结果，而非整体市场的回暖。

这意味着，半导体产业链的下游分化正在加剧。高端和AI相关领域处于供给紧张状态，而中低端消费电子仍在消化库存。这种分化对投资者的含义是：不能再用“半导体”这个整体概念来做投资决策，必须深入到具体制程、具体应用和具体客户结构。

5. 报告尚未完全回答的关键问题：资本开支的回报率能否支撑如此大规模的投资

JPM在这份报告中大幅上调了存储器资本开支预测，未来三年的总支出预计达到4500亿美元，较此前预测的3000亿美元增加了50\%。DRAM的年底产能将从2025年的约192万片WSPM增加到2028年的280万片，其中60\%的增量产能将分配给HBM。

\subsection{[R037] 摩根斯坦利：AI需求的结构性韧性正在重塑台湾科技股的定价逻辑}
{\small\color{refgray}归类：AI/算力：从模型竞赛到成本信任，供应链紧俏重塑定价权}

摩根斯坦利：AI需求的结构性韧性正在重塑台湾科技股的定价逻辑

五月营收数据出炉，台湾科技板块整体同比增长39\%，环比增长4\%，表面上是AI需求持续抵消消费电子季节性疲软的又一个佐证。但这份摩根斯坦利研报真正值得关注的信号，不是总量增长本身，而是增长质量的分化：数据中心硬件和智能手机供应链超出预期，而IC代工全线miss。这意味着市场此前对AI溢价的定价逻辑可能需要重新校准——真正受益的不是所有与AI沾边的环节，而是那些具备订单可见性优势或产业链议价权的公司。

这份报告发布后，台湾加权指数和相关科技股已累计上涨11\%，但摩根斯坦利的推荐清单却高度集中：MediaTek、Macronix、Delta、Bizlink、Accton、Wiwynn、Yageo、Zhen Ding。这八只股票的共同特征，并非同属AI赛道，而是各自在供应链中拥有不可替代的位置或清晰的订单能见度。换言之，市场正在从一个“AI概念普涨”的阶段，进入一个“AI业绩必须兑现”的阶段。

以下是我们从这份研报中提炼出的五个核心洞察，以及一个报告尚未完全回答的关键问题。

1. IC代工集体miss揭示的不是需求问题，而是产能定价权的再分配

五月营收中，IC代工板块成为最突出的低于预期区域。台积电、联电、世界先进、力积电等主要代工厂无一例外地miss了市场预期，其中台积电5月营收416,975百万新台币，低于预期的451,798百万新台币，miss幅度达8\%。联电同样miss 6\%，世界先进miss 2\%，力积电是唯一beat的例外，但幅度仅有2\%。

这些数据合在一起意味着什么？不是AI需求减弱了，而是代工环节的产能紧张程度正在缓解。台积电的miss尤其值得关注——作为全球最先进的制程供应商，其营收连续低于预期可能暗示两个结构性变化：其一，AI芯片客户正在将部分订单从先进制程向更成熟的制程迁移，以优化成本结构；其二，非AI领域的消费电子、车用芯片需求复苏速度慢于预期，导致代工厂的产能利用率未能达到市场乐观假设。

对投资者而言，这意味着过去两年“代工厂因产能紧缺而享有定价权”的逻辑正在松动。代工环节的议价能力正在从卖方市场向买方市场缓慢倾斜。那些依赖代工产能的IC设计公司，反而可能因此获得成本改善的空间。

与代工板块形成鲜明对比的是，数据中心硬件成为五月营收的最大亮点之一。Gold Circuit beat预期14\%，KingSlide beat幅度高达43\%，Giga-Byte beat 17\%，Hon Precision虽然miss但营收同比增长84\%。这些数据的共同指向是：AI基础设施投资仍在加速，且已经从GPU采购阶段扩展到网络、散热、连接器、PCB等配套环节。

但摩根斯坦利在推荐清单中并未覆盖所有数据中心硬件公司，而是选择了Accton、Bizlink、Wiwynn。这三家公司的共同特征是什么？Accton是网络交换设备供应商，Bizlink是高速连接器厂商，Wiwynn是服务器系统集成商——它们都处于AI数据中心建设中“必须经过”的节点，且订单能见度通常超过两到三个季度。

这意味着，数据中心硬件的增长故事已经不再是一个“所有参与者都受益”的Beta故事，而是一个需要判断“谁在供应链中不可替代”的Alpha故事。随着AI基础设施投资从爆发期进入持续建设期，订单能见度将成为区分股价表现的核心变量。

3. 智能手机供应链的beat信号值得深挖：不是周期复苏，是结构替换

五月营收中，智能手机供应链同样出现了超出预期的数据。鸿海（Hon Hai）营收859,409百万新台币，beat预期9\%，同比增长40\%；臻鼎（Zhen Ding）beat 2\%，同比增长37\%；国巨（Yageo）beat 3\%，同比增长47\%。这些数据放在消费电

[... middle omitted ...]

个独立于手机大盘的结构性增长曲线。

4. 报告尚未完全回答的关键问题：AI需求能否持续支撑当前估值溢价

尽管五月营收数据整体强劲，但一个无法回避的问题是：当前台湾科技股的估值溢价是否已经充分反映了这些增长？台湾加权指数自5月4日以来已上涨11\%，而同期营收增速为4\%环比、39\%同比。股价涨幅远超营收环比增速，这意味着市场已经在为未来几个季度的增长提前定价。

报告本身并未直接回答估值是否合理的问题，但它提供了一个重要的观察框架：摩根斯坦利推荐的股票并非营收增速最高的公司，而是那些“具备健康订单可见性或领先地位”的公司。这实际上是在暗示，市场正在从“买增长”切换到“买确定性”。那些订单能见度低、依赖一次性订单或概念炒作的股票，可能面临估值回调的风险。

这个问题的答案，取决于AI基础设施投资的可持续性。如果AI需求在2026年下半年出现任何放缓迹象，当前估值溢价中最脆弱的部分将是那些订单能见度不足半年的公司。反之，如果AI需求持续超预期，则当前估值可能仍有上行空间。

\subsection{[R038] 摩根斯坦利：大宗商品周期的分化正在重塑中国消费股的盈利格局}
{\small\color{refgray}归类：地产与消费：分化中的结构性机会与成本端出清}

大宗商品价格从来不是消费股投资的核心叙事——直到它开始系统性地改变利润表的每一个科目。

这份摩根斯坦利在5月底发布的原材料价格图鉴，表面上是一张月度商品价格追踪表，但真正值得决策者关注的，不是锡涨了还是铜跌了，而是一个更结构性的信号：中国消费行业正在经历一轮“成本端的分化式出清”。上游价格不再同步涨跌，而是沿着各自供需逻辑走向截然不同的方向。这种分化，正在把一批公司推入利润修复通道，同时让另一批公司面临被成本压垮的风险。

市场目前的定价，更多聚焦在需求端的弱复苏上，对成本端正在发生的结构性变化定价不足。这正是这份报告最值得细读的地方。

原奶价格在5月环比微涨0.2\%，全年累计下跌约1\%。这个数字本身并不惊艳，但结合供给端的减产趋势，意义完全不同。

报告明确指出，上游供给削减后，原奶价格已经趋于稳定，并预计2026年将逐步回升。这意味着持续了两年的原奶成本下行周期正在接近尾声。对于伊利和蒙牛这两家龙头来说，这既是挑战也是机会。

挑战在于，成本上升将压缩低端产品的利润空间。但机会在于，原奶价格上涨会加速中小乳企的退出。伊利和蒙牛均将2026年液态奶恢复增长作为目标，并预期在原奶价格回升过程中，能从中小玩家手中夺取市场份额。

更重要的是，报告提到“2026年存货拨备将大幅减少”。过去两年，原奶价格持续下跌导致乳企面临大量存货减值压力，这部分非经营性损失严重拖累了净利润。一旦存货拨备压力消退，净利率的修复弹性可能超出市场预期。

这里的关键判断是：原奶价格企稳不是一个简单的成本拐点，而是一个行业竞争格局的转折点——成本下行期是中小企业的“续命期”，成本上行期才是龙头收割份额的真正窗口。

生猪价格5月环比微涨0.8\%，但整体仍处于历史低位。市场对生猪板块的关注点通常集中在猪价何时反弹，但这份报告提供了一个更精细的视角：猪价持续低迷的真正意义，是加速行业产能出清。

报告对牧原股份的判断值得反复品味：“尽管短期盈利承压，但猪价持续下跌应会加速行业产能合理化，并推动价格拐点提前到来，我们预计拐点在2026年中左右。”

这意味着，当前的低猪价不是终点，而是触发下一轮上行周期的必要条件。牧原作为行业成本领先者，在行业亏损期仍能维持相对优势，一旦猪价回升，其盈利弹性将远超同行。

报告给出的牧原牛市情景目标价较当前股价有116\%的上行空间，这一数字本身就隐含了分析师对产能出清后价格弹性的高度预期。

但需要警惕的是，这一判断的核心假设是“产能出清能够如期推进”。中国生猪养殖行业的产能出清历来充满变数，散户退出速度、规模企业逆势扩张意愿、政策干预等因素都可能改变节奏。这是报告没有完全展开的关键风险点。

饮料行业是这份报告中最值得拆解的子板块，因为成本端的分化在这里表现得最为极致。

PET价格5月环比下跌3\%，但全年仍累计上涨4\%，且与油价高度相关。报告对康师傅和统一企业的判断非常谨慎：PET价格维持高位，预计从二季度末到三季度开始，将对两家公司的饮料业务利润率形成压力。同时，棕榈油和小麦价格上涨也将对方便面业务构成压力，尽管“大体可控”。

相比之下，东鹏饮料的成本优势显得极为突出。报告指出，东鹏已经锁定了2026年全年的PET成本和约半年的糖成本，且锁定价格均低于去年水平。这意味着，在同行面临成本上行压力时，东鹏反而享有确定的成本优势。

这种成本锁定能力的差异，本质上反映的是公司规模、议价能力和供应链管理水平的差距。东鹏能够在年初就锁定全年PET成本，说明其与供应商的关系、订单规模和预付款安排都具备行业领先水平。这种能力在成本上行周期中，可以直接转化为利润率和市场份额的优势。

对于啤酒板块，报告给出的判断相对乐观：“整体成本压力可控”。尽管铝价上涨，但其他原材料成本保持稳定，加上渠道结构改善、消费者促销减少

[... middle omitted ...]

果消费者因金价下跌而推迟购买、或转向按克计价产品，那么销量下滑可能完全抵消利润率改善的正面影响。

报告指出，市场情绪短期内可能仍将承压，因为资本市场开始对需求产生怀疑，而长期增长潜力则依赖于品牌建设。但品牌建设的效果需要时间验证，短期内缺乏可见的催化剂。

这里有一个报告没有完全回答的问题：老铺黄金的品牌溢价是否足以在需求疲软期维持销量？如果消费者对固定价格黄金产品的需求弹性高于预期，那么即使金价下跌，销量下滑的幅度也可能超出当前模型的假设。这是投资者需要自行判断的关键变量。

读完这份报告，一个更深层的问题浮现出来：原材料价格变化对消费股的影响，是否仅仅停留在成本端？

答案是否定的。成本变化会引发一连串的“二阶效应”，而报告受限于篇幅，对这些效应的展开有限。

第一个二阶效应是竞争格局的变化。成本上升期，中小企业被迫退出或收缩，龙头企业获得份额；成本下降期，中小企业获得喘息空间，竞争加剧。当前原奶和生猪行业正处于成本上行前的“黎明期”，竞争格局的变化才刚刚开始。

\subsection{[R039] Bernstein：人形机器人真正的竞争壁垒不是技术，而是分化能力}
{\small\color{refgray}归类：未被主题摘要引用}

Bernstein：人形机器人真正的竞争壁垒不是技术，而是分化能力

人形机器人赛道正在经历一场奇特的拥挤。超过150家玩家挤满了供应链的每一个环节，而且这个名单每个月都在变长。按照常规逻辑，这指向一个低门槛市场，不值得严肃对待。但Bernstein这份研报给出了一个截然相反的判断：真正的问题不是“进入门槛”，而是“分化空间”。走进人形机器人这扇门很容易，但门后的房间天花板极高。

这份报告以四个角色——科学家、发明家、工程师、园丁——构建了一套极为精炼的竞争分化框架。它不是在罗列哪家公司做了什么事，而是在回答一个更本质的问题：在这个产业从实验室走向商业化的过程中，什么样的能力结构才能让一家企业不被淹没，反而在拥挤中建立起真正的护城河。

我们需要认真对待这个框架，因为它揭示了一个反直觉的事实：人形机器人产业最稀缺的资源，不是某个单一技术突破，而是同时驾驭多种角色的能力。而这恰恰是当前市场定价中尚未被充分认知的部分。

1. 产业拥挤不是信号，真正值得关注的是谁在拥挤中拥有“多角色”能力

超过150家玩家的存在，让很多人得出“门槛太低”的判断。但Bernstein的分析框架告诉我们，这恰恰是理解产业竞争格局的起点，而非终点。

科学家、发明家、工程师、园丁——这四个角色分别对应四种截然不同的能力。科学家负责探索未知，解决“大脑模型”这类前沿科学问题；发明家负责将功能需求转化为形态设计和技术规格，赋予机器人“小脑”的运动能力；工程师负责把部件做得更好更便宜，在技术原理已被业界掌握的前提下追求极致的精度、强度、可靠性；园丁则负责构建生态，成为产业链各类参与者的引力中心。

关键洞察在于：大多数玩家只能胜任其中一个角色。而真正具备投资价值的公司，往往是那些能够身兼多角的选手。Bernstein在报告中明确表示，机器人OEM比零部件供应商在结构上拥有更多分化维度，因为OEM可以自然地整合大脑模型开发，将产品升级为软硬件一体化的高级套装。

这意味着，当我们评估一家公司时，不应该只看它“有没有人形机器人产品”，而要看它在这个四角色框架中占据了几个位置。单一角色的公司面临的是“在工程卓越性的狭窄路径上被迫竞争”，而多角色公司则能在更宽的维度上建立差异化。

2. 科学家角色决定了产业的天花板，但商业化时间线仍然高度不确定

大脑模型的演进是这份报告中最引人注目的部分。从LLM到VLA（视觉-语言-行动）模型，再到各种世界模型，范式在短短几年内快速迭代。Bernstein用一张密集的图表展示了2024至2026年间世界行动模型的爆炸式演进，涉及数十个学术研究项目。

但报告并没有给出乐观的时间表。它的表述非常克制：“隧道尽头已有光亮，但前方是什么仍不确定，隧道依然漫长。”科学家们正在努力带领我们穿越隧道，将未知变为已知，但这个过程需要多久，目前没有答案。

这里有一个重要的投资含义：市场可能会高估短期内大脑模型突破的速度，同时低估长期内它带来的结构性变化。当前许多公司的估值中已经隐含了“大脑模型将很快成熟”的假设，但Bernstein的框架提醒我们，科学家角色仍然处于“探索未知”的阶段，而不是“工程优化”的阶段。这意味着，依赖大脑模型突破作为核心叙事的企业，其估值面临较大的时间线风险。

相比之下，那些在发明家和工程师角色上已经建立清晰能力的公司，其价值兑现路径更为确定。它们不需要等待大脑模型的终极突破，就可以在现有技术条件下实现商业闭环。

3. 发明家角色是当前阶段最稀缺的能力，因为它需要同时驾驭广度和深度

报告对发明家角色的描述值得细读。发明家“在口袋里装满了各种东西”，他制造了各种尺寸和形态的人形机器人，有的有腿，有的有轮子，甚至有一个有八只手臂。他不断绘制应用场景地图，将功能需求转化为形态因素和技术要求，再分解

[... middle omitted ...]

视觉系统、数字显微镜、控制与驱动、3D机器视觉、测量系统，再到视觉引导机器人和工业AI。它的“口袋里有超过6000件物品”，而且还在以惊人的速度不断掏出新东西。

这个案例说明了一个关键点：发明家角色的竞争壁垒不在于单一产品的技术领先，而在于持续产生新产品的能力和速度。这是一种组织能力，而不是技术能力。它需要公司内部有一套系统化的机制，能够快速识别应用场景、定义产品规格、完成设计迭代并推向市场。

对于投资者而言，这意味着评估一家公司时，不能只看它当前的产品线有多强，更要看它推出新产品的历史节奏和组织机制。Keyence之所以能在30倍以上的PE估值下仍然被Bernstein给予“跑赢”评级，正是因为它在这个维度上建立了难以复制的优势。

4. 园丁角色是长期价值的终极来源，但需要先成为其他角色的高手

园丁是这份报告中最具想象力的角色。他不是在亲自做所有事情，而是搭建了支架、播下种子，然后偶尔照料植物。他的植物长势茂盛，上面挂着“集成商、技能包、数据实验室”等标签。

\subsection{[R040] JPM：AI服务器与PC的背离正在重塑台湾ODM的投资逻辑}
{\small\color{refgray}归类：未被主题摘要引用}

市场对台湾ODM板块的关注，往往停留在“AI服务器高增长”这一单一叙事上。但JPM最新发布的5月营收追踪报告揭示了一个更复杂的图景：AI服务器确实在加速，但传统服务器、PC、VGA/主板等板块已经出现了明显的需求分化和利润压力。这份报告的真正价值，不在于确认AI服务器的景气度——这已是共识——而在于量化了非AI业务的疲软程度，并指出了哪些公司的产品组合将在这种背离中受益或承压。

对于产业决策者和投资者而言，当前的关键问题不是“AI服务器还能涨多少”，而是“当PC和传统硬件的利润开始收缩时，哪些公司的估值安全垫足够厚”。

JPM对NVL72机架出货量的追踪显示，2026年全年出货量预估为65,000至70,000台，其中2季度单季出货量预计达到19,000台，较1季度的17,000台环比增长约12\%。这一增长主要由GB300的持续放量驱动，5月数据显示出货节奏符合预期。

但报告也提示了一个潜在隐患：由于产品从GB300向VR200的过渡，3季度可能出现“air pockets”（出货空窗期）。VR200系统的正式放量预计要到4季度才会开始。这意味着，AI服务器供应链在3季度的环比增速可能阶段性放缓，而这对于已经将高增长预期计入股价的公司而言，是一个需要警惕的短期风险。

更值得关注的是AWS ASIC服务器（Trainium 3项目）的进展。JPM在5月观察到组件拉货需求开始启动，包括服务器滑轨、电源、机箱、CCL/PCB等环节，并预计6月开始系统级出货。这对Wiwynn（纬颖）而言是明确的正面信号，因为ASIC服务器正是其差异化竞争的核心领域。报告将Wiwynn列为服务器ODM中的首选，排名高于Quanta、Hon Hai和Wistron。

2. 传统服务器需求被低估：供给瓶颈才是真正的约束，而非需求不足

与AI服务器的市场热度相比，传统通用服务器的表现容易被忽视。但JPM的数据显示，传统服务器的出货势头同样强劲，5月持续增长，2季度环比增速预计达到两位数，与Lotes（连接器厂商）给出的2季度服务器ODM出货预测（环比+15\%）一致。

真正关键的信息来自供应链反馈：传统服务器的需求增速实际在30-40\%区间，但受限于CPU供应瓶颈，实际出货增速只能达到20\%左右。这意味着，一旦CPU供应在2H26得到改善，传统服务器的出货增速有望进一步上行，并延续到2027年。订单积压的存在，为这一判断提供了支撑。

对于投资者而言，这意味着传统服务器相关的组件供应商——尤其是ASPEED（服务器管理芯片）和Lotes（连接器）——将成为这一结构性缺口的直接受益者。JPM明确将这两家公司列为“general server strength”的关键受益标的，并给予ASPEED优于大盘（OW）评级。

3. PC和VGA/主板正在经历需求与利润的双重挤压，2H26前景堪忧

这是本份报告中最具警示性的部分。JPM对PC ODM的2季度出货预期基本为环比持平或微增，但来自ODM的初步反馈显示，2H26的展望低于季节性水平，上下半年出货比例可能达到55\%对45\%，这意味着下半年将出现两位数的同比下滑。

驱动这一判断的核心逻辑是“价格弹性效应”（price elasticity impact）。报告没有详细展开，但结合上下文可以推断：内存价格上涨导致整机成本上升，从而抑制了终端消费者的换机意愿。这一效应在2H26将更加明显，并可能进一步压缩PC品牌的利润率。

VGA和主板的情况更为严峻。JPM预计2季度VGA/主板出货量环比和同比均出现两位数下滑，且这一疲软态势将延续到2H26。原因在于缺乏新产品刺激和终端需求低迷。这对于Micro-Star（微星）和ASUSTek（华硕）而言是明确的负面信号——微星约40-50\%的收

[... middle omitted ...]

比增长9\%。这一反转的背后，是新品发布数量减少和潜在的需求疲软——同样受到内存价格上涨带来的价格弹性效应影响。

这意味着，当前iPhone供应链的强劲表现，可能只是库存周期和促销活动带来的暂时性脉冲，而非终端需求的真实改善。对于依赖iPhone业务的EMS厂商而言，2H26的业绩指引可能面临下修风险。

5. 报告尚未完全回答的关键问题：HVDC的渗透节奏和利润分配机制

JPM在报告中提到了800V HVDC（高压直流）数据中心电源的进展——预计在4Q26/1Q27开始小批量导入，比此前预期的2H27有所提前。但报告也坦诚，HVDC并非VR200机架的必需配置，预计在VR200周期中的渗透率仅为10-20\%，且早期部署将偏向±400V而非800V。

第一，HVDC的渗透率假设是否过于保守？如果数据中心对能效的要求持续提升，HVDC的采用速度可能超出预期，这将直接利好Delta（台达电子）——JPM已将其列为AI数据中心电源/散热TAM增长和潜在产品涨价的受益者。

\subsection{[R041] 摩根斯坦利：中国机场股的核心问题不是客流，而是盈利结构尚未完成重置}
{\small\color{refgray}归类：未被主题摘要引用}

摩根斯坦利：中国机场股的核心问题不是客流，而是盈利结构尚未完成重置

市场对机场股的关注点长期停留在“客流恢复”这一个维度上。但摩根斯坦利在2026年6月发布的最新研报中，提出了一个更令人不安的判断：即便客流数据看起来尚可，中国三大机场的盈利修复路径已经因为三个结构性因素而显著偏离了市场预期。

这份研报不是简单的行业更新，它实际上是在重新定价中国机场资产的估值逻辑。摩根斯坦利同时下调了三大机场的目标价，降幅均达到40\%左右，并将广州白云机场的评级从“持有”直接下调至“减持”。背后的核心逻辑是：市场正在低估能源冲击对客流的持续压制、免税业务转型期的收入真空，以及新产能投产后成本分担机制的不确定性。

这不是一个“等待反转”的故事。这是一个盈利结构需要被重新审视的时刻。

摩根斯坦利的数据显示，中国航空客运量自2026年4月以来明显承压，直接原因是全球能源价格飙升推高了机票价格。四月全国航空客运量同比仅增长0.4\%，其中国内航线甚至出现了0.2\%的同比下滑。这个数字与一季度6.5\%的增速相比，是一个急剧的减速。

更关键的信号出现在出行方式的替代上。铁路客运量在同期加速增长，四月同比增速达到10.5\%，而一季度仅为5.5\%。摩根斯坦利明确指出，这是旅客从航空向铁路转移的明确证据。五一假期的数据进一步印证了这一趋势：航空客运量同比下降5.7\%，而铁路客运量同比增长4.6\%。

对于机场股来说，这意味着客流量增长的基本面正在被侵蚀。上海机场、白云机场和首都机场的国内客运增速从一季度的6\%至10\%，全面放缓至四月的3\%至4\%。国际航线同样未能幸免，增速从一季度的7\%至24\%回落至四月的4\%至18\%。其中上海机场表现最弱，研报认为与其日本航线占比较高有关。

这个趋势的含义很清晰：如果能源价格维持高位，航空出行相对于高铁的竞争力将持续被削弱。而机场的收入结构中，无论是航空性收入还是非航收入，最终都依赖于客流量的增长。当客流量增速从“恢复性高增长”切换为“结构性低增长”时，市场对机场股的估值中枢需要相应下调。

2025年12月签署的新免税合同，从条款上看确实对机场更有利。打破单一运营商垄断、引入竞争、调整租金结构为固定租金加销售提成，以及上海机场通过合资公司持有运营商49\%股权——这些变化在长期应该能提升机场在免税价值链中的议价权。

但摩根斯坦利揭示了一个被市场忽略的短期矛盾：合同切换本身会带来收入真空。

上海机场的新合同包含三个月的免租期，这直接导致一季度免税租金收入降至1.76亿元，同比暴跌49\%。首都机场也出现了人均免税消费额的下滑。新运营商需要时间建立产品品类、优化供应链，而旧运营商的退出又留下了短期空白。

研报判断，人均免税消费额已经基本见底，但“有意义的复苏需要时间”。这个“时间”到底有多长，报告没有给出明确答案。这里存在一个关键的不确定性：即便消费额触底，反弹的斜率仍取决于新运营商的运营能力和宏观经济环境。如果消费复苏慢于预期，机场免税收入恢复到历史水平可能需要数年，而不是几个季度。

对于投资者而言，这意味着机场股的非航收入在2026年全年都可能低于市场预期，而市场此前对免税业务的乐观情绪需要被重新校准。

白云机场的情况最为复杂，也最能说明为什么摩根斯坦利将其评级下调至“减持”。

2025年10月，白云机场三期和第五跑道正式投入运营，前期资本开支超过500亿元，由母公司承担。问题在于，白云机场与母公司之间的成本分摊框架至今尚未最终确定。目前的过渡安排是约9亿元的资产使用费，但最终的方案可能包含基础租金、资产收购和收入分成等多层结构，对利润的影响可能大于当前安排。

与此同时，收入端面临双重压力：能源成本高企抑制出行需求，加上T1航站楼从2026年5月开始关闭改造，抵消了T3新增产能带来的增量收入

[... middle omitted ...]

习惯向线上转移的背景下，中国机场的长期定价权是否正在被侵蚀？

免税合同的改善是机场重新获得议价权的一个信号，但这份议价权能否转化为可持续的利润增长，取决于两个尚未验证的假设：一是人均免税消费额能否真正回升到疫情前水平，二是机场能否通过引入新品类和新运营商来创造增量消费，而不仅仅是恢复存量。

此外，首都机场被剔除出港股通后估值持续承压，这一结构性因素在研报中提及但未深入分析。对于A股投资者来说，这提醒我们：机场股的流动性溢价正在发生变化，估值体系可能需要重新锚定。

基于这份研报的逻辑，评估中国机场股时，传统的“客流恢复进度”指标已经不够。需要建立一个新的观察框架，聚焦三个维度：

第一，能源价格与出行替代的弹性关系。当航空票价相对于高铁票价的价差扩大到什么程度时，旅客会显著转向铁路？这个临界点决定了机场客流增速的下限。

第二，免税收入的结构性拐点。不是看总收入的同比变化，而是看新运营商的产品上架进度、人均消费额的环比趋势，以及固定租金加提成模式下的实际收入弹性。

\subsection{[R042] 摩根斯坦利：市场低估了eBay和Etsy在智能体电商时代的独特优势}
{\small\color{refgray}归类：未被主题摘要引用}

摩根斯坦利：市场低估了eBay和Etsy在智能体电商时代的独特优势

智能体电商（Agentic Commerce）正在从概念走向落地。当行业焦点集中在亚马逊和沃尔玛如何应对这一变革时，摩根斯坦利的最新研报给出了一个反直觉的判断：在中小市值电商领域，被市场长期视为“结构性输家”的eBay和Etsy，实际上拥有被低估的结构性优势。

这份报告的核心洞察并非简单的“买入建议”，而是一个关于资产重新定价的框架。市场对这两家公司的定价，仍然停留在关键词搜索时代的竞争逻辑中，而智能体电商的范式转换，将重新定义哪些资产真正稀缺、哪些商业模式真正具备防御性。

摩根斯坦利提出的“5I框架”——Inventory（库存）、Infrastructure（基础设施）、Innovation（创新）、Incrementality（增量空间）、Income Statement Impact（损益影响）——本质上是在回答一个根本问题：当消费者从“搜索关键词”转向“对话式购物助手”时，什么才是真正的护城河？

传统电商的竞争逻辑建立在流量获取和转化效率上。谁能在搜索结果页获得更高排名，谁就能获得更多订单。这种逻辑下，亚马逊凭借规模效应和履约网络占据了绝对优势。但智能体电商改变了游戏规则：购物助手不再只是匹配关键词，而是理解意图、比较价格、发现长尾商品。

这意味着，库存的独特性和定价优势，将比配送速度更成为竞争的关键变量。而恰恰在这两点上，eBay和Etsy拥有亚马逊无法复制的结构性优势。

2. eBay的库存护城河：90\%的独特库存，是智能体电商时代的稀缺资产

市场对eBay的长期偏见，源于将其视为“二手货交易平台”的刻板印象。但摩根斯坦利的数据揭示了一个被忽视的事实：eBay拥有约25亿条在线listing，其中约90\%是二手商品或非当季商品。这意味着，这些商品在绝大多数其他平台上根本找不到。

智能体电商的核心能力之一，是能够更好地理解和匹配消费者的长尾需求。当消费者说“我想找一款1990年代生产的瑞士军刀”或“我需要一个特定尺寸的复古相框”时，传统关键词搜索几乎无法有效响应，但智能体助手可以通过自然语言理解和多轮对话，精准地从eBay的库存中匹配出结果。

这不仅仅是体验提升，而是供需匹配效率的质变。报告指出，eBay的库存天然具有定价优势——二手和非当季商品通常比新品便宜，这使得eBay在智能体助手的比价推荐中更容易获得优先展示。此外，智能体还能减少卖家的上架摩擦：eBay的AI卖货工具“Magical Listings”已显示，上架时间减少25\%以上，每位卖家的GMV提升超过10\%。

这些变化正在转化为财务数据。eBay在2025年上半年实现了约12\%的固定汇率GMV增长，远超市场预期。但市场将此视为短期现象，摩根斯坦利则认为这是结构性改善的开始。如果智能体电商成为主流，eBay有望实现中高个位数到低双位数的可持续GMV增长，并在其牛市情景下，2030年EBITDA较基准情景高出30\%。

3. Etsy的悖论：长期风险真实存在，但中期的机会窗口被市场完全忽略

Etsy的情况更为复杂，也更值得深思。市场对Etsy的担忧并非空穴来风：超过75\%的头部卖家同时在多个渠道销售，这意味着Etsy的库存并非独有；约25\%的抽成率是全行业最高，这给了卖家强烈的动机去引导交易到更低成本的渠道。

最令人不安的类比是线下旅行社：它们曾经能够对高度同质化的产品收取溢价，但一旦出现更好的用户体验和更有竞争力的定价（在线旅游），市场份额迅速蒸发。如果消费者直接向Gemini这样的通用智能体说出“给我妈妈找一份生日礼物”，而智能体直接从Etsy卖家那里获取商品信息并完成交易，那么Etsy作为“需求聚合器”的价值就会被架空。

但摩根斯坦利指出，

[... middle omitted ...]

机会率先打造针对“特殊商品”的专属智能体体验，如果能够将每位买家的GMS提高5美元，整体GMS将增长4\%——而市场目前对Etsy的4年GMS复合增长率预期仅为3\%。

这意味着，5美元的增量相当于超过一年的预期增长。在Etsy的牛市情景下，2030年GMS和EBITDA可以分别较基准情景高出29\%和40\%。

4. 报告没有完全回答的关键问题：智能体电商的采用速度和渠道迁移路径

这份报告的价值在于提供了一个清晰的评估框架，但它也留下了几个悬而未决的问题。

第一，智能体电商的采用速度。报告承认，这一判断基于智能体购物成为主流的假设，但这一过程可能比预期更慢，也可能更快。如果消费者对智能体助手的接受度低于预期，那么eBay和Etsy的牛市情景就难以实现。

第二，平台内智能体与通用智能体的竞争格局。报告认为初期平台内智能体更具优势，但这一假设需要验证。如果谷歌或OpenAI推出的通用购物助手能够更快地解决支付和物流问题，那么Etsy面临的“旅行社式”风险就会提前到来。

\subsection{[R043] 摩根斯坦利：印度钢铁的强势周期远未结束，但真正值得关注的是供给侧的再定价}
{\small\color{refgray}归类：工业与材料：供给约束与新兴需求驱动的再定价}

摩根斯坦利：印度钢铁的强势周期远未结束，但真正值得关注的是供给侧的再定价

当一家投行在Sensex指数年内下跌约13\%的背景下，其覆盖的钢铁股却逆势上涨约12\%，这本身就值得每一个产业决策者和资产配置者停下来问一句：这轮上涨的驱动力到底是什么？是短期季节性需求的脉冲，还是某种更深层的结构性变化？

摩根斯坦利最新发布的印度材料行业研报给出了一个明确的判断：印度钢铁的强势周期尚未结束。但真正值得深入解读的，不是“继续看好”这个结论本身，而是支撑这个结论的三个结构性支柱——本地化政策护城河、中国供给侧的外溢效应、以及全球地缘冲突下的成本传导机制。这三者叠加，正在重塑印度钢铁行业的定价逻辑，而市场对此的定价可能仍然不够充分。

这份报告发布于2026年6月，正值印度季风季节来临前的传统需求淡季。国内热轧卷（HRC）价格从4月高点回落约11\%，市场情绪出现分化。正是在这个看似需要谨慎的时点，摩根斯坦利选择上调多家钢铁公司的盈利预测和目标价。这种逆势判断背后，是对行业中期格局的重新审视。

以下是我们从这份报告中提炼出的五个关键洞察，以及一个报告尚未完全回答的关键问题。

市场习惯用需求侧的逻辑来理解钢铁股——基建投资、制造业扩张、房地产复苏。但摩根斯坦利这份报告传递的核心信号是：这轮钢铁股表现优异的根本原因，是供给侧正在经历一次罕见的、多重力量叠加的锁定。

首先是印度本地的贸易保护措施。报告明确指出，保障性关税的延期是支撑国内HRC价格从去年12月中旬低点反弹约26\%的关键因素。这不是一个临时性的政策工具，而是印度政府在“印度制造”战略下持续推进的产业本地化政策的延续。

其次是中国的中期“反内卷”主题。中国的钢铁出口在经历了数年的高强度竞争后，正在出现结构性放缓。摩根斯坦利认为，这并非短期波动，而是中国钢铁行业自身供给侧改革的延续。中国出口的软化，意味着全球钢铁市场的竞争格局正在发生根本性转变，印度钢企的国际定价权正在增强。

第三是成本端的支撑。尽管中东冲突推高了炸药和物流成本，但报告指出，印度钢企的原材料采购对中东地区的直接依赖有限。更重要的是，铁矿石和焦煤价格的高位运行，反而为具备资源优势的印度钢企构建了一个成本护城河——那些没有自有矿山的竞争者将更难生存。

这三个供给侧力量叠加，意味着印度钢铁行业正在经历一个“供给收缩-政策护城河-成本优势”的三重锁定。这才是这轮周期与以往最大的不同。

报告发布时，印度钢铁行业正处于季风季节前的传统需求淡季。国内HRC价格相对于进口平价已出现约8\%的折价，价差从4月高点回落约11\%。对于只看短期数据的人来说，这似乎是见顶的信号。

但摩根斯坦利的判断恰恰相反：季节性回调恰恰是结构性逻辑的验证点。

报告的逻辑链条非常清晰：季风季节的需求走弱是历史常态，但一旦需求在季风后恢复，国内价格将重新获得支撑。更重要的是，保障性关税的存在意味着进口不会大量涌入来填补缺口，国内钢企的定价权不会因为短期需求波动而被削弱。

这意味着，当前价差的收窄是一个“可预期的、可逆的”季节性现象，而非结构性恶化。对于中长期投资者而言，这种季节性回调恰恰提供了更好的入场时机。

从更宏观的视角看，印度钢铁行业的价差正在进入一个“下有底、上有弹性”的新阶段。底部由保障性关税和中国的出口限制支撑，弹性则来自印度国内需求的恢复性增长。这种定价环境，与过去那种“需求好则价差扩、需求差则价差崩”的周期模式，已经有了本质区别。

3. 中东冲突对印度钢企的影响被高估了，但一个二阶风险值得警惕

市场对中东冲突的担忧主要集中在成本端——油价上涨推高能源和物流成本，可能侵蚀钢企利润。但摩根斯坦利的分析显示，这种担忧可能被夸大了。

报告指出，印度钢企的大部分原材料采购并不依赖于中东地区。铁矿石主要来自国内和澳大利

[... middle omitted ...]

本冲击更大。

摩根斯坦利目前认为这个风险可控，但报告并未给出具体的量化分析。这正是值得读者持续跟踪的关键变量——中东局势的演变，可能成为打破当前供给侧锁定逻辑的最大外部变量。

摩根斯坦利在这份报告中维持了对JSW Steel、Jindal Steel和Tata Steel的增持评级，同时维持了对SAIL的低配评级。这种分化背后，反映的是印度钢铁行业内部正在加剧的盈利分化。

报告显示，摩根斯坦利上调了F27财年的钢铁价差预期约2\%，但不同公司的盈利调整幅度差异显著。Jindal Steel的F27 EBITDA预测上调了5\%，而JSW Steel和Tata Steel分别下调了11\%和4\%。SAIL的F27 EBITDA预测却上调了10\%。

这种看似矛盾的数据背后，隐藏着不同的成本结构和经营杠杆。拥有自有矿山、成本控制能力强的公司，在价差扩张周期中能够获得更大的盈利弹性。而那些依赖外购原料、成本结构僵化的公司，即使行业价差改善，其盈利改善的幅度也会被成本端侵蚀。

\subsection{[R044] NOM：光器件供应链的瓶颈不在产能，而在核心衬底的锁定能力}
{\small\color{refgray}归类：AI/算力：从模型竞赛到成本信任，供应链紧俏重塑定价权}

当市场还在讨论AI算力需求对光模块出货量的拉动时，一个更根本的约束正在浮出水面：光器件的核心原材料——磷化铟（InP）衬底——已经进入实质性紧缺阶段。这不是一个遥远的预警，而是已经传导至行业头部公司采购策略的当下事实。

NOM在最新发布的研报中指出，全球主要光器件生产商Lumentum Holdings在6月9日的投资者交流中，明确表示需要寻找新的InP衬底供应商。这距离该公司5月5日财报会议上“长期合同保障供应稳定”的表态，仅隔了一个月。一个月内，一家行业龙头从“供应可控”转向“必须找新来源”，意味着供需紧张程度已经出现了质的跃迁。

这份报告的核心判断是：InP衬底的供应格局正在从“相对均衡”转向“结构性偏紧”，而这一变化的真正赢家，不是那些组装能力最强的光模块厂商，而是掌握衬底产能的垂直整合型企业。具体而言，NOM认为住友电工（Sumitomo Electric Industries）拥有全球最大的InP衬底产能，并且同时具备光器件芯片的自制能力，这种“材料+器件”的双重优势正在成为竞争壁垒，而市场可能尚未充分定价这一结构性变化。

以下，我们将从供应链动态、竞争格局、以及尚未被充分回答的问题三个层面，拆解这份研报的核心逻辑。

1. 从“长期合同”到“寻找新供应商”的转变，是供需格局的拐点信号

一份长期合同的意义在于锁定价格和供应量。当一家企业从依赖长期合同转向主动寻找替代供应商，通常只有两种可能：一是现有供应商无法满足增量需求，二是现有供应商的报价或条款已经不可接受。无论哪种情况，都指向同一个结论——InP衬底市场已经进入卖方市场。

Lumentum的转变发生在短短一个月内。5月5日财报会议上，该公司还表示长期合同能够确保InP衬底的稳定采购。但到了6月9日，随着光器件需求的进一步上升，其表态已经变为“需要寻找替代供应商”。NOM在研报中评价认为，这表明光器件需求最近正在进一步走强。

这一判断值得认真对待。InP衬底是制造连续波激光二极管（CW-LD）和电吸收调制集成激光器（EML）芯片的核心基材，而这些器件正是当前数据中心高速光模块（400G/800G）以及未来1.6T模块中不可或缺的组件。AI集群对算力互联带宽的需求，正在以超出预期的方式拉动光器件的出货量。当终端需求以季度环比20\%以上的速度增长时，上游衬底的产能扩张周期（通常需要12-18个月）根本无法同步跟上。

所以，Lumentum的“一个月变脸”不是个案，而是整个产业链供需再定价的缩影。对于投资者而言，关键信号不在于Lumentum做了什么，而在于它为什么不得不做——这意味着InP衬底的供应弹性已经耗尽。

2. 住友电工的“材料+器件”双重身份，正在转化为不对称竞争优势

NOM在这份研报中明确给出了一个投资结论：看好住友电工（评级为Buy），其核心逻辑在于公司拥有全球最大的InP衬底产能，并且同时生产CW-LD和EML芯片，部分器件还对外销售。这种垂直整合模式，在供应链紧张的环境下，正在转化为实实在在的竞争优势。

第一层是材料供应保障。在InP衬底紧缺的背景下，住友电工可以优先保障自身器件产线的原材料供应。这与那些完全依赖外部衬底采购的竞争对手形成了鲜明对比。后者不仅要面对价格上涨，还可能面临供应中断的风险。NOM在研报中特别指出，住友电工在核心材料供应上的优势，使其能够比竞争对手更灵活地扩大光器件产能。

第二层是利润率优势。当衬底价格上升时，住友电工作为衬底生产商可以直接受益；同时，其器件业务的成本端受衬底涨价的影响也远小于竞争对手。这种“双重受益”的结构，意味着在InP衬底紧缺周期中，住友电工的盈利能力可能比市场预期的更具弹性。

NOM认为，这种优势“将进一步扩大”。这里的逻辑是：随着紧缺程度加深，衬底的议价权将从器件厂商向上游材料厂商转移。住友电工恰好站在这个转移方向的正中央。

3. Lumentum的供应链多元化策略，揭示了一个尚未被充分定价的行业趋势

Lumentum寻找替代供应商的动作，并非仅限于日本厂商。NOM在研报中披露，Lumentum还计划增加从英国供应商的采购。2025年12月，该公司已经与英国企业IQE签署了多年的InP外延片采购协议，并且可能会进一步增加从该供应商的采购量。

这一信息的意义在于：它揭示了光器件行业正在经历一次供应链地理重构。长期以来，日本企业在InP衬底领域占据主导地位，住友电工、三菱化学等公司合计占据了全球大部分产能。但Lumentum的举动表明，终端客户正在主动推动供应商多元化，以降低对单一区域的依赖。

这并不意味着日本企业的地位会被削弱。恰恰相反，在供给紧张的环境下，拥有成熟产能和长期客户关系的日本厂商反而更具议价能力。Lumentum的多元化采购更多是一种“对冲”而非“替代”。但这一趋势本身值得关注：它可能会加速非日本厂商在InP衬底领域的产能建设，从而在中长期改变行业竞争格局。

对于住友电工而言，短期内的关键变量不是市场份额被侵蚀，而是其能否利用当前的供应优势，与更多客户签订长期供货协议，从而锁定未来的收入可见性。NOM的报告没有直接讨论这一点，但这是从供应链重构逻辑中可以合理推导出的下一步判断。

4. 报告尚未完全回答的关键问题：衬底紧缺的可持续性与替代技术风险

\subsection{[R045] 美国银行：市场正在为一场“通胀高于增长”的资产重新定价做准备}
{\small\color{refgray}归类：宏观与利率：通胀结构分裂与政策信号再校准 / 风险偏好：资金流拥挤与‘通胀高于增长’的再定价}

美国银行：市场正在为一场“通胀高于增长”的资产重新定价做准备

这份报告最值得看的判断，不是某个资产类别下周的涨跌，而是美国银行首席投资策略师Michael Hartnett团队通过资金流数据揭示的一个正在成形的宏观叙事：通胀正在重新成为主导变量，而资产配置却仍停留在“看涨惯性”中。当通胀预期与政策应对之间的落差被市场意识到时，资产价格的重估可能比大多数人预期的更剧烈。

报告的核心信号有两个层面。第一个是数据层面的确认：美国通胀已经超过4\%，而特朗普在通胀问题上的支持率跌至27\%，低于拜登任内低点。第二个是行为层面的警示：尽管美国银行的牛熊指标已经连续四周发出“卖出信号”，但资金仍在涌入股票，尤其是科技股——上周科技基金录得创纪录的123亿美元流入。这种“知道风险但仍在加仓”的张力，正是资产定价中最值得警惕的状态。

为什么现在重要？因为历史提供了清晰的参照系。报告特别点出了1994年的类比：当时美联储在滞后于曲线后被迫大幅加息，最终导致债券收益率飙升、股市陷入多月的交易区间，直到墨西哥比索危机和橙县破产事件才完成了去杠杆。今天的情景——通胀走高、就业市场紧张、政策制定者仍在强调“美国优先”——与1994年有着令人不安的相似性。

以下是我们从这份报告中提炼的五个核心洞察，它们共同指向一个判断：当前资产配置的“看涨惯性”正在积累风险，而真正的机会可能出现在那些已经被资金抛弃、但基本面正在改善的角落。

1. 资金流显示市场正在形成“科技单边市”，但历史表明这种拥挤往往以剧烈调整收场

上周的数据是惊人的。全球股票基金净流入315亿美元，其中美国股票占174亿美元，已是连续第11周净流入——这是自2025年12月以来最长的连续流入记录。而在这174亿美元中，科技板块独占了123亿美元，创下历史单周最高纪录。其中，Direxion每日半导体三倍做多ETF和iShares半导体ETF分别吸引了30亿和29亿美元。

但与此同时，报告中的牛熊指标已经从8.7升至8.8，离“极度看涨”的10仅一步之遥。自2002年以来，该指标共有17次触发“卖出信号”，平均而言，全球股票在随后2-3个月内下跌2-3\%，命中率约60\%，最大回撤可达15-20\%。

这里有一个值得深思的矛盾：资金正在以创纪录的速度涌入科技，尤其是半导体，但牛熊指标却表明市场整体已经过于拥挤。这意味着什么？意味着科技板块的上涨可能已经不再是基本面驱动的，而是“怕踏空”心态下的被动配置。当所有人都坐在船的一侧时，任何外部冲击都可能导致船体剧烈晃动。

报告没有完全展开的一个问题是：如果科技资金流出现逆转，哪些板块最可能成为承接资金的目的地？从美国银行私人客户过去四周的ETF资金流来看，材料、MLP和TIPS是买入最多的，而科技本身反而是净卖出的。这说明“聪明钱”已经在悄悄调整仓位，但散户和被动资金仍在涌入科技。

2. 通胀正在从“尾部风险”变成“核心矛盾”，但资产价格尚未充分反映这一转变

报告中最引人注目的数据之一，是美国通胀率已经超过4\%，而特朗普在通胀问题上的支持率仅为27\%，低于拜登任期内28\%的低点。与此同时，美国原油库存降至48天，接近45年来的最低水平。这组数据放在一起，讲述了一个清晰的叙事：供给侧的约束正在推高物价，而政策制定者对此的应对空间有限。

更值得关注的是报告中的历史对照。当CPI超过4\%时，标普500指数在随后3个月平均下跌4\%，随后6个月平均下跌7\%。更重要的是，当失业率低于通胀率时——目前美国失业率4.3\%，通胀4.2\%——历史表明这往往伴随着美联储加息周期，如1966年、1973年、1990年、2000年、2008年和2021年。这些年份在华尔街的集体记忆中都不算愉快。

报告还提供了一个前瞻性的路径图：如果美国CPI月度

[... middle omitted ...]

币基金在过去五周累计流出66亿美元，创下历史纪录。

这组数据意味着什么？至少有两个可能的解释。第一个是“风险偏好转移”：资金从避险资产（黄金）和另类资产（加密货币）流向风险资产（科技股）。第二个是“流动性收缩”：这些资产的流出可能反映了投资者在科技股上涨后需要重新平衡组合，或者是在通胀上升背景下被迫减持非生息资产。

报告提到，黄金的去杠杆过程与1994年的情景有相似之处。当时，黄金价格在美联储加息周期中经历了大幅调整。如果历史重演，黄金可能还有进一步下跌的空间。但这里有一个报告没有完全展开的关键问题：黄金和加密货币的资金流出是暂时的战术调整，还是结构性趋势转变的开始？如果是后者，那么这些资产的估值中枢可能需要重新审视。

4. 新兴市场正在出现“抄底信号”，但需要区分哪些国家真正具备韧性

一个值得注意的积极信号是，新兴市场股票基金在连续九周净流出后，上周首次转为净流入，金额为45亿美元。其中，韩国股票基金录得59亿美元净流入，是2026年3月以来的最大单周流入。

\subsection{[R046] 摩根斯坦利：阿尔茨海默病的真正拐点不在新药，而在“脑穿梭”能否重构成本曲线}
{\small\color{refgray}归类：生物医药与健康：中国创新资产需求刚性，阿尔茨海默病成本曲线待重构}

摩根斯坦利：阿尔茨海默病的真正拐点不在新药，而在“脑穿梭”能否重构成本曲线

这份报告最值得看的判断，不是阿尔茨海默病市场规模有多大，而是：美国正在经历一场由人口结构决定的健康支出被动升级，而阿尔茨海默病是这场升级中最集中的成本放大器。 摩根斯坦利在这份题为《Can Innovation Ease Alzheimer's Growing Burden?》的研报中，系统性地论证了一个让投资者重新审视生物科技板块底层逻辑的观点——真正值得关注的不是任何一个单药能否成功，而是“脑穿梭”技术平台能否将阿尔茨海默病从不可控的财政黑洞，转化为可管理的治疗路径。

这不是一份常规的行业综述。它从人口统计的硬约束出发，推演到医疗支出的机械性增长，再落到生物技术能否改变这一轨迹。它的核心洞察是：我们正在把增加的寿命，系统性地转化为增加的失能年数。阿尔茨海默病是这个转化过程中最高效的“催化剂”。

摩根斯坦利的分析起点，是一个容易被忽视的宏观事实：美国人口总量在2050年达到约3.64亿峰值后，将进入长期缓慢下降，到2099年基本回到当前水平。这意味着，未来75年美国医疗支出的增长，几乎完全由年龄结构的“顶部加重”驱动，而非人口总量的扩张。

数字本身已经足够说明问题。目前美国老年人（65岁以上）占人口约18.8\%，到2099年将升至30.5\%。而老年人的年均个人医疗支出约为22,400美元，是成年人的2.4倍、儿童的5倍以上。仅凭年龄结构变化，美国个人医疗支出占GDP的比重将从当前的11.7\%被动推升至2050年的约13.1\%。

报告估算，当前美国阿尔茨海默病相关市场规模约为500亿美元，基于约700万患者、80\%为轻中度、其中30\%经筛查后适合治疗的假设。而到2050年，仅阿尔茨海默病患者的绝对数量就将从约610万翻倍至约1270万。这不是需求在增长，而是“需求池”在不可逆地扩大。

这里的关键洞察是：我们通常将医疗支出增长归因于新药、新技术、或过度医疗。但摩根斯坦利提醒，最被低估的因素是人口结构的机械性推力。这个推力不受政策、经济周期或保险设计的影响，它几乎是确定性的。

当前最热门的医药叙事无疑是GLP-1受体激动剂。从减肥到心血管保护，从NASH到慢性肾病，这些药物几乎被描绘成解决一切代谢相关问题的万能钥匙。但摩根斯坦利在这份报告中划出了一条清晰的边界：GLP-1对阿尔茨海默病无效。

这个判断基于一个硬数据：诺和诺德的口服司美格鲁肽在EVOKE/EVOKE+三期临床试验中，未能减缓早期症状性阿尔茨海默病的进展。试验于2025年11月终止。报告明确指出，“肥胖/糖尿病管线主导了当前的长寿讨论，但它们不太可能缓解尾端神经系统的负担。”

这一结论的So what含义是双重的。对于投资者而言，它意味着GLP-1领域的成功并不能线性外推到大脑健康领域。对于产业决策者而言，它意味着阿尔茨海默病的治疗突破不能依靠“搭便车”式的代谢风险改善，而需要专门针对神经退行性机制的创新。

报告用一个历史类比强化了这一观点：制药创新确实多次弯曲了死亡率曲线——疫苗、抗生素、他汀、抗逆转录病毒、靶向肿瘤药、GLP-1——但阿尔茨海默病始终是“尾端剩余负担”中最大的那个谜题。它不是另一个可被同一套工具解决的疾病。

当前已获批的疾病修饰疗法——Biogen/Eisai的Leqembi和礼来的Kisunla——验证了淀粉样蛋白清除可以减缓认知衰退，约27\%的CDR-SB评分改善，约0.45分的18个月差异。但摩根斯坦利认为，淀粉样蛋白清除已经接近其疗效天花板。

报告的逻辑是：淀粉样蛋白清除率已经接近最大化，更快的清除或更深的清除不太可能带来更好的临床结局。即使通过“脑穿梭”技术递送更好的抗淀粉样蛋白抗体，它仍然作用于同样的Aβ级联反应，因此继承

[... middle omitted ...]

 mRNA的脑穿梭寡核苷酸）和Arrowhead的ARO-MAPT（一种BBB穿透性siRNA）是仅有的两个脑穿梭增强的tau项目。如果tau数据能够证实其作为疾病修饰靶点的价值，这将是比淀粉样蛋白清除更根本的突破。

4. 脑穿梭：从“递送失败”到“递送是可控的障碍”的范式转换

报告认为，Denali的Avlayah（tividenofusp alfa）在2026年3月获得美国加速批准用于Hunter综合征（MPS II）的神经系统表现，是脑穿梭技术的一个标志性事件。这是首个获得FDA批准的脑穿梭项目，它实质性地将“大脑生物利用度”从一个不可逾越的障碍，转化为一个可工程化解决的问题。

在2026年之前，一个项目可能因为分子从未达到治疗浓度而失败——对于被动脑穿透率低于0.1\%的大分子生物药，这几乎是常态。Denali的批准改变了这一前提。报告指出，脑穿梭平台使大分子生物药、酶替代疗法、寡核苷酸变得可递送，但“它并不使它们在其认知分子靶点上天然地更有效或更无效。”

\section{Disclaimer}
{\scriptsize\color{refgray}
This community is a paid learning and information-sharing community created for general educational purposes only. The membership fee is charged solely for access to community discussions, educational content, organization, and operational costs. It is not an advisory fee, management fee, performance fee, brokerage fee, or compensation for personalized financial, investment, legal, tax, or professional advice.\par
I participate and share information solely as an individual and for educational discussion among community members. I am not acting as a financial adviser, investment adviser, broker, dealer, portfolio manager, legal adviser, tax adviser, accountant, fiduciary, or any other licensed professional adviser. Nothing shared in this community constitutes, and should not be interpreted as, financial advice, investment advice, legal advice, tax advice, a recommendation, solicitation, offer, endorsement, instruction, or guarantee to buy, sell, hold, short, trade, or invest in any security, cryptocurrency, fund, derivative, strategy, or other financial product.\par
All content shared in this community, including messages, discussions, links, files, charts, examples, market commentary, watchlists, AI-generated or AI-polished summaries, and third-party materials, is general, impersonal, and educational in nature. It does not take into account any individual's financial situation, investment objectives, risk tolerance, investment horizon, liquidity needs, tax status, legal restrictions, or suitability requirements. No content should be relied upon as suitable for any specific person, account, portfolio, or financial decision.\par
Information shared in this community may be incomplete, inaccurate, outdated, speculative, AI-generated, AI-edited, or based on third-party internet sources that have not been independently verified. I make no representation or warranty as to the accuracy, completeness, timeliness, reliability, or suitability of any information shared.\par
Financial markets involve substantial risk, including the possible loss of principal. Past performance, historical data, hypothetical examples, simulated results, backtests, personal opinions, or educational examples do not guarantee future results. Each member is solely responsible for conducting their own research, exercising independent judgment, and making their own decisions. Before making any financial, investment, legal, or tax decision, members should consult qualified and licensed professionals.\par
Participation in this community, payment of any membership fee, or communication with me or other members does not create any adviser-client, fiduciary, professional, contractual, agency, partnership, employment, or other advisory relationship. By joining or participating in this community, you acknowledge and agree that you are solely responsible for your own decisions, actions, transactions, profits, losses, risks, and consequences, and that I shall not be liable for any direct, indirect, incidental, consequential, financial, legal, tax, or other loss or damage arising from or related to any content, discussion, information, or material shared in this community.\par
}
\end{document}
